%% ======================================================================
%% Section 1, preceded by the front matter (\maketitle and the abstract,
%% assigned by the outline to the writer of Section 1).  Uses only the
%% macros and environments of the preamble; defines nothing.
%% ======================================================================
\maketitle

\begin{abstract}
The Erd\H{o}s--Gallai conjecture (Erd\H{o}s Problem \#184) asserts that the edge
set of every $n$-vertex simple graph can be partitioned into $O(n)$ cycles and
single edges. The best published bound is $O(n\logs n)$, due to \BMnames{}
\cite{BM} (arXiv:2211.07689). This manuscript presents a \candidate{} proof of
the bound $O(n)$: the linear bound is established here \emph{only as a
candidate}.

The candidate argument has three ingredients beyond \cite{BM}. First,
\emph{$\tau$-rules} placed inside the recursion of \BMnames's Lemma~14, together
with a cap on the degrees of guest vertices, produce a hierarchy $\HBtp$ in which
hub concentration is negligible for every run. Second, \emph{standalone parts are
chained} into cycles through junction paths that run inside lent random colour
classes of ancestors at least two rounds back (a transparent vortex, TPV,
together with a junction synchronisation that splits every large class into
cherries, JS-LC); only parity leftovers, single edges, remain. Third, \emph{each
end of a leftover single edge receives a private junction edge}, which joins it
to a junction vertex (possibly shared with other ends) in a sparse random pool
inside its ancestor's expander; the resulting simple quotient graphs
have at most a vanishing fraction of $n$ vertices in total (the fraction tends to
$0$ as the degree threshold $\Dstar$ grows), every decomposition of a quotient
lifts back, and the quotients are handled by induction on the number of
vertices.

This is a candidate proof. All checking so far was done by AI referees
(statement-level referees and two AI red teams, which found no fatal or major
error and a list of minor fixes, all applied here). It has not been verified by
human experts and must not be regarded as a proof of the Erd\H{o}s--Gallai
conjecture until it is. Two further AI audits of the chaining engine inside
JS-LC likewise found no fatal or major error. A subsequent clean-room AI review
of this manuscript found no fatal error but one major gap (a per-vertex bound on
arc ends in the P-vortex lemma, used in the chaining of light parts); it is
repaired here without new mathematics, and a second clean-room AI review
found the repair correct and no error of the same kind elsewhere. Three further
clean-room AI reviews and a cross-model clean-room AI review (AI referees running
on a different model) each found no fatal or major error. This is
version~6.1 of the text (26 September 2026): version~6 with the repairs of
version~6.1 described below. Version~6 simplifies version~5 for a formal
verification: an elementary long-cycle bound replaces Theorem~2 of \cite{BM},
a Bernstein inequality proved in full replaces Freedman's inequality, Haxell's
and Hall's theorems replace the Aharoni--Haxell and max-flow min-cut theorems,
the vortex lemmas use finite probability spaces, the galactic conditions are
stated as eventualities, and the dependency data are completed; the numerical
constants of the final bound ($1085$ and $1091$) do not change. The parts written or changed for version~6 have had
one clean-room AI review (G0, 2026-09-26: eight AI referee lenses and an AI
judge), which found no fatal or major error; its minor and cosmetic fixes are
applied here and have not been re-reviewed, and a proof of the stated form of
Haxell's theorem (Cited result~\ref{s1:citHaxell}) was added after G0 and has had no
referee. Version~6.1 applies the repairs found by the triage of a formalization in
Lean (an integral $M_l$, declared arguments for the fixed rules of the round step,
and wording and proof-step repairs); no constant changes, no statement used
downstream is weakened (Lemma~\ref{s6:lemJSLC}, whose statement named an object
internal to its proof, is restated in the form its only consumer uses), and these
repairs have had no referee yet. No human expert has refereed any
part of this text. The constants are galactic and
are not computed: the degree threshold $\Dstar$ is only required to be sufficiently large (each
condition on it holds for all sufficiently large $\Dstar$), and it is of tower type
($\log_2\log_2\Dstar\ge2^8$ is one of these conditions).
\end{abstract}

\section{Introduction, conventions, cited inputs, constants, dependency structure}\label{s1:sec}

For a graph $G$ let $\f(G)$ be the least number of cycles and single edges that
partition $E(G)$, and let $\f(n)$ be the maximum of $\f(G)$ over graphs $G$ on
$n$ vertices (Definition~\ref{s1:defObject}). Erd\H{o}s and Gallai \cite{EG66}
conjectured that $\f(n)=O(n)$. Repeatedly removing a longest cycle gives
$\f(n)=O(n\log n)$ \cite{EG66}; Conlon, Fox and Sudakov \cite{CFS14} proved
$\f(n)=O(n\log\log n)$; and \BMnames{} \cite{BM} proved $\f(n)=O(n\logs n)$,
which is the best published bound. Lov\'asz \cite{Lov68} proved that every
$n$-vertex graph decomposes into at most $n/2$ paths and cycles.

This manuscript presents a \candidate{} proof of $\f(n)=O(n)$, stated as
Theorem~\ref{s1:thmMain} and proved (as a candidate) in Section~\ref{s7:sec}.
Its status is described in Remark~\ref{s1:remStatus} and, at the end of the
manuscript, in the final subsection of Section~\ref{s7:sec}, ``\ref{s7:ssecNotEstablished}''. Section~\ref{s1:sec}
contains the overview, the conventions, every cited result (each published result
of \cite{BM} is quoted in the manuscript's notation, with its number in \cite{BM}
and the line range of the arXiv text against which it was checked), one
elementary decomposition bound and one standard concentration inequality
that are proved in full instead of being cited (Fact~\ref{s1:factEG0} and
Lemma~\ref{s1:lemBBD}), the constants
and the single list of galactic conditions, the symbol table, the ledger of
red-team fixes, and the dependency table of all numbered statements.

\begin{theorem}[Main Theorem; candidate]\label{s1:thmMain}
Let $\eps,\sigma,\Cp,\Aexp,\Nzero,\Dstar$ be as in
Definition~\ref{s1:defConstants}, with $\Dstar$ satisfying
\ref{s1:condG1}--\ref{s1:condG4} of Condition~\ref{s1:condGamma}. Put
\[
\Czero\defeq \frac{\Dstar}{2}+\cTot+\epsOne(\Dstar)
\qquad\text{and}\qquad
\cEG\defeq\max\Bigl\{\frac{\Czero}{1-2\thetaQ},\,\frac{\Nzero}{2}\Bigr\},
\]
where $\epsOne(\Dstar)$ and $\thetaQ=\epsTwo(\Dstar)$ are the explicit functions
of $\Dstar$ defined in Proposition~\ref{s7:propCost} and
Lemma~\ref{s7:lemUHsplit}. Then every graph $G$ on $n$ vertices satisfies
\[
\f(G)\le\cEG\, n .
\]
In particular $\f(n)=O(n)$.
\deps{\ref{s7:thmMainProof} (forward pointer; in Table~\ref{s1:tabDAG} the proof
is the node \texttt{s7:thmMainProof}).}
\end{theorem}

The candidate proof of Theorem~\ref{s1:thmMain} is given in
Section~\ref{s7:sec}, as Corollary~\ref{s7:thmMainProof}; its status is
described in Remark~\ref{s1:remStatus}. Theorem~\ref{s1:thmMain} is the only
statement of this manuscript that is stated before the statements it depends
on.

\begin{remark}[status of this manuscript]\label{s1:remStatus}
\begin{enumerate}
\item[(i)] This is a \candidate{} proof, not an established theorem. Nothing in
this manuscript should be cited as a proof of the \EG{} conjecture.

\item[(ii)] Every referee so far is an AI agent. Each link of the chain, as
formulated in the project's working notes, was checked by two independent AI
referees (or by one AI referee together with an independent AI re-derivation).
Two AI red teams then audited the whole chain. Red team~1 returned the verdict
``survives'' (no fatal or major error) with the minor fixes
\RTa{M1}--\RTa{M12}; red team~2 returned the verdict ``survives'' (no fatal or
major error) with the minor fixes \RTb{J1}--\RTb{J8} for the composition and
\RTb{I1}--\RTb{I16} for the inputs. After red team~2, two further independent
AI audits re-derived the routing engine inside JS-LC (HCC-P, PAR, MED and EQ-LPT
as used for cherry systems under joint routing) and again found no fatal or
major error; their minor fixes (keep the $\max(1,\cdot)$ in the depth $k_i$,
state the joint-routing facts explicitly, cite no scratch computation) are also
applied. Then a clean-room AI review of this manuscript as a single text
(clean-room review~1: eight AI referee lenses that read only the manuscript and
\cite{BM}, AI verification of their findings, and an AI final judge) returned
verdict~(B): no fatal error, \emph{one major gap}, and minor findings, none of
which affects correctness. The major gap, PV-ARC-END, was found independently by
two lenses: the per-vertex arc-end bound $8\vJ+8+\vb\le2^9L^8$ claimed in
Lemma~\ref{s4:lemPV}(c) misses the arcs that stripping leaves in the finish of
the P-vortex, so the pair multiplicity used in Steps~5--6 of
Lemma~\ref{s5:lemParent} was not justified. It has been repaired, without new
mathematics, as described in the entry \textsf{CR1-PV} of
Table~\ref{s1:tabFixes}: Lemma~\ref{s4:lemPV}(c) now bounds the number of arc
ends at a vertex by its degree, so a vertex lies in at most $M_l-1$ pairs, and
the U-lent classes are path connected for the multiplicity parameter
$t_Y=\lceil\lambda_r^{1.6}\rceil\ge M_l$ (Definition~\ref{s3:defCOL}; rows~4
and~7 of Table~\ref{s3:tabCOLJV} re-derived). \emph{The repair has been re-reviewed}: clean-room review~2 (two AI lenses on the
repaired statements, one of them adversarial with computational tests, two AI
lenses auditing every site where a per-vertex multiplicity feeds a
path-connectivity parameter~$t$, and an AI judge) returned verdict~(A): the
repair is correct, no such site has an error, and only presentational findings
remain. All fixes are applied in this text; Table~\ref{s1:tabFixes} records
where. Clean-room reviews~3--5 (AI referee lenses reading only
the manuscript and \cite{BM}, and an AI judge) each returned verdict~(A): no
fatal or major error, and presentational findings only.
A sixth review round, a cross-model clean-room review (six AI section referees
running on a different model from the earlier reviews) together with three
conceptual AI lenses (whether the proof proves too much, traces on concrete graph
families, and the key idea versus known barriers) and an AI judge, also returned
verdict~(A), with minor and cosmetic findings only, all applied in this text.
In blinded calibrations of this review process on copies of the text with
planted errors, the AI referees flagged every planted error that lay inside the
sections they were asked to review ($9$ of $9$ distinct plants), missed the one
plant placed outside those sections, and raised no false fatal or major alarm
in $21$ passes over the unmodified text.
Such AI review cannot exclude conceptual gaps that leave no local
inconsistency in the text, or blind spots shared by the model family that wrote
the proof and the one that reviewed it; verification by human experts is
required.
Correlated blind spots between AI referees are possible: agreement among
AI referees is not independent evidence in the sense in which agreement among
human experts is.

\item[(iii)] This manuscript, as a single text, has had six AI review rounds
(clean-room reviews~1--5 and the cross-model round of (ii)); the repair \textsf{CR1-PV} has had one
\rc{1}; and no human expert has verified any part of it. Several statements are
new formulations made for this manuscript (repackagings of refereed material),
and several were changed by \textsf{CR1-PV}; the former are marked \rc{0} and
the latter \rc{1} in Table~\ref{s1:tabDAG}.
The present version~6 of the text simplifies version~5 for a formal
verification. The parts written or changed for version~6 postdate the six
review rounds. Before the review described next, they had been checked at most
by the AI reviews of the individual changes (the integration corrections below
and the final repair rounds of the changes not at all). Then a clean-room AI
review of the version-6 changes, gate G0 of 2026-09-26 (eight AI referee
lenses that read only the manuscript and the cited papers, each covering a part
of the changes, and an AI judge), returned verdict~(A): no fatal or major
error, one minor finding (incomplete dependency rows), two bookkeeping findings
classed as minor, and cosmetic findings. It covered every statement that was
marked for version~6 and the unmarked changed text named in the legend of
Table~\ref{s1:tabDAG}. G0 is an AI review, like all earlier ones: no human
expert has refereed any part of this text. In
Table~\ref{s1:tabDAG}, the statements that were added for version~6, or whose
statement or proof was rewritten for it, are marked ``one clean-room AI review
(G0, 2026-09-26)''; changes
that concern only single proof steps, dependency entries, pointers, names,
titles or explanatory text are not necessarily marked. The fixes of G0 are
described at the end of this item. The changes are described below,
item by item, and those of item R6 in Section~\ref{s1:ssecDAG}; the project
files named in the legend of Table~\ref{s1:tabDAG} record every change
exactly. The replacement of
\cite[Theorem~2]{BM} by the elementary Fact~\ref{s1:factEG0} (item R1) wrote or
changed the following: Fact~\ref{s1:factEG0} (new); the per-step
bound~\eqref{s4:eqTPVstep}, the finish and the final count of
Lemma~\ref{s4:lemTPV}; the per-step bound~\eqref{s4:eqPVstep}, the finish and
part~(a) of Lemma~\ref{s4:lemPV}; the displayed bound, the finish and the count
of Theorem~\ref{s4:thmVXp}; the size condition~(iii) in the proofs of these
three statements, which now follows from size condition~(i); the sentence of
the introduction of Section~\ref{s4:sec} that says where simplicity of graphs
is used; Remark~\ref{s5:remConstants}(a); the table and the
last sentence of Remark~\ref{s1:remBMused}; and item~(4) of
Remark~\ref{s7:remNonCirc}.
It also removed from Definition~\ref{s1:defConstants} the requirement on
$\Nzero$ and the item that concerned the constant of \cite[Theorem~2]{BM},
which is no longer used.
The replacement of Freedman's martingale inequality by the Bernstein
inequality of Lemma~\ref{s1:lemBBD}, which is proved in full (item R2 of the
simplification of the text for version~6), wrote Lemma~\ref{s1:lemBBD} with
its proof and the concentration step in case~(b) of the proof of
Lemma~\ref{s3:lemL17s}; both are marked ``one clean-room AI review (G0,
2026-09-26)'' in Table~\ref{s1:tabDAG}.
The replacement of the Aharoni--Haxell theorem \cite[Theorem~6]{BM} by
Haxell's theorem (item R3 of the simplification of the text for version~6)
added Cited result~\ref{s1:citHaxell} with the derivation that follows it,
and wrote a new claim and a new final step for the proof of
Lemma~\ref{s3:lemL9rho}, whose statement is unchanged; no constant changes.
The statement of Cited result~\ref{s1:citHaxell}, including its factor
$2q-1$, has not yet been checked against the text of \cite{Hax95}; it agrees
with Haxell's theorem as quoted in the later literature, with $r=q+1$ the
maximum edge size. The
derivation and the new parts of the proof have had the AI reviews of this
change and one clean-room AI review (G0), and they are marked ``one clean-room
AI review (G0, 2026-09-26)'' in Table~\ref{s1:tabDAG}. Item R3 also updated
the pointers to Theorem~6 of \cite{BM} and to the new text in
Remark~\ref{s1:remBMused}, whose table now marks Theorem~6 as not used and
whose dependency line no longer lists it, and in the description of the
proof of \cite[Lemma~9]{BM} in Cited result~\ref{s1:citLem9}. Only pointers
and a dependency entry changed there, so R3 adds no version-6
marker to the rows of these two statements in Table~\ref{s1:tabDAG}.
Item R4 of the simplification of the text for version~6 gives a new proof of
Lemma~\ref{s3:lemHB}, whose statement is unchanged: the flow argument is
replaced by one application of Hall's theorem, and Cited
result~\ref{s1:citHall} replaces the max-flow min-cut theorem cited in
version~5. The new proof has had the AI reviews of this change and one
clean-room AI review (G0), and it is marked ``one clean-room AI review (G0,
2026-09-26)'' in Table~\ref{s1:tabDAG}. Item R4 also added to
Remark~\ref{s1:remNotUsed} that the Hall-type statements HALL and HALL$'$
listed there are not Hall's theorem; this explanatory sentence carries no
marker.

The replacement of the auxiliary randomness in Section~\ref{s4:sec} by finite
probability spaces (item R5 of the simplification of the text for version~6)
truncated the levels of Lemma~\ref{s4:lemTPV}, Lemma~\ref{s4:lemPV} and
Theorem~\ref{s4:thmVXp} at the number of steps. The new level has the law of
the old level capped at that number: for every $j$ up to that number the
probability that a level is at least $j$ is unchanged, and the probability of
the old values above that number is moved to that number itself. Every
probability estimate of the three proofs depends on the levels only through
their values capped at that number, and the deterministic arguments hold for
any levels with values in $\{0,1,2,\dots\}$. Item R5 also added the monotone comparison (MC) to
the observations of Section~\ref{s4:ssecConv}; re-proved with it the
probability bound for the good event (G2) of the three statements, in place of
a thinning by auxiliary coins, and the binomial domination (B), in place of a
coupling through uniform variables on $[0,1]$; recorded the level range and
the proof-local symbols in Convention~\ref{s4:convVortex}; and added a
paragraph on finite probability spaces after the observations. No probability
bound and no constant changes. These parts have had the AI reviews of this
change and one clean-room AI review (G0). They are
marked ``one clean-room AI review (G0, 2026-09-26)'' in Table~\ref{s1:tabDAG},
in the rows of
Convention~\ref{s4:convVortex} and of the three statements; the paragraph on
finite probability spaces has no row of its own and has had the same reviews.

Item R7 of the simplification of the text for version~6 completed the
dependency lines and Table~\ref{s1:tabDAG}. It added entries to the dependency
lines and rows of 40 statements (the fixes of G0 described below completed five
rows, four of them among these 40 and the fifth that of
Lemma~\ref{s7:lemLift}, so that entries were added to the rows of 41 statements
for version~6); removed from the dependency line and row of
Remark~\ref{s1:remBMused} the three results that the remark marks as not used
directly; moved the forward pointers of Definition~\ref{s1:defConstants} out of
its dependency line; marked the items that are not used; added the citations of Cited
results~\ref{s1:citHall} and~\ref{s1:citMarkov} to the proof of
Lemma~\ref{s7:lemWellDef}(iii) and to step~(f) of
Construction~\ref{s7:consRound}; corrected, in Condition~\ref{s1:condGamma},
where \ref{s1:condG2}(b),(c) are proved, and, in the introduction of
Section~\ref{s3:sec}, the list of the results that the section uses; and added
the paragraphs on the meaning of an entry and on forward references in
Section~\ref{s1:ssecDAG}. No hypothesis, conclusion, inequality or constant
changed. These parts have had the AI reviews of this change and one
clean-room AI review (G0). In Table~\ref{s1:tabDAG}, the rows of
Construction~\ref{s7:consRound} and Lemma~\ref{s7:lemWellDef} are marked ``one
clean-room AI review (G0, 2026-09-26)'' for the added citation; the other
statements that R7 changed carry no R7 marker (except
Condition~\ref{s1:condGamma}, where the statement of where
\ref{s1:condG2}(b),(c) are proved was corrected), because only their dependency
entries, their names or titles, or explanatory sentences changed.

The items R1--R7 were combined into this text on 2026-09-26 (the integration
of version~6). The integration also made the following corrections, which were
suggested by the final AI reviews of the items or found in the formalization in
Lean, and which change no constant: Lemma~\ref{s3:lemCOLJV}(ii) now asserts the
column-3 inequalities of Table~\ref{s3:tabCOLJV} at every
$\mu\ge\log\log\Dstar$, as \Gc{1}(f) gives, and the proof of
Lemma~\ref{s3:lemCOL}(b) cites it for row~2 at $\lambda=\lambda_r$; the remark
in Cited result~\ref{s1:citDef11} assumes $\varepsilon>0$; Cited
result~\ref{s1:citLem25} assumes $m\ge2$, which its only application
($m\ge2^{40}$) satisfies; Fact~\ref{s1:factEG0}(b) is stated for the edge set of
a (simple) graph; Step~5 of the proof of Lemma~\ref{s3:lemL17s} cites the form of
Chernoff's bound with a lower bound on the mean (Cited
result~\ref{s1:citChernoffGen}(a)), which it uses; two short proof passages were reworded (the case $\Pl=\emptyset$ in the proof of
Lemma~\ref{s4:lemPV} and the justification in Lemma~\ref{s2:lemTower}(b)); two
auxiliary symbols were renamed; the sentence on the non-use of
$(\Gamma2)$(b),(c) was strengthened; and some wording and bookkeeping in the
explanatory text, in Table~\ref{s1:tabOrder} and in the outline used for the
formalization was corrected. These corrections had no review before G0; they
have had one clean-room AI review (G0). They are recorded in the
project file
\texttt{proofs/\allowbreak{}manuscript/\allowbreak{}v6work/\allowbreak{}INTEGRATION\_v6.md}.
The rows of the two Cited results are marked ``one clean-room AI review (G0,
2026-09-26)'', and so are, since the fixes of G0, the change to
Lemma~\ref{s3:lemCOLJV}(ii) in the row of that lemma and the change to step~(b)
of the proof of Lemma~\ref{s3:lemCOL} in the row of that lemma. G0 did not
identify the two renamed auxiliary symbols: its two independent compilations
rule out dangling references, but not a clash in meaning.

\emph{The fixes of G0.} The judge of G0 listed three minor fixes and a list of
cosmetic fixes. All of them are applied in this text, and they are recorded in
the project file
\texttt{proofs/\allowbreak{}manuscript/\allowbreak{}g0/\allowbreak{}FIXES.md}.
They complete the dependency rows of Theorem~\ref{s6:thmMIXC},
Lemma~\ref{s7:lemLift}, Proposition~\ref{s7:propCost},
Lemma~\ref{s7:lemWellDef} and Corollary~\ref{s7:thmMainProof}; add the status
markers of Lemma~\ref{s3:lemCOLJV}(ii) and Lemma~\ref{s3:lemCOL}(b); record that
\ref{s1:condG2}(b),(c) follow from \ref{s1:condG1} alone; and correct wording,
symbol names, pointers and bookkeeping elsewhere. No hypothesis, inequality or
constant changes; the statement of Lemma~\ref{s5:lemDemoted} now records the
probability bound $1-\eta_{\mathrm{VX}}\ge1/2$ that Theorem~\ref{s4:thmVXp}
gives, in place of the weaker bound $1/2-o(1)\ge1/3$. These fixes were applied
after G0 and have not been re-reviewed. After G0, a complete proof of the
stated form of Cited result~\ref{s1:citHaxell} (Haxell's alternating-tree
argument) was also added; it has had no referee yet, and it is marked so in
Table~\ref{s1:tabDAG}.

\emph{Version~6.1 (P2 triage).} The triage of the formalization of this text in
Lean (the P2 triage of 2026-09-26, based on blueprints of every section) found no
fatal error and no false statement on which the proof relies, but several places
where a statement is true while its proof or its wording has a gap; one
meta-remark about a proof, item~(c) of Theorem~\ref{s3:thmT16s}, was false as
written and had no consumer. Version~6.1 repairs them. No constant changes (in
particular not $\cTot$, $\cTotR$, $\cTPV$, $\cPV$ and $\cVX$ of the cost line),
and no statement used downstream is weakened: the one statement whose new form
is weaker than its old one is Lemma~\ref{s6:lemJSLC}, whose v6 form referred to
an object defined only inside its proof and which is restated in the form its
only consumer uses (see (4)); the sharper v6 bound is still proved (the paragraph
\emph{The sharper bound} at the end of its proof). The changes are:
(1)~in (R2) of Definition~\ref{s2:defHBtp}, $M_l$ is now the ceiling
$\lceil\max(2^{40},2^{16}T\log^4T)\rceil$, a natural number, as its uses as a
number of classes and colours ($\KJS_l=M_l^2$, the palettes $[3M_l]$ and
$[4M_l]$) require; consequently the proof of Proposition~\ref{s2:propDegRec}
shows $M_l\le d_l^2$ with the extra $+1$, bound~(B6) in the proof of
Lemma~\ref{s3:lemCOLJV} is re-derived for the ceiling (the stated bound is
unchanged), the proof of Lemma~\ref{s2:lemCap}(ii) says $M_l\ge\max(\dots)$, and
the symbol table records the ceiling;
(2)~Construction~\ref{s7:consRound} now declares which data each of its fixed
rules may read (a new paragraph \emph{Fixed rules}: the rules of steps (b) and
(c) read only the past, those of (d), (e1) and (e2) read the past and the lists
but never the orders $\prec_u$ of the round, and the rank order of (g) and
$\Dec_l$ in (h) are unrestricted), and the claims of steps (c) and (e2), the
proofs of Lemma~\ref{s7:lemWellDef}(iii) and Lemma~\ref{s7:lemCC}(i), and the
opening paragraph of the subsection on quotient size and payments cite these
restrictions; without them those claims could fail for adversarial rules;
(3)~item~(c) of Theorem~\ref{s3:thmT16s}, a meta-remark about the proof, is
corrected: it now says that Step~0 of the proof uses the full hypothesis on $s$
(the preliminary part of the proof is named Step~0);
(4)~the statement of Lemma~\ref{s6:lemJSLC} bounds the number of cycles by
$126n/M_l+1.5\sum m_{Y,l}$, the form used by its only consumer
(Theorem~\ref{s6:thmMIXC}), and the split cost $\sco_{Y,l}$ and the number
$g_{Y,l}$ of cherry classes, which exist only inside the proof, moved into the
proof; this form is implied by, but weaker than, the v6 bound
$126n/M_l+\sum\sco_{Y,l}$, whose sharper version $15n/M_l+\sum\sco_{Y,l}$ remains
proved and is recorded at the end of the proof;
(5)~cosmetic repairs: the sketch of Cited result~\ref{s1:citLem14} now justifies
its inequality~(9) also for small $n$; the proof of Lemma~\ref{s6:lemHCCglob}
cites the argument of Fact~\ref{s1:factAdd}(b) rather than its statement; and
Section~\ref{s2:sec} states the standing assumption that $\Dstar$ satisfies
\Gc{1}--\Gc{4}.
These patches have had no referee of any kind yet (neither AI nor human). In
Table~\ref{s1:tabDAG} the changed statements are marked ``changed in v6.1 (P2
triage): no referee yet''; the project file
\texttt{proofs/\allowbreak{}manuscript/\allowbreak{}v61/\allowbreak{}PATCHES.md}
records every change with its rationale.

\item[(iv)] The most likely places for a hidden error are:
\begin{enumerate}
\item[(a)] the routing engine inside JS-LC (Lemmas \ref{s6:lemHCCP},
\ref{s6:lemPAR}, \ref{s6:lemMED} and \ref{s6:lemEQLPT} as used for cherry
systems at junction multiplicity $2M_l-2$ in Lemma~\ref{s6:lemJSLC}); red
team~2 named it the single most likely place, and it has since been re-derived by
the two AI audits mentioned in (ii), but by no human;
\item[(b)] Lemma 14$^\tau$ (Lemma~\ref{s2:lem14tau}) and part~(iv) of
Theorem~\ref{s6:thmCONCL}: together they modify the core Lemma~14 of
\cite{BM} instead of using it as a black box, and red team~1 named this the
weakest link;
\item[(c)] the composition JV$^{+*}$ of Section~\ref{s7:sec}, which is the
newest link;
\item[(d)] the statements changed by the repair \textsf{CR1-PV}
(Lemma~\ref{s4:lemPV}(c), Lemma~\ref{s5:lemChild}(c),
Definition~\ref{s5:defStages}, Lemma~\ref{s5:lemParent},
Definition~\ref{s3:defCOL}, Lemma~\ref{s3:lemCOL}(c) and
Lemma~\ref{s3:lemCOLJV}), which have had one AI re-review, and more generally
every place where a per-vertex multiplicity feeds a path-connectivity
parameter~$t$: property (P4) of the vortex runs of Section~\ref{s4:sec}, Step~6
of Lemma~\ref{s5:lemParent}, and the joint routing of JS-LC at junction
multiplicity $2M_l-2$. The gap PV-ARC-END was of this kind;
\item[(e)] the factor $2q-1$ of Haxell's theorem (Cited
result~\ref{s1:citHaxell}), as long as it has not been checked against the
text of \cite{Hax95}; its only application, in
Lemma~\ref{s3:lemL9rho}, has slack of a factor of about $h/2$, where $h$ is the path-length bound of
that proof, and the proof of
the stated form written out after the citation has had no referee yet.
\end{enumerate}

\item[(v)] Refutation targets. Any one of the following would refute the
candidate proof as written:
\begin{enumerate}
\item[(1)] a valid $\HBtp$ run (Definition~\ref{s2:defHBtp}) and a designation
with $\sum_{(Y,l)}m_{Y,l}=\omega(n)$;
\item[(2)] a valid run on which JS-LC emits two $\Jhub$-edges of one class at one
hub in one round;
\item[(3)] a lifted sub-layer cycle that repeats a vertex;
\item[(4)] an edge covered twice, in particular a junction edge used both by a
lifted cycle and by the device of its owner;
\item[(5)] a per-round cost term that is none of the three kinds listed in the
paragraph ``Per-round costs here'' of \ref{s1:ssecOverview}; in particular, a
per-round cost term of order $\Theta(n)$ that is not a once-per-vertex charge.
\end{enumerate}
\end{enumerate}
\deps{none.}
\fixes{the stale status line flagged by red team~2 (\textsf{RT2}, status line of
the JV$^{+*}$ write-up); the status stated here supersedes it.}
\end{remark}

\subsection*{Overview: the architecture, and where the $\logs$ disappears}
\namedlabel{s1:ssecOverview}{the overview}

This overview is prose; it contains no statement. Every claim it makes about the
present argument is proved (as part of the candidate proof) at the place it points
to; the paragraph on \cite{BM} describes their published proof.

\paragraph{The chain.} The candidate argument is a chain of five links:
\[
\text{UNI-MIX}\ \longrightarrow\ \text{MIX-C on }\HBtp\text{ with no VX-parts}
\ \longrightarrow\ \text{JV}^{+*}\ \longrightarrow\ \text{UH}^{*}
\ \longrightarrow\ \text{layered HI}''.
\]
UNI-MIX and MIX-C decompose all edges except a set of parity leftovers; JV$^{+*}$
joins the ends of the leftovers by private junction edges to pooled junction
vertices, and so routes them into simple quotient graphs; UH$^{*}$ controls the size of these quotients (round-split pools and
sub-layered ultra hubs); the layered induction HI$''$ closes the argument.
Section~\ref{s2:sec} builds the hierarchy $\HBtp$ (Lemma~14$^\tau$, Lemma SEP,
the overlap bounds, rule GC, ORIGIN$^\tau$, the tower facts);
Section~\ref{s3:sec} proves the random splitting and linking tools
(Lemma~15$^+$, Theorem~16$^*$, Lemma HB) and defines the lending colouring
COL-JV; Section~\ref{s4:sec} proves the vortex tools (TPV, the four-phase
P-vortex, Theorem VX$^+$); Section~\ref{s5:sec} treats light parts (zones, the
child and parent sides of the Theorem-U machinery, Lemma K-RED);
Section~\ref{s6:sec} treats standalone parts (GATE, MED, EQ-LPT, PAR, HCC-P,
Theorems CONC and CONC-L, Lemma JS-LC, Lemma J$^+$ and the MIX-C assembly for an
abstract J-consumer); Section~\ref{s7:sec} constructs the junction-vertex
quotient, bounds its size and the cost, and proves Theorem JV$^{+*}$,
Theorem HI$''$ and the Main Theorem.

\paragraph{Where \BMnames{} lose the factor $\logs n$.} The proof of Theorem~2
of \cite{BM} iterates their Lemma~26 about $\logs n$ times: each iteration
removes long cycles, splits the remaining edges into expanders by Lemma~14 of
\cite{BM} (with $s=0$), and decomposes each (necessarily small) expander into
cycles and few edges; the average degree of the leftover drops from $d$ to a
polylogarithm of $d$. Each iteration costs $\Theta(n)$ objects, charged to all
vertices that are active in it, and there are $\logs n$ iterations.

\paragraph{Per-round costs here.} The hierarchy $\HBtp$ has the same round
structure (rounds $l=1,\dots,R$, with $R=O(\logs n)$). In the candidate
argument every per-round cost term is of one of three kinds.
\begin{enumerate}
\item[(a)] \emph{Decaying in the round index.} It is bounded by a constant
multiple of one of $n/M_l$, $n/\log P_l$, $n\lambda_r^{-1/2}$, $n/P_r$,
$n M_l^3\log\lambda_{l-2}/\lambda_{l-2}^{95}$ (the parameters are those of
Definition~\ref{s2:defHBtp}) and $n\,\eta(\lambda_{l-2})$, where
$\eta(x)=\log_2\bigl(2\Aexp\log_2(\Aexp\log_2x)\bigr)/(\log_2x)^2$; the last
bounds the number $\mathrm{cap}_l$ of visit-capping pieces of the Theorem-U
chaining (Lemma~\ref{s5:lemParent}(ii), Step~4 of its proof). The hub charge is
of this kind: a vertex can be a hub in many rounds, and the round-$l$ term is
$\sum_{Z\in\Std_l}|A_Z|\le15.2\,\eps n/\log P_l$ (Proposition~\ref{s2:propOV},
(K2)), which sums over all rounds to at most $\epsA n$
(Lemma~\ref{s2:lemTower}(e)). Because the round
parameters grow like a tower towards earlier rounds, each of these sums over all
rounds to at most a quantity that depends only on $\Dstar$, tends to $0$ as
$\Dstar\to\infty$, and is uniform in $R$ (tower-lacunary sums,
Lemma~\ref{s2:lemLacunary}; the terms
$nM_l^3\log\lambda_{l-2}/\lambda_{l-2}^{95}$ sum to at most
$2F(\log_2\Dstar)\,n$, with $F(x)=(3184+30400\log_2x)\,x^{-90.2}$, by
Lemma~\ref{s7:lemUHsplit}(iv), and the terms $\mathrm{cap}_l$ are part of
$\epsChain(\Dstar)\,n$ in Lemma~\ref{s5:lemParent}(ii)).
\item[(b)] \emph{A once-per-vertex global charge.} Fresh ports (a vertex is fresh
in at most two rounds, and TPV pays $\cTPV$ per fresh port: $\cTPV\cdot2n=\cFresh
n$); parentless light incidences (at most two per vertex, each paying $\cPV$:
$\cPV\cdot 2n$); unpaired fresh legs (at most $2n$); long cycles (at most $n$); and the final leftover
$E_0$ (fewer than $\Dstar n/2$ edges).
\item[(c)] \emph{Per-ancestor terms with a factor $R-r$.} Some terms are charged
to an ancestor $Y$ of round $r$ once in every later round, so that in total they
carry a factor $R-r$: the $\tau$-, $d^*$- and guest-cap terms of the
concentration bound (Theorems~\ref{s6:thmCONC}(iii) and~\ref{s6:thmCONCL}(iv)),
and the lost-parent mass $\mathrm{lp}=\sum_{Y\in\Bad}|Y|\,(R-r(Y))$
(Definition~\ref{s5:defStages}). The factor satisfies
$R-r\le2\logs d_r+2\le\log\lambda_r$ (Lemma~\ref{s2:lemTower}(a),(d)), so it
depends on $d_r$ alone, and it is outweighed by the decay of the rest of the term
in $\lambda_r$ (for $\mathrm{lp}$, by $\Prob(Y\in\Bad)\le|Y|^{-2}$,
Lemma~\ref{s5:lemE1}(c)). Hence these terms, too, are at most (for
$\mathrm{lp}$: in expectation) $n$ times a function of $\Dstar$ that tends to
$0$ as $\Dstar\to\infty$ and is uniform in $R$
(Theorem~\ref{s6:thmCONCL}(iv), Lemma~\ref{s5:lemExpect}). This is why
$\logs\Dstar$, and not $\logs n$, appears in $\epsCONC$.
\end{enumerate}

\paragraph{The three mechanisms.}
\begin{enumerate}
\item[(a)] \emph{No hub concentration.} The $\tau$-rules inside the recursion of
Lemma~14 of \cite{BM}, together with rule GC, give
$\sum_{(Y,l)}m_{Y,l}\le\epsCONC(\Dstar)\,n$ for every valid run and every
designation (Theorem~\ref{s6:thmCONCL}(iv)). The chaining cost of JS-LC is at
most $1.5\sum_{(Y,l)}m_{Y,l}$ plus a decaying term, hence a vanishing fraction
of $n$.
\item[(b)] \emph{TPV and JS-LC.} Classed ports are never retired by the vortex
of their own part; their edges become beads, which are chained into cycles
through junction paths inside lent classes of an ancestor of round at most
$l-2$. Only parity leftovers remain; they form the J-set $J_l$ of
Lemma~\ref{s6:lemJplus}.
\item[(c)] \emph{Private junction edges.} Each end of a J-item receives a
private junction edge, which joins it to a junction vertex in a pool of density
$M_l^{-2}$ inside the expander of its ancestor, which has minimum degree larger
than $s_r/2\gg M_l$. A junction vertex may serve several ends, but every
junction edge serves at most one (Lemma~\ref{s7:lemLift}(iii)). The quotient
$Q_l$ is simple, every decomposition of $Q_l$ lifts to a decomposition of the
corresponding edges of $G$ (Lemma~\ref{s7:lemLift}), and
$\sum_l|V(Q_l)|\le\thetaQ n\le n/4$ (Lemma~\ref{s7:lemUHsplit}). The
Erd\H{o}s--Gallai bound for the smaller graphs $Q_l$ is supplied by minimality in
the layered induction HI$''$ (Theorem~\ref{s7:thmHI}).
\end{enumerate}

\paragraph{The cost line.} On the outcome chosen in Section~\ref{s7:sec}
(Proposition~\ref{s7:propCost}),
\[
\f(G)\ \le\ \Bigl(\frac{\Dstar}{2}+\underbrace{\cKRED}_{\text{K-RED}}
+\underbrace{\cFresh}_{\text{TPV, fresh}}+\underbrace{2}_{\text{fresh legs}}
+\epsOne(\Dstar)\Bigr)n\;+\;2\sum_{l}\f(Q_l),
\]
with $\cKRED+\cFresh+2=\cTot$. Together with $\sum_l|V(Q_l)|\le\thetaQ n$ and
$\thetaQ\le1/4$, the induction of Theorem~\ref{s7:thmHI} gives
$\f(G)\le\Czero n+2\cEG\thetaQ n\le\cEG n$.

\subsection{Conventions}\label{s1:ssecConv}

\begin{convention}[graphs, logarithms, neighbourhoods, balls]\label{s1:convGraphs}
\begin{enumerate}
\item[(a)] Graphs are finite and simple, unless they are explicitly called
multigraphs (auxiliary multigraphs occur only in Sections~\ref{s5:sec}
and~\ref{s7:sec}). $G$ always denotes the input graph and $n\defeq|V(G)|$. For
any graph $H$ we write $|H|\defeq|V(H)|$, as in \cite{BM}. For $k\in\mathbb N$,
$[k]\defeq\{1,\dots,k\}$.
\item[(b)] $\log=\log_2$ throughout, as in \cite[l.~141]{BM}. We put
$\log^{[0]}x\defeq x$ and $\log^{[k]}x\defeq\log\bigl(\log^{[k-1]}x\bigr)$, and
$\logs x$ is the least integer $k\ge0$ with $\log^{[k]}x\le1$
\cite[l.~144--145]{BM}.
\item[(c)] For a vertex $v$ of a graph $H$, $N_H(v)$ is its set of neighbours
and $d_H(v)=\deg_H(v)$ its degree; $\delta(H)$ and $\Delta(H)$ are the minimum
and maximum degree. For $U\subseteq V(H)$, $\Nbr_H(U)$ is the set of vertices of
$V(H)\setminus U$ that have a neighbour in $U$. $H[U]$ is the induced subgraph;
$H-U$ and $H\setminus U$ both denote $H[V(H)\setminus U]$; for
$F\subseteq E(H)$, $H-F$ is obtained by deleting the edges of $F$ (keeping all
vertices). For disjoint $A,B\subseteq V(H)$, $E_H(A,B)$ is the set of edges with
one end in $A$ and the other in $B$, and $e_H(A,B)\defeq|E_H(A,B)|$; for an
edge set $F$ and a vertex $v$, $\deg_F(v)$ is the number of edges of $F$ at $v$.
\item[(d)] A \emph{path through} $V$ is a path all of whose interior vertices lie
in $V$; its ends are unrestricted. For $U,V\subseteq V(H)$ and an integer
$i\ge0$, $\Ball{i}{H}{U,V}$ is the set of vertices of $V$ that can be reached by
a path through $V$ of length at most $i$ starting at a vertex of $U$ (the
starting vertex need not lie in $V$); thus $\Ball{i}{H}{U,V}\subseteq V$
\cite[l.~379--385]{BM}.
\item[(e)] Families of pairs of vertices are \emph{multisets}: a pair may occur
several times, and ``every vertex lies in at most $t$ pairs'' counts occurrences
(Cited result~\ref{s1:citDef7}).
\end{enumerate}
\deps{none.}
\end{convention}

\begin{definition}[objects and $\f$]\label{s1:defObject}
An \emph{object} is a cycle (of length at least $3$) or a single edge. A
\emph{decomposition} of an edge set $F$ is a partition of $F$ into the edge sets
of objects; a decomposition of a graph $H$ is a decomposition of $E(H)$.
$\f(F)$ is the least number of objects in a decomposition of $F$ (so
$\f(\emptyset)=0$); $\f(H)\defeq\f(E(H))$; and
$\f(n)\defeq\max\{\f(H):|V(H)|=n\}$.
\deps{none.}
\end{definition}

\begin{fact}[elementary properties of $\f$]\label{s1:factAdd}
\begin{enumerate}
\item[(a)] $\f(F)\le|F|$ for every edge set $F$.
\item[(b)] If $F_1,F_2$ are disjoint edge sets, then $\f(F_1\cup F_2)\le
\f(F_1)+\f(F_2)$.
\item[(c)] If $H$ is the vertex-disjoint union of graphs $H_1$ and $H_2$, then
$\f(H)=\f(H_1)+\f(H_2)$.
\item[(d)] Adding or deleting isolated vertices does not change $\f$. In
particular $\f(n)$ is non-decreasing in $n$.
\end{enumerate}
\deps{\ref{s1:defObject}.}
\end{fact}
\begin{proof}
(a) The single edges of $F$ form a decomposition of $F$ into $|F|$ objects.
(b) The union of a decomposition of $F_1$ and a decomposition of $F_2$ is a
decomposition of $F_1\cup F_2$, because $F_1\cap F_2=\emptyset$.
(c) By (b), $\f(H)\le\f(H_1)+\f(H_2)$. Conversely, the edge set of every object
induces a connected graph, so every object of a decomposition of $E(H)$ has all
its edges in $E(H_1)$ or all in $E(H_2)$; hence every decomposition of $E(H)$
splits into a decomposition of $E(H_1)$ and one of $E(H_2)$, and
$\f(H)\ge\f(H_1)+\f(H_2)$.
(d) $\f(H)$ depends only on $E(H)$, which does not change. Adding an isolated
vertex to an $n$-vertex graph gives an $(n+1)$-vertex graph with the same $\f$,
so $\f(n+1)\ge\f(n)$.
\end{proof}

\begin{fact}[the long-cycle bound]\label{s1:factEG0}
\begin{enumerate}
\item[(a)] Every graph $H$ with $h\ge1$ vertices has a decomposition into at
most $h(\log h+1)$ objects, at most $h-1$ of which are single edges. In
particular $\f(H)\le h(\log h+1)$.
\item[(b)] Let $N>1$ and $L\defeq\log N$, and let $\alpha$ and $\beta$ be real
numbers with $0\le\alpha\le N$, $\beta>0$ and $4\beta\le L$. If $F$ is the
edge set of a (simple) graph, so that $F$ has no loops and no parallel edges,
and every edge of $F$ has both ends in a set $W$ with $|W|\le\beta\alpha/L$,
then $F$ has a decomposition into at most $\beta\alpha$ objects, at most $|W|$
of which are single edges.
\end{enumerate}
The argument, iterated removal of a long cycle, is the elementary
$O(n\log n)$ bound of Erd\H{o}s and Gallai \cite{EG66} (see also
\cite[l.~88--90]{BM}); it is proved here in full. It uses that graphs are
simple (Convention~\ref{s1:convGraphs}(a)) in two places: in the claim, where a
vertex of degree at least $d$ has at least $d$ neighbours, and through
$|E(H)|\le\binom h2$.
Part~(b) is the form in which the vortex runs of Section~\ref{s4:sec} are
finished.
\deps{\ref{s1:convGraphs}, \ref{s1:defObject}.}
\end{fact}
\begin{proof}
(a) Put $m\defeq|E(H)|$, so that $m\le\binom h2\le h^2/2$.

\emph{Claim. If a graph $H'$ with vertex set $V(H)$ has $m'\ge h$ edges, then
$H'$ contains a cycle of length at least $m'/h$.} Put
$d\defeq\max\{2,\lceil m'/h\rceil\}$. Then $d-1\le m'/h$: if
$\lceil m'/h\rceil\ge2$ then $d-1=\lceil m'/h\rceil-1<m'/h$, and otherwise
$d-1=1\le m'/h$ because $m'\ge h$. Delete from $H'$, one at a time, a vertex
whose degree in the current graph is at most $d-1$, as long as such a vertex
exists. Every edge of $H'$ disappears exactly once, together with the first of
its two ends to be deleted. Suppose that all $h$ vertices are deleted. The last
of them has degree $0$ when it is deleted, and each of the others has degree at
most $d-1$ when it is deleted; hence $m'\le(h-1)(d-1)\le(h-1)m'/h<m'$, since
$m'\ge h\ge1$, a contradiction. So the deletions stop at a subgraph $H^*$ of $H'$
with at least one vertex and minimum degree at least $d\ge2$. Let
$v_0v_1\cdots v_r$ be a longest path in $H^*$. Every neighbour of $v_0$ in $H^*$
is one of $v_1,\dots,v_r$, since otherwise the path could be extended. As $v_0$
has at least $d$ neighbours in $H^*$, the largest index $i$ with
$v_0v_i\in E(H^*)$ satisfies $i\ge d\ge2$, so $v_0v_1\cdots v_iv_0$ is a cycle of
$H'$ of length $i+1>d\ge\lceil m'/h\rceil\ge m'/h$. This proves the claim.

\emph{Removal of long cycles.} Put $H_0\defeq H$ and $m_i\defeq|E(H_i)|$. For
$i=0,1,\dots$, as long as $m_i\ge h$, choose by the claim (applied to
$H'\defeq H_i$) a cycle $C_i$ of $H_i$ of length at least $m_i/h$, and put
$H_{i+1}\defeq H_i-E(C_i)$. As $m_{i+1}<m_i$, this stops, at some $t\ge0$ with
$m_t\le h-1$. The cycles $C_0,\dots,C_{t-1}$ together with the $m_t$ edges of
$H_t$, taken as single edges, form a decomposition of $E(H)$ into $t+m_t$
objects, of which $m_t\le h-1$ are single edges. For $i<t$ we have
$m_{i+1}=m_i-|E(C_i)|\le(1-1/h)m_i$; hence $m_i\le(1-1/h)^im$ for $0\le i\le t$.

We use $(1-1/h)^h\le1/2$. For $h=1$ the left-hand side is $0$. For $h\ge2$ put
$x\defeq1/h$; then $(1-x)^h(1+x)^h=(1-x^2)^h\le1$ and, by the binomial theorem,
$(1+x)^h\ge1+hx=2$, so $(1-x)^h\le1/2$.

If $t=0$, the decomposition has $m_0\le h-1\le h(\log h+1)$ objects, as
$\log h\ge0$. Let $t\ge1$, and write $t-1=qh+r'$ with integers $q\ge0$ and
$0\le r'\le h-1$. Since $C_{t-1}$ was chosen, $m_{t-1}\ge h$; as $t-1\ge qh$ and
$0\le1-1/h\le1$,
\[
h\ \le\ m_{t-1}\ \le\ (1-1/h)^{t-1}m\ \le\ \bigl((1-1/h)^h\bigr)^qm\ \le\ 2^{-q}m\
\le\ 2^{-q}h^2/2 .
\]
Hence $2^{q+1}\le h$, i.e.\ $q+1\le\log h$, and $t=qh+r'+1\le(q+1)h\le h\log h$.
So the decomposition has $t+m_t\le h\log h+h-1\le h(\log h+1)$ objects.

(b) Note that $L>0$ as $N>1$. If $W=\emptyset$, then $F=\emptyset$, and the
empty decomposition has $0\le\beta\alpha$ objects. Otherwise $h\defeq|W|\ge1$,
and $h\le\beta\alpha/L\le\beta N/L\le N/4$ because $\alpha\le N$, $\beta>0$ and
$4\beta\le L$; hence $\log h\le\log N-2=L-2$. By (a), applied to the graph with
vertex set $W$ and edge set $F$, the set $F$ has a decomposition into at most
$h(\log h+1)\le h(L-1)\le hL\le\beta\alpha$ objects, at most $h-1\le|W|$ of
which are single edges.
\end{proof}

\subsection{Cited results of \BMnames{}}\label{s1:ssecCited}

Every published result of \cite{BM} used in this manuscript is stated below in the
manuscript's notation, with its number in \cite{BM}. Each statement was checked
against the plain-text extraction of arXiv:2211.07689v2; the footnote of each
item gives the line range. Where this manuscript modifies a proof of \cite{BM}
(Lemma~14, Propositions~12 and~13, Lemma~17, Proposition~18, Lemma~19, Lemma~9)
or extracts an explicit bound from a proof (Lemma~25), the relevant steps of the
proof are summarised as well; the modified proofs themselves are written out in
full where they are used. All items of this subsection are published results
(referee status ``pub'' in Table~\ref{s1:tabDAG}).

\begin{citedresult}[{\cite[Section~2.1, Section~2.6]{BM}}]\label{s1:citNotation}
\cite{BM} uses the following notation, which is also the notation of this
manuscript (Convention~\ref{s1:convGraphs}). For a graph $G$ and $v\in V(G)$,
$d_G(v)$ is the degree and $N_G(v)$ the neighbourhood of $v$; $\Delta(G)$ is the
maximum degree; for $U\subseteq V(G)$, $\Nbr_G(U)$ is the set of vertices of
$V(G)\setminus U$ with a neighbour in $U$; $G[V]$ is the induced subgraph,
$G\setminus V\defeq G[V(G)\setminus V]$, and $G-F$ is obtained by deleting the
edge set $F$; $|G|$ is the number of vertices. All logarithms have base~$2$;
$\log^{[0]}n=n$, $\log^{[k]}n$ is the $k$-fold iterated logarithm, and $\logs n$
is the least $k\ge0$ with $\log^{[k]}n\le1$. A path through $V$ is a path with
all internal vertices in $V$, and $\Ball{i}{G}{U,V}$ is the set of vertices of
$V$ that can be reached by a path through $V$ of length at most $i$ starting at a
vertex of $U$ (the starting vertex is not required to lie in $V$), so that
$\Ball{i}{G}{U,V}\subseteq V$.\bmcheck{123--145 and 379--385}
\deps{none (published).}
\end{citedresult}

\begin{citedresult}[{\cite[Theorem~4]{BM} (Chernoff)}]\label{s1:citChernoff}
Let $n$ be an integer, $0\le\delta,p\le1$, $X\sim\Bin(n,p)$ and $\mu\defeq\Ex
X=np$. Then
\[
\Prob\bigl(X>(1+\delta)\mu\bigr)\le e^{-\delta^2\mu/3}
\qquad\text{and}\qquad
\Prob\bigl(X<(1-\delta)\mu\bigr)\le e^{-\delta^2\mu/2}.
\]
\bmcheck{318--322}
\deps{none (published).}
\end{citedresult}

\begin{citedresult}[{\cite[Definition~7]{BM}}]\label{s1:citDef7}
A graph $G$ is \emph{$(\ell,t)$-path connected through} a vertex set $V\subseteq
V(G)$ if, for every collection $\mathcal P\subseteq\binom{V(G)}{2}$ in which every
vertex lies in at most $t$ pairs, there are edge-disjoint paths $P_{x,y}$,
$\{x,y\}\in\mathcal P$, such that each $P_{x,y}$ is an $xy$-path through $V$ of
length at most $\ell$.

\emph{Multiset clause} (quoted in substance from \cite{BM}): here
$\binom{V(G)}{2}$ denotes the \emph{multiset} of pairs of distinct vertices of
$G$; in particular the same pair may occur several times in $\mathcal P$. Thus a
pair occurring $k$ times receives $k$ pairwise edge-disjoint paths, and the
bound $t$ counts occurrences. The ends $x,y$ are unrestricted (they need not lie
in $V$); only interior vertices must lie in $V$.\bmcheck{341--355}
\deps{none (published).}
\fixes{\RTb{I14}.}
\end{citedresult}

\begin{citedresult}[{\cite[Proposition~8]{BM}}]\label{s1:citProp8}
Let $1\le\ell,t\le n$, let $G$ be an $n$-vertex graph and let $V\subseteq V(G)$
satisfy $|V|\ge4t-2$ and $|\Ball{\ell}{G}{U,V}|>|V|/2$ for every $U\subseteq
V(G)$ with $|U|=t$. If $x_1,\dots,x_{2t-1},y_1,\dots,y_{2t-1}$ are distinct
vertices of $G$, then for some $j\in[2t-1]$ there is an $x_jy_j$-path in $G$
through $V$ of length at most $4\ell\log n$.\bmcheck{386--392}
\deps{none (published).}
\end{citedresult}

\begin{citedresult}[{\cite[Lemma~9]{BM}}]\label{s1:citLem9}
Let $n\ge2$ and $1\le\ell,k\le n$. Let $G$ be an $n$-vertex graph and
$V\subseteq V(G)$ with $|V|\ge n/8+1$, and suppose that for every non-empty
$U\subseteq V(G)$ and every $F\subseteq E(G)$ with
$|F|\le2^9k|U|(\ell\log n)^2$ we have $|\Ball{\ell}{G-F}{U,V}|>|V|/2$. Then $G$
is $(4\ell\log n,k)$-path connected through $V$.

\emph{Structure of the proof in \cite{BM}} (modified in Lemma~\ref{s3:lemL9rho},
which uses Haxell's theorem, Cited result~\ref{s1:citHaxell}, in place of
Theorem~6).
Let $\mathcal P$ be a collection of pairs in which every vertex lies in at most
$k$ pairs, let $r\defeq|\mathcal P|$, and index the pairs by $i\in[r]$. Let
$H_i$ be the hypergraph on $E(G)$ whose edges are the edge sets of the paths
through $V$ of length at most $4\ell\log n$ joining the $i$-th pair. For
$I\subseteq[r]$ let $M_I$ be a maximal matching of $\bigcup_{i\in I}H_i$; the
claim is $|M_I|\ge4\ell\log n\cdot|I|$. If not, the set $F$ of edges covered by
$M_I$ has $|F|<(4\ell\log n)^2|I|$. Let $I'\subseteq I$ be maximal such that no
vertex lies in two pairs of $I'$; then $2k|I'|\ge|I|$ and $|I'|\le n/2$. With
$t\defeq\lceil|I'|/16\rceil$ one has $2t-1\le|I'|$, $4t-2\le\lceil n/8\rceil+1\le
|V|$ and $|F|<8kt(8\ell\log n)^2$, so every $t$-set $U$ satisfies
$|\Ball{\ell}{G-F}{U,V}|>|V|/2$; Proposition~8 (Cited
result~\ref{s1:citProp8}), applied to $G-F$ and $2t-1$ pairs of $I'$, gives a
path in $G-F$ joining a pair of $I'$, i.e.\ an edge of some $H_j$ disjoint from
$M_I$, contradicting maximality. Theorem~6 of \cite{BM} (Aharoni--Haxell
\cite{AH00}) then gives the paths. Since the pairs are indexed by $i\in[r]$, and the inequality
$2k|I'|\ge|I|$ counts the pairs of $I$ with multiplicity (every pair of $I$
meets a vertex of a pair of $I'$, and each of the $2|I'|$ vertices of pairs of
$I'$ lies in at most $k$ pairs of $I$), the argument applies verbatim to
multisets $\mathcal P$.\bmcheck{449--488}
\deps{none (published).}
\end{citedresult}

\begin{citedresult}[{\cite[Definition~11]{BM}}]\label{s1:citDef11}
An $n$-vertex graph $G$ is an \emph{$(\varepsilon,s)$-expander} if for every
$U\subseteq V(G)$ and every $F\subseteq E(G)$ with $1\le|U|\le2n/3$ and
$|F|\le s|U|$ we have
\[
|\Nbr_{G-F}(U)|\ge\frac{\varepsilon|U|}{\log^2n}.
\]
\emph{Remark (from \cite{BM}; hypotheses made explicit):} if $n\ge2$ and $\varepsilon>0$ then, for every
$v\in V(G)$, the pair $U=\{v\}$, $F=\{\text{edges at }v\}$ violates this
inequality, so necessarily $|F|>s$; hence $\delta(G)>s$ for every
$(\varepsilon,s)$-expander $G$ with $\varepsilon>0$ on at least two vertices.
(For $\varepsilon=0$ every graph satisfies the inequality; wherever this remark
is used below, $\varepsilon\ge2^{-7}$.)\bmcheck{579--594}
\deps{none (published).}
\end{citedresult}

\begin{citedresult}[{\cite[Proposition~12]{BM}}]\label{s1:citProp12}
For $U\subseteq V(G)$ and $d>0$ let $N_{G,d}(U)$ be the set of vertices of
$V(G)\setminus U$ with at least $d$ neighbours in $U$. Let $G$ be an $n$-vertex
$(\varepsilon,s)$-expander, $U\subseteq V(G)$ with $1\le|U|\le2n/3$, and $F$ a set
of at most $s|U|/2$ edges. Then for every $0<d\le s$, either (a)
$|\Nbr_{G-F}(U)|\ge s|U|/(2d)$, or (b) $|N_{G-F,d}(U)|\ge\varepsilon|U|/\log^2n$.

\emph{Proof in \cite{BM}} (its counting is modified in
Proposition~\ref{s3:propP13s}). If (a) fails, the set
$X\defeq\Nbr_{G-F}(U)\setminus N_{G-F,d}(U)$ has fewer than $s|U|/(2d)$
vertices, and the set $F'$ of edges of $G-F$ between $U$ and $X$ has
$|F'|\le d|X|\le s|U|/2$. Then $|F\cup F'|\le s|U|$ and
$\Nbr_{G-F-F'}(U)=N_{G-F,d}(U)$, so the expansion of $G$ gives
(b).\bmcheck{620--642}
\deps{none (published).}
\end{citedresult}

\begin{citedresult}[{\cite[Proposition~13]{BM}}]\label{s1:citProp13}
There is $n_0$ such that the following holds for $n\ge n_0$,
$1\ge\varepsilon\ge2^{-9}$ and $s\ge8\log^{13}n$. Let $G$ be an $n$-vertex
$(\varepsilon,s)$-expander, $U\subseteq V(G)$ with $|U|\le2n/3$, and $F$ a set of
at most $s|U|/4$ edges. Then $G-F$ contains either
\begin{enumerate}
\item[(a)] $|U|/\log^7n$ vertex-disjoint stars, each with $\log^9n$ leaves, with
centre in $U$ and all leaves in $V(G)\setminus U$, or
\item[(b)] a bipartite subgraph $H$ with classes $U$ and $X\subseteq
V(G)\setminus U$ such that $|X|\ge\varepsilon|U|/(2\log^2n)$, every vertex of
$X$ has degree at least $\log^4n$ in $H$, and every vertex of $U$ has degree at
most $2\log^9n$ in $H$.
\end{enumerate}
\emph{Structure of the proof in \cite{BM}} (modified in
Proposition~\ref{s3:propP13s}). Take a maximal family of vertex-disjoint stars as
in (a), with centre set $C\subseteq U$ and leaf set $L$. If (a) fails, then
$|C|\le|U|/\log^7n$, $|L|\le|U|\log^2n$, and by maximality no vertex of
$U\setminus C$ has $\log^9n$ neighbours in $G-F$ outside $U\cup L$, whence
$|\Nbr_{G-F}(U\setminus C)|<2|U|\log^9n$. With $d\defeq\log^4n$ and
$\Delta\defeq2\log^9n$, go through the vertices $v$ of $V(G)\setminus U$ in turn
and add a star with centre $v$ and $d$ leaves in $U$ whenever the degrees in $U$
stay at most $\Delta$; $X$ is the set of centres added and $H$ the union of the
stars. The saturated set $U'$ (vertices of $U\setminus C$ of $H$-degree exactly
$\Delta$) satisfies $|U'|\le2d|L|/\Delta\le\varepsilon|U|/(8\log^2n)$, because a
saturated vertex has at least $\Delta/2$ of its $H$-neighbours in $L$. With
$B\defeq C\cup U'$, Proposition~12 applied to $U\setminus B$ excludes option (a)
there, and every vertex of $N_{G-F,d}(U\setminus B)\setminus B$ lies in $X$,
giving $|X|\ge\varepsilon|U|/(2\log^2n)$.\bmcheck{643--713}
\deps{none (published).}
\end{citedresult}

\begin{citedresult}[{\cite[Lemma~14]{BM}}]\label{s1:citLem14}
Given an $n$-vertex graph $G$, a non-negative integer $s$ and
$\varepsilon\le2^{-5}$, one can delete at most $4sn\log n$ edges of $G$ so that
the remaining edges can be partitioned into graphs $G_1,\dots,G_r$ with
$\sum_{i=1}^r|G_i|\le2n$, each $G_i$ an $(\varepsilon,s)$-expander. (For $s=0$
nothing is deleted.)

\emph{Structure of the proof in \cite{BM}} (modified in
Lemma~\ref{s2:lem14tau}). The proof is by induction on $n$ under the stronger
bound $\sum_i|G_i|\le2n-2n/(2+\log n)$, which is at least $n$; the case $n=1$ is
trivial, and if $G$ is an $(\varepsilon,s)$-expander one takes $G_1=G$.
Otherwise there is a witness $(U,F)$: $1\le|U|\le2n/3$, $F\subseteq E(G)$,
$|F|\le s|U|$ and $|\Nbr_{G-F}(U)|<\varepsilon|U|/\log^2n$. Put
\[
G_1\defeq G[U\cup\Nbr_{G-F}(U)]-F,\qquad
G_2\defeq G\setminus U-E(G_1)-F,
\]
which edge-partition $G-F$, and $n_1\defeq|G_1|=|U\cup\Nbr_{G-F}(U)|$,
$n_2\defeq|G_2|=n-|U|$. Then
\begin{align*}
&\text{(6)}\quad n_1+n_2=n+|\Nbr_{G-F}(U)|<n+\varepsilon n_1/\log^2n;\\
&\text{(7)}\quad n_2<n\ \text{ and }\ n_1\le|U|+\varepsilon|U|/\log^2n\le
2n/3+\varepsilon n\le3n/4<n.
\end{align*}
The induction hypothesis gives sets $E_i\subseteq E(G_i)$ with $|E_i|\le4sn_i\log
n_i$ and partitions of $G_i-E_i$ into $(\varepsilon,s)$-expanders $G_{i,j}$ with
$\sum_j|G_{i,j}|\le2n_i-2n_i/(2+\log n_i)$. Using $\log n_1\le\log(3n/4)<\log
n-2/5$:
\begin{align*}
&\text{(8)}\quad |F|+|E_1|+|E_2|\le s\bigl(|U|+4n_1\log n_1+4n_2\log n_2\bigr)
\le s\bigl(4(n_1+n_2)\log n-\tfrac35n_1\bigr)\le4sn\log n,\\
&\text{(9)}\quad \frac{2n_1}{2+\log n_1}>\frac{2n_1}{2+\log
n}+\frac{n_1}{10\log^2n},
\end{align*}
where the last step of (8) uses (6). For $n\ge3$, (9) follows from $\log
n_1<\log n-2/5$ and $n_1<n$: then $\frac2{2+\log n_1}-\frac2{2+\log
n}>\frac{0.8}{(2+\log n)^2}\ge\frac1{10\log^2n}$, where the last step holds
because $\log n\ge\log3>2/(2\sqrt2-1)$. For $n=2$ we have $n_1=1$ by (7), and
(9) reads $1>\frac23+\frac1{10}$. Combining (9) with (6) gives
$\sum_{i,j}|G_{i,j}|\le(n_1+n_2)(2-\frac{2}{2+\log n})-\frac{n_1}{10\log^2n}\le
2n-\frac{2n}{2+\log n}$. The deleted edges are $F\cup E_1\cup E_2$.
\bmcheck{714--821}
\deps{none (published).}
\end{citedresult}

\begin{citedresult}[{\cite[Lemma~15]{BM}}]\label{s1:citLem15}
Let $n,k,s\in\mathbb N$ and $0<\varepsilon\le1$, and let $G$ be an $n$-vertex
$(\varepsilon,s)$-expander with $s\ge2^{12}\varepsilon^{-1}k^2\log^4n$. Then
there are edge-disjoint graphs $G_1,\dots,G_k$ with $E(G)=\bigcup_iE(G_i)$, each
an $\bigl(\varepsilon/4,\sqrt{s\varepsilon}/(8k\log n)\bigr)$-expander with vertex
set $V(G)$.

This result is quoted for comparison only and is \emph{not used}: it is replaced
by Lemma~\ref{s3:lemL15p} (Lemma~15$^+$), which is proved directly and loses
neither the factor $4$ in $\varepsilon$ nor the square root in
$s$.\bmcheck{822--836}
\deps{none (published).}
\end{citedresult}

\begin{citedresult}[{\cite[Theorem~16]{BM}}]\label{s1:citThm16}
Let $G$ be an $n$-vertex $(\varepsilon,s)$-expander with
$1\ge\varepsilon\ge2^{-7}$ and $s\ge\log^{135}n$, and let $V\subseteq V(G)$
contain each vertex independently with probability $1/3$. Then with high
probability $G$ is $(4\log^5n,2^8\log^5n)$-path connected through $V$.

This theorem is not used as a black box. Its proof (split $E(G)$ by Lemma~15
into $k=2^{17}\log^{42}n$ expanders; apply Lemma~19 in each; for a given $U$ and
$F$ with $|F|\le2^{17}|U|\log^{15}n$ some class contains at most
$|U|/\log^{27}n$ edges of $F$; conclude by Lemma~9 with $k=2^8\log^5n$ and
$\ell=\log^4n$) is modified in Theorem~\ref{s3:thmT16s}.\bmcheck{944--951 and
1327--1360}
\deps{none (published).}
\end{citedresult}

\begin{citedresult}[{\cite[Lemma~17]{BM}}]\label{s1:citLem17}
Let $n\ge2$ and let $G$ be an $n$-vertex $(\varepsilon,s)$-expander with
$2^{-9}\le\varepsilon\le1$ and $s\ge8\log^{13}n$. Let $U\subseteq V(G)$ satisfy
$|\Nbr_G(U)|\ge|U|\log^{24}n$, let $F\subseteq E(G)$ satisfy $|F|\le|U|$, and let
$V\subseteq V(G)$ contain each vertex independently with probability $1/3$. Then
with probability $1-e^{-\Omega(|U|\log^2n)}$ we have
$|\Ball{\log^4n}{G-F}{U,V}|>|V|/2$.

\emph{Proof parameters and structure in \cite{BM}} (modified in
Lemma~\ref{s3:lemL17s}). $\ell\defeq\log^4n$, $q\defeq3/11$, and $p$ is defined
by $(1-p)^{\ell-1}=11/12$, so that $p\ge1/(15\log^4n)$; $V$ is realised as
$V_1\cup\dots\cup V_\ell$ with independent $V_i$, each vertex in $V_i$ with
probability $p$ for $i<\ell$ and $q$ for $i=\ell$. $B_i$ is the set of vertices
reachable from $U$ by a path of $G-F$ of length at most $i$ whose interior lies
in $V_1\cup\dots\cup V_{i-1}$; it is determined by $U,V_1,\dots,V_{i-1}$, every
vertex of $\Nbr_{G-F}(B_i)$ with a $G-F$-neighbour in $B_i\cap V_i$ lies in
$B_{i+1}$, and $B_\ell\cap V_\ell\subseteq\Ball{\ell}{G-F}{U,V}$. The growth of
$B_i$ is shown via Proposition~13: in case (a) by Chernoff's bound on the star
centres, in case (b) by a bounded-differences inequality for the number of
vertices of $X$ with an $H$-neighbour in $V_i$; the final stage uses Chernoff's
bound with density $q$.\bmcheck{1046--1059 and 1068--1233}
\deps{none (published).}
\end{citedresult}

\begin{citedresult}[{\cite[Proposition~18]{BM}}]\label{s1:citProp18}
Let $n\ge2$, $0<\varepsilon\le1$ and $s\ge\log^{24}n$. Let $G$ be an $n$-vertex
$(\varepsilon,s)$-expander and $U\subseteq V(G)$ with $|U|\le2n/3$. Then there is
$U'\subseteq U$ with $|\Nbr_G(U')|\ge|U'|\log^{24}n$ and
$|U'|\ge\varepsilon|U|/(3\log^{26}n)$.

\emph{Proof in \cite{BM}} (modified in Lemma~\ref{s3:lemP18s}). Put
$\theta\defeq\log^{24}n$ and let $U'\subseteq U$ be maximal with
$|\Nbr_G(U')|\ge\theta|U'|$. Let $F$ be the set of edges from $U\setminus U'$ to
$V(G)\setminus(U'\cup\Nbr_G(U'))$. Expansion applied to $U$ and $F$ gives
$\varepsilon|U|/\log^2n\le|\Nbr_{G-F}(U)|\le|\Nbr_G(U')|<(|U'|+1)\theta+1\le
3|U'|\theta$.

\emph{Remark} (\RTb{I12}). The maximality step gives, for every $v\in U\setminus
U'$, \emph{fewer than $\theta+1$} neighbours outside $U'\cup\Nbr_G(U')$ (adding
$v$ to $U'$ removes at most the one vertex $v$ from the neighbourhood); \cite{BM}
write ``at most $\theta$'', which is exact only when $\theta$ is an integer.
Hence $|F|\le(\theta+1)|U\setminus U'|\le(\theta+1)|U|$, and the proof needs
$s\ge\theta+1$ rather than $s\ge\theta$. Lemma~\ref{s3:lemP18s} is formulated
with $s_1\ge\theta_*+1$ accordingly.\bmcheck{1234--1261}
\deps{none (published).}
\fixes{\RTb{I12}.}
\end{citedresult}

\begin{citedresult}[{\cite[Lemma~19]{BM}}]\label{s1:citLem19}
Let $G$ be an $n$-vertex $(\varepsilon,s)$-expander with
$2^{-9}\le\varepsilon\le1$ and $s\ge2\log^{24}n$, and let $V\subseteq V(G)$
contain each vertex independently with probability $1/3$. Then with probability
$1-o(1/n)$, for every non-empty $U\subseteq V(G)$ and every $F\subseteq E(G)$ with
$|F|\le|U|/\log^{27}n$, we have $|\Ball{\log^4n}{G-F}{U,V}|>|V|/2$. (The
statement in \cite{BM} says ``every $U$''; the empty set is excluded implicitly.)

\emph{Structure of the proof in \cite{BM}} (modified in
Lemma~\ref{s3:lemP18s}). Call $U'$ well-expanding if
$|\Nbr_G(U')|\ge|U'|\log^{24}n$. A union bound of the failure probability of
Lemma~17 over all pairs $(U',F)$ with $U'$ well-expanding and $|F|\le|U'|$ gives
failure probability $o(1/n)$. For $|U|\le2n/3$, Proposition~18 gives a
well-expanding $U'\subseteq U$ with $|U'|\ge\varepsilon|U|/(3\log^{26}n)\ge|F|$,
and the ball is monotone in $U$; for $|U|>2n/3$ one passes to a subset $\bar U$
with $n/2\le|\bar U|\le2n/3$.\bmcheck{1262--1326}
\deps{none (published).}
\end{citedresult}

\begin{citedresult}[{\cite[Theorem~21]{BM} (Lov\'asz \cite{Lov68})}]\label{s1:citThm21}
Every $n$-vertex graph can be decomposed into at most $n/2$ paths and
cycles.\bmcheck{1419--1422}
\deps{none (published).}
\end{citedresult}

\begin{citedresult}[{\cite[Corollary~22]{BM}}]\label{s1:citCor22}
Every graph can be decomposed into paths such that each vertex is an end of at
most two of the paths. (The proof in \cite{BM} adds a new vertex joined to every
vertex of even degree, applies Theorem~21, and deletes the new
vertex.)

\emph{Consequence used in this manuscript.} If every edge of a graph $H$ has both
ends in a set $W$, then $E(H)$ decomposes into at most $|W|$ paths. Indeed, apply
the corollary to the graph $(W,E(H))$: every path has two distinct ends in $W$
and every vertex of $W$ is an end of at most two paths, so twice the number of
paths is at most $2|W|$.\bmcheck{1423--1442}
\deps{none (published).}
\end{citedresult}

\begin{citedresult}[{\cite[Lemma~25]{BM}, explicit form}]\label{s1:citLem25}
Let $\varepsilon\ge2^{-5}$ and $m\ge\max(2,2^{30}/\varepsilon^2)$ (so
$m\ge2^{40}$ suffices at $\varepsilon=2^{-5}$). Every $m$-vertex
$(\varepsilon,0)$-expander contains a cycle of length at least
$\varepsilon^2m/(18\log^4m)$. (The bound $m\ge2$, which excludes the
one-vertex graph, follows from $m\ge2^{30}/\varepsilon^2$ unless
$\varepsilon^2>2^{29}$; without it the statement would be false for such
$\varepsilon$ at $m=1$.)

\emph{Explicit form from the proof in \cite{BM}.} The statement of \cite{BM}
says ``a cycle of length $\Omega(m/\log^4m)$''. The proof runs a depth-first
search, finds a moment at which the unexplored set $U$ and the processed set $R$
have equal size, so that the current path $P$ contains
$\Nbr_G(U)$ and has at least $\varepsilon m/(3\log^2m)$ vertices; it splits $P$
into consecutive segments $X,Y,Z$ with $|X|,|Z|\ge|P|/3$ and
$\varepsilon^2m/(18\log^4m)\le|Y|<\varepsilon^2m/(9\log^4m)$. If $X$ and $Z$
are joined by a path avoiding $Y$, a shortest such path together with a segment
of $P$ is a cycle containing $Y$, of length at least $|Y|\ge
\varepsilon^2m/(18\log^4m)$; otherwise expansion of the side of size at most
$m/2$ gives a contradiction. This explicit bound is the form used in
Lemma~\ref{s2:lemCap}.\bmcheck{1502--1558}
\deps{none (published).}
\end{citedresult}

\begin{remark}[where the results of \cite{BM} are used]\label{s1:remBMused}
\begin{center}
\small
\begin{tabular}{p{6.2cm}p{8.6cm}}
\textbf{Result of \cite{BM}} & \textbf{Used in}\\
\hline
Definition~7 (\ref{s1:citDef7}) & Sections~\ref{s3:sec}, \ref{s4:sec}, \ref{s5:sec}, \ref{s6:sec}\\
Lemma~9, Proposition~8 (\ref{s1:citLem9}, \ref{s1:citProp8}) & Lemma~\ref{s3:lemL9rho}\\
Theorem~6 (Aharoni--Haxell) & not used (Lemma~\ref{s3:lemL9rho} uses Haxell's theorem \cite{Hax95}, Cited result~\ref{s1:citHaxell}, instead)\\
Definition~11 (\ref{s1:citDef11}) & everywhere\\
Propositions~12--13, Lemma~17, Proposition~18, Lemma~19 & Section~\ref{s3:sec} (proof of Theorem~16$^*$)\\
Lemma~14 and its proof (\ref{s1:citLem14}) & Section~\ref{s2:sec}\\
Corollary~22 (\ref{s1:citCor22}) & Section~\ref{s4:sec}\\
Lemma~25 (\ref{s1:citLem25}) & Lemma~\ref{s2:lemCap}\\
Theorem~4 (\ref{s1:citChernoff}) & Sections~\ref{s3:sec}, \ref{s4:sec}\\
Lemma~15, Theorem~16, Theorem~21 & not used directly (Lemma~15 is replaced by Lemma~\ref{s3:lemL15p}; Theorem~16 by Theorem~\ref{s3:thmT16s}; Theorem~21 enters only through Corollary~22)\\
\hline
\end{tabular}
\end{center}
Theorem~2 of \cite{BM} (the bound $O(n\logs n)$) is not used, and neither are
Theorem~24 and Lemma~26 of \cite{BM}, which are used only in its proof: the vortex runs
of Section~\ref{s4:sec} are finished with the elementary long-cycle bound of
Fact~\ref{s1:factEG0}.
The three results of the row ``not used directly'' are not dependencies of this
remark.
\deps{\ref{s1:citNotation}, \ref{s1:citChernoff},
\ref{s1:citDef7}, \ref{s1:citProp8}, \ref{s1:citLem9}, \ref{s1:citDef11},
\ref{s1:citProp12}, \ref{s1:citProp13}, \ref{s1:citLem14},
\ref{s1:citLem17}, \ref{s1:citProp18}, \ref{s1:citLem19},
\ref{s1:citCor22}, \ref{s1:citLem25}, \ref{s1:factEG0}.}
\end{remark}

\subsection{Other cited and standard results}\label{s1:ssecOther}

The following standard results are used. Where a short derivation of the exact
form used here is needed, it is written out.
One of them, Lemma BBD below (a Bernstein inequality for functions of
independent indicators with bounded differences), is not cited but proved in
full. Its probabilistic part uses only finite sums. Its analytic part uses
only elementary properties of the exponential function (such as $e^v\ge1+v$
for real $v$) and the fact that a differentiable function whose derivative is
non-negative on an interval is non-decreasing there; the one elementary
one-variable inequality that it needs is proved there.

\begin{citedresult}[{Chernoff bounds for sums of independent indicators; \cite[Theorems~2.1 and~2.8]{JLR00}}]\label{s1:citChernoffGen}
Let $X=\sum_{i=1}^mI_i$ be a sum of independent indicator variables with
$\Prob(I_i=1)=p_i$ (a Poisson-binomial variable), and let $0\le\delta\le1$.
\begin{enumerate}
\item[(a)] If $\mu\ge\Ex X$ then $\Prob\bigl(X\ge(1+\delta)\mu\bigr)\le
e^{-\delta^2\mu/3}$; if $0\le\mu\le\Ex X$ then $\Prob\bigl(X\le(1-\delta)\mu\bigr)
\le e^{-\delta^2\mu/2}$. In particular the two bounds of Cited
result~\ref{s1:citChernoff} hold verbatim with $\mu=\Ex X$.
\item[(b)] If $\Ex X\le\mu$ then for every integer $j\ge1$,
\[
\Prob(X\ge j)\le\frac{\mu^j}{j!}\le\Bigl(\frac{e\mu}{j}\Bigr)^j.
\]
\end{enumerate}
\deps{none (standard).}
\end{citedresult}
\begin{proof}[Derivation of the forms used]
(a) By \cite[Theorems~2.1 and~2.8]{JLR00}, with $\lambda\defeq\Ex X$ and $t\ge0$,
$\Prob(X\ge\lambda+t)\le\exp\bigl(-t^2/(2(\lambda+t/3))\bigr)$ and
$\Prob(X\le\lambda-t)\le\exp\bigl(-t^2/(2\lambda)\bigr)$. For the upper tail put
$t\defeq(1+\delta)\mu-\lambda\ge\delta\mu$; the exponent
$t^2/(2(\lambda+t/3))$ is increasing in $t$ and, for fixed $\mu$, decreasing in
$\lambda$ (as $\lambda$ decreases, $t$ increases and
$\lambda+t/3=\frac23\lambda+\frac13(1+\delta)\mu$ decreases), so it is at least
its value at $\lambda=\mu$, namely $\delta^2\mu/(2(1+\delta/3))\ge\delta^2\mu/3$
since $\delta\le1$. For the lower tail write $\lambda=\mu+a$ with $a\ge0$ and
put $t\defeq\lambda-(1-\delta)\mu=\delta\mu+a\ge0$. Then
$t^2=\delta^2\mu^2+2\delta\mu a+a^2\ge\delta^2\mu^2+\delta^2\mu a=\delta^2\mu\lambda$
(as $\delta\le2$), so $t^2/(2\lambda)\ge\delta^2\mu/2$ (the case $\lambda=0$ is
trivial).
(b) If $X\ge j$ then some $j$-set $S\subseteq[m]$ has $I_i=1$ for all $i\in S$,
so $\Prob(X\ge j)\le\sum_{|S|=j}\prod_{i\in S}p_i\le(\sum_ip_i)^j/j!\le\mu^j/j!$,
because every product $\prod_{i\in S}p_i$ occurs $j!$ times in the expansion of
$(\sum_ip_i)^j$ and all terms are non-negative. Finally $j!\ge(j/e)^j$.
\end{proof}

\begin{citedresult}[{Markov's inequality and its two uses}]\label{s1:citMarkov}
If $X\ge0$ and $a>0$ then $\Prob(X\ge a)\le\Ex X/a$. Consequently:
\begin{enumerate}
\item[(a)] for $X\ge0$, $\Prob(X\le3\Ex X)\ge2/3$;
\item[(b)] for $X_1,X_2\ge0$, $\Prob(X_1\le4\Ex X_1\text{ and }X_2\le4\Ex X_2)\ge
1/2$;
\item[(c)] (a) and (b) hold with $\Prob$ and $\Ex$ replaced by the conditional
probability and expectation given any event of positive probability, or given
any fixed outcome of variables that are independent of the variables being
drawn.
\end{enumerate}
\deps{none (standard).}
\end{citedresult}
\begin{proof}[Derivation]
(a) If $\Ex X=0$ then $X=0$ almost surely. Otherwise
$\Prob(X>3\Ex X)\le\Prob(X\ge3\Ex X)\le1/3$.
(b) Similarly $\Prob(X_i>4\Ex X_i)\le1/4$ for $i=1,2$, and a union bound gives
failure probability at most $1/2$.
(c) The same argument in the conditional probability space; if the conditioning
fixes variables independent of those being drawn, the conditional law of the
latter is their unconditional law.
\end{proof}

\begin{lemma}[Lemma BBD: Bernstein inequality for bounded differences]\label{s1:lemBBD}
Let $m\ge0$ be an integer, let $p_1,\dots,p_m\in[0,1]$, and let
$I_1,\dots,I_m$ be independent random variables with $\Prob(I_k=1)=p_k$ and
$\Prob(I_k=0)=1-p_k$; write $I\defeq(I_1,\dots,I_m)$. Let
$\Psi\colon\{0,1\}^m\to\mathbb{R}$, and let $b\ge0$ and
$c_1,\dots,c_m\in[0,b]$ be such that $|\Psi(x)-\Psi(x')|\le c_k$ whenever
$x,x'\in\{0,1\}^m$ differ only in the $k$-th coordinate. Let $\beta>0$ with
$\beta\ge\sum_{k=1}^mp_k(1-p_k)c_k^2$. Then for every $a\ge0$,
\[
\Prob\bigl(\Psi(I)\ge\Ex\Psi(I)+a\bigr)\le\exp\Bigl(-\frac{a^2}{2(\beta+ba/3)}\Bigr)
\quad\text{and}\quad
\Prob\bigl(\Psi(I)\le\Ex\Psi(I)-a\bigr)\le\exp\Bigl(-\frac{a^2}{2(\beta+ba/3)}\Bigr).
\]
\deps{none (standard).}
\end{lemma}
\begin{proof}
This is a standard Bernstein-type inequality; we prove it by the
exponential-moment method. All sums in this proof are finite. For
$0\le k\le m$, $x=(x_1,\dots,x_k)\in\{0,1\}^k$ and
$z=(z_{k+1},\dots,z_m)\in\{0,1\}^{m-k}$ put
\[
\mathrm{wt}_k(x)\defeq\prod_{j=1}^{k}\bigl(p_jx_j+(1-p_j)(1-x_j)\bigr),\qquad
\mathrm{wt}'_k(z)\defeq\prod_{j=k+1}^{m}\bigl(p_jz_j+(1-p_j)(1-z_j)\bigr)
\]
(empty products are $1$), and write $(x,z)\in\{0,1\}^m$ for the
concatenation. These weights are non-negative,
$\sum_{x\in\{0,1\}^k}\mathrm{wt}_k(x)=\prod_{j\le k}\bigl(p_j+(1-p_j)\bigr)=1$, and
likewise $\sum_{z\in\{0,1\}^{m-k}}\mathrm{wt}'_k(z)=1$. By independence,
$\Prob(I=x)=\mathrm{wt}_m(x)$ for $x\in\{0,1\}^m$. Hence
$\Prob(I\in S)=\sum_{x\in S}\mathrm{wt}_m(x)$ for $S\subseteq\{0,1\}^m$, and
$\Ex g(I)=\sum_{x}\mathrm{wt}_m(x)g(x)$ for $g\colon\{0,1\}^m\to\mathbb{R}$.

\emph{Step 1: partial averages.} For $0\le k\le m$ and $x\in\{0,1\}^k$ put
\[
\Psi_k(x)\defeq\sum_{z\in\{0,1\}^{m-k}}\mathrm{wt}'_k(z)\,\Psi(x,z).
\]
Then
$\Psi_m=\Psi$, and, as $\mathrm{wt}'_0=\mathrm{wt}_m$, $\Psi_0$ is the constant
$\sum_x\mathrm{wt}_m(x)\Psi(x)=\Ex\Psi(I)$. Let $1\le k\le m$ and
$y\in\{0,1\}^{k-1}$. Split the sum defining $\Psi_{k-1}(y)$ according to
the first coordinate $z_k$ of $z=(z_k,z')$, and use
$\mathrm{wt}'_{k-1}(z_k,z')=\bigl(p_kz_k+(1-p_k)(1-z_k)\bigr)\mathrm{wt}'_k(z')$. This gives
\begin{equation}\label{s1:eqBBDavg}
\Psi_{k-1}(y)=p_k\Psi_k(y,1)+(1-p_k)\Psi_k(y,0).
\end{equation}
Put $\Xi_k(y)\defeq\Psi_k(y,1)-\Psi_k(y,0)
=\sum_{z}\mathrm{wt}'_k(z)\bigl(\Psi(y,1,z)-\Psi(y,0,z)\bigr)$. The points
$(y,1,z)$ and $(y,0,z)$ differ only in the $k$-th coordinate, so
$|\Xi_k(y)|\le\sum_z\mathrm{wt}'_k(z)c_k=c_k$. By~\eqref{s1:eqBBDavg},
\begin{equation}\label{s1:eqBBDdiff}
\Psi_k(y,x_k)-\Psi_{k-1}(y)=(x_k-p_k)\,\Xi_k(y)\qquad(x_k\in\{0,1\}).
\end{equation}

\emph{Step 2: an elementary inequality.} For every real $u$,
\begin{equation}\label{s1:eqBBDphi}
\Bigl(1-\frac u3\Bigr)e^u\le1+\frac{2u}3+\frac{u^2}6 .
\end{equation}
Indeed, let $\chi(u)$ be the right side minus the left side. Then
$\chi(0)=0$ and
$\chi'(u)=\frac23+\frac u3-\bigl(\frac23-\frac u3\bigr)e^u$, so
$\chi'(0)=0$; and $\chi''(u)=\frac13\bigl(1-(1-u)e^u\bigr)\ge0$,
because $1-u\le e^{-u}$. Hence $\chi'$ is non-decreasing, so
$\chi'\le0$ on $(-\infty,0]$ and $\chi'\ge0$ on $[0,\infty)$. Thus
$\chi$ is non-increasing on $(-\infty,0]$ and non-decreasing on
$[0,\infty)$, and $\chi\ge\chi(0)=0$. Now let $0\le\zeta<3$ and $u\le\zeta$.
Then $1-u/3\ge1-\zeta/3>0$ and
$1+\frac{2u}3+\frac{u^2}6=\bigl(1-\frac u3\bigr)(1+u)+\frac{u^2}2$. Dividing
\eqref{s1:eqBBDphi} by $1-u/3$ gives
\begin{equation}\label{s1:eqBBDexp}
e^u\le1+u+\frac{u^2}{2(1-u/3)}\le1+u+\frac{u^2}{2(1-\zeta/3)} .
\end{equation}

\emph{Step 3: one coordinate.} Let $\eta\ge0$ with $\eta b<3$ (it is chosen
in Step~5), and put $\zeta\defeq\eta b$. Let $1\le k\le m$ and
$y\in\{0,1\}^{k-1}$, and write $p\defeq p_k$ and $\Xi\defeq\Xi_k(y)$.
The numbers $\eta(1-p)\Xi$ and $-\eta p\Xi$ are at most
$\eta|\Xi|\le\eta c_k\le\eta b=\zeta$. Apply~\eqref{s1:eqBBDexp} to both,
multiply by $p\ge0$ and $1-p\ge0$ respectively, and add. Since
$p(1-p)-(1-p)p=0$ and $p(1-p)^2+(1-p)p^2=p(1-p)$, this gives
\begin{equation}\label{s1:eqBBDstep}
p\,e^{\eta(1-p)\Xi}+(1-p)\,e^{-\eta p\Xi}
\le1+\frac{\eta^2p(1-p)\Xi^2}{2(1-\zeta/3)}
\le\exp\Bigl(\frac{\eta^2p_k(1-p_k)c_k^2}{2(1-\zeta/3)}\Bigr),
\end{equation}
where the last step uses $\Xi^2\le c_k^2$ and $1+v\le e^v$.

\emph{Step 4: induction over the coordinates.} For $0\le k\le m$ put
\[
\mathcal{E}_k\defeq\sum_{x\in\{0,1\}^k}\mathrm{wt}_k(x)\exp\bigl(\eta(\Psi_k(x)-\Psi_0)\bigr).
\]
Then $\mathcal{E}_0=1$, and $\mathcal{E}_m=\Ex\exp\bigl(\eta(\Psi(I)-\Ex\Psi(I))\bigr)$.
Let $1\le k\le m$, and write $x\in\{0,1\}^k$ as $(y,x_k)$ with
$y\in\{0,1\}^{k-1}$; then $\mathrm{wt}_k(y,1)=p_k\mathrm{wt}_{k-1}(y)$ and
$\mathrm{wt}_k(y,0)=(1-p_k)\mathrm{wt}_{k-1}(y)$. By~\eqref{s1:eqBBDdiff}, and then
by~\eqref{s1:eqBBDstep} and the non-negativity of all weights,
\begin{align*}
\mathcal{E}_k&=\sum_{y\in\{0,1\}^{k-1}}\mathrm{wt}_{k-1}(y)\,e^{\eta(\Psi_{k-1}(y)-\Psi_0)}
\Bigl(p_ke^{\eta(1-p_k)\Xi_k(y)}+(1-p_k)e^{-\eta p_k\Xi_k(y)}\Bigr)\\
&\le\mathcal{E}_{k-1}\exp\Bigl(\frac{\eta^2p_k(1-p_k)c_k^2}{2(1-\zeta/3)}\Bigr).
\end{align*}
By induction on $k$, and since $\sum_kp_k(1-p_k)c_k^2\le\beta$,
$\mathcal{E}_m\le\exp\bigl(\eta^2\beta/(2(1-\zeta/3))\bigr)$.

\emph{Step 5: the upper tail.} Let
$S\defeq\{x\in\{0,1\}^m:\Psi(x)\ge\Psi_0+a\}$. For $x\in S$ we have
$\exp\bigl(\eta(\Psi(x)-\Psi_0-a)\bigr)\ge1$, as $\eta\ge0$, and all terms
below are non-negative. Hence
\begin{align*}
\Prob\bigl(\Psi(I)\ge\Ex\Psi(I)+a\bigr)&=\sum_{x\in S}\mathrm{wt}_m(x)
\le e^{-\eta a}\sum_{x\in\{0,1\}^m}\mathrm{wt}_m(x)e^{\eta(\Psi(x)-\Psi_0)}\\
&=e^{-\eta a}\mathcal{E}_m
\le\exp\Bigl(-\eta a+\frac{\eta^2\beta}{2(1-\eta b/3)}\Bigr).
\end{align*}
Now choose $\eta\defeq a/(\beta+ba/3)\ge0$. Since $\beta>0$, we have
$\eta b/3=(ba/3)/(\beta+ba/3)<1$, so $\eta b<3$ as required in Step~3, and
$1-\eta b/3=\beta/(\beta+ba/3)$. Therefore
\[
-\eta a+\frac{\eta^2\beta}{2(1-\eta b/3)}
=-\frac{a^2}{\beta+ba/3}+\frac{a^2}{(\beta+ba/3)^2}\cdot\frac{\beta\,(\beta+ba/3)}{2\beta}
=-\frac{a^2}{2(\beta+ba/3)},
\]
which is the first bound.

\emph{Step 6: the lower tail.} The function $-\Psi$ satisfies the hypotheses
with the same $b$, $c_1,\dots,c_m$ and $\beta$, and
$\Ex(-\Psi(I))=-\Ex\Psi(I)$. So the first bound for $-\Psi$ is the second
bound for $\Psi$.
\end{proof}

\begin{citedresult}[{Hall's theorem \cite{Hal35}}]\label{s1:citHall}
Let $\mathcal J$ be a finite set and, for every $a\in\mathcal J$, let
$\mathcal T(a)$ be a finite set. If
$\bigl|\bigcup_{a\in\mathcal J'}\mathcal T(a)\bigr|\ge|\mathcal J'|$ for every
$\mathcal J'\subseteq\mathcal J$, then there is an injective map $\chi$ from
$\mathcal J$ to $\bigcup_{a\in\mathcal J}\mathcal T(a)$ with
$\chi(a)\in\mathcal T(a)$ for every $a\in\mathcal J$ (a system of distinct
representatives).
\deps{none (standard).}
\end{citedresult}

\begin{citedresult}[{Haxell's condition for matchability \cite{Hax95}}]\label{s1:citHaxell}
Let $q\ge1$ be an integer, and let $\mathcal H$ be a hypergraph (a set of
subsets, called edges, of a finite vertex set) whose vertex set is the
disjoint union of two sets $\mathcal X$ and $\mathcal Y$, such that every edge
$e$ satisfies $|e\cap\mathcal X|=1$ and $|e\cap\mathcal Y|\le q$. Suppose that
for every $\mathcal X'\subseteq\mathcal X$ and every $Z\subseteq\mathcal Y$ with
$|Z|\le(2q-1)(|\mathcal X'|-1)$ there is an edge $e$ of $\mathcal H$ with
$e\cap\mathcal X\subseteq\mathcal X'$ and $e\cap Z=\emptyset$. Then
$\mathcal H$ has an \emph{$\mathcal X$-saturating matching}, i.e.\ a set of
pairwise disjoint edges such that every vertex of $\mathcal X$ lies in one of
them.

For $\mathcal X'=\emptyset$ there is no such $Z$, because $(2q-1)(0-1)<0$. For
$q=1$ (edges with at most two vertices) the result is Hall's theorem; a vertex
$a$ with $\{a\}\in\mathcal H$ is matched by that edge. In the later literature
the result of \cite{Hax95} is quoted for hypergraphs whose edges have at most
(in some sources: exactly) $r$ vertices, one of them in $\mathcal X$, with the
bound $(2r-3)(|\mathcal X'|-1)$. With $r=q+1$, the maximum edge size, this bound
is $(2(q+1)-3)(|\mathcal X'|-1)=(2q-1)(|\mathcal X'|-1)$: the stated form agrees
with Haxell's theorem as quoted there, and the case $q=1$ (Hall's theorem)
confirms this reading of the parameters. The statement, including the factor
$2q-1$, has not yet been checked against the text of \cite{Hax95}
(Table~\ref{s1:tabDAG}). The only application, in the proof of
Lemma~\ref{s3:lemL9rho}, verifies the hypothesis for every $Z$ with
$|Z|<q^2|\mathcal X'|$; so that application stays valid if $2q-1$ is replaced
by any positive factor at most $q^2$. Unlike the other cited items of this
subsection, which are standard textbook results, this is a published research
result (status ``pub.'' in Table~\ref{s1:tabDAG}). The derivation below shows
that the stated form follows from the special case in which every edge meets
$\mathcal Y$ in exactly $q$ vertices; so the stated form holds if \cite{Hax95}
proves either of the two quoted forms. After the derivation, a complete proof of
the stated form is written out (Haxell's alternating-tree argument), so that
the application does not rest on the unchecked citation; this proof was added
after the clean-room review G0 and has had no referee yet.
\deps{none (published).}
\end{citedresult}
\begin{proof}[Derivation from the case $|e\cap\mathcal Y|=q$]
Let $\mathcal H$ be as stated, and let $\mathcal X_0$ be the set of
$a\in\mathcal X$ such that $\{a\}$ is an edge. Every edge $e$ with
$e\cap\mathcal X\not\subseteq\mathcal X_0$ has $e\cap\mathcal Y\ne\emptyset$
(otherwise $e=\{a\}$ with $a\in\mathcal X_0$). For each such $e$ take a set
$N_e$ of $q-|e\cap\mathcal Y|$ new vertices, the sets $N_e$ pairwise disjoint
and disjoint from $\mathcal X\cup\mathcal Y$, and put $e^+:=e\cup N_e$. Let
$\mathcal H^+$ be the hypergraph on
$(\mathcal X\setminus\mathcal X_0)\sqcup\mathcal Y^+$, where
$\mathcal Y^+:=\mathcal Y\cup\bigcup_eN_e$, whose edges are these sets $e^+$.
Every edge of $\mathcal H^+$ meets $\mathcal X\setminus\mathcal X_0$ in exactly
one vertex and $\mathcal Y^+$ in exactly $q$ vertices.

We check the hypothesis for $\mathcal H^+$. Let
$\mathcal X'\subseteq\mathcal X\setminus\mathcal X_0$ and
$Z^+\subseteq\mathcal Y^+$ with $|Z^+|\le(2q-1)(|\mathcal X'|-1)$. For every
$z\in Z^+\setminus\mathcal Y$ let $f_z$ be the unique edge of $\mathcal H$ with
$z\in N_{f_z}$, and choose a vertex $b_z\in f_z\cap\mathcal Y$. Put
$Z:=(Z^+\cap\mathcal Y)\cup\{b_z:z\in Z^+\setminus\mathcal Y\}\subseteq\mathcal Y$.
Then $|Z|\le|Z^+|$, so the hypothesis for $\mathcal H$ gives an edge $e$ with
$e\cap\mathcal X\subseteq\mathcal X'$ and $e\cap Z=\emptyset$. As
$|e\cap\mathcal X|=1$, the set $e\cap\mathcal X$ is a single vertex of
$\mathcal X'\subseteq\mathcal X\setminus\mathcal X_0$; so
$e\cap\mathcal X\not\subseteq\mathcal X_0$, the set $e^+$ is an edge of
$\mathcal H^+$, and
$e^+\cap(\mathcal X\setminus\mathcal X_0)=e\cap\mathcal X\subseteq\mathcal X'$.
Moreover $e^+\cap Z^+=\emptyset$: a vertex of $e^+\cap Z^+$ lies either in
$e\cap\mathcal Y$, and then in $Z\cap e$; or it is some $z\in N_e$, and then
$f_z=e$ and $b_z\in e\cap Z$. Both contradict $e\cap Z=\emptyset$.

By the case $|e\cap\mathcal Y|=q$, $\mathcal H^+$ has an
$(\mathcal X\setminus\mathcal X_0)$-saturating matching
$\{e_1^+,\dots,e_k^+\}$. The edges $e_1,\dots,e_k$ of $\mathcal H$ are pairwise
disjoint, as $e_i\subseteq e_i^+$, and they cover
$\mathcal X\setminus\mathcal X_0$, as
$e_i\cap\mathcal X=e_i^+\cap(\mathcal X\setminus\mathcal X_0)$. The edges
$\{a\}$, $a\in\mathcal X_0$, are disjoint from each other and from every $e_i$,
since $e_i\cap\mathcal X\subseteq\mathcal X\setminus\mathcal X_0$. So
$\{e_1,\dots,e_k\}\cup\{\{a\}:a\in\mathcal X_0\}$ is an $\mathcal X$-saturating
matching of $\mathcal H$.
\end{proof}
\begin{proof}[Proof of the stated form]
This proof was added after G0 and has had no referee yet. It is Haxell's
alternating-tree argument. For an edge $e$ let $a(e)$ denote
the vertex of $e\cap\mathcal X$. We use induction on $|\mathcal X|$. If
$\mathcal X=\emptyset$, the empty set of edges is an $\mathcal X$-saturating
matching. Let $|\mathcal X|\ge1$ and fix $a_0\in\mathcal X$. The edges of
$\mathcal H$ that do not contain $a_0$ form a hypergraph on
$(\mathcal X\setminus\{a_0\})\sqcup\mathcal Y$ that satisfies the hypothesis
with $\mathcal X\setminus\{a_0\}$ in place of $\mathcal X$: for
$\mathcal X'\subseteq\mathcal X\setminus\{a_0\}$, the edge $e$ given by the
hypothesis for $\mathcal H$ has $e\cap\mathcal X\subseteq\mathcal X'$, so
$a_0\notin e$. By the induction hypothesis this hypergraph has an
$(\mathcal X\setminus\{a_0\})$-saturating matching.

Call a set $\mathcal M$ of pairwise disjoint edges of $\mathcal H$ \emph{good} if
no edge of $\mathcal M$ contains $a_0$ and every vertex of
$\mathcal X\setminus\{a_0\}$ lies in an edge of $\mathcal M$. Since every edge
meets $\mathcal X$ in exactly one vertex, the map $m\mapsto a(m)$ is then a
bijection from $\mathcal M$ onto $\mathcal X\setminus\{a_0\}$. We have just seen
that a good set exists. It suffices to find a good $\mathcal M$ and an edge $x$
with $a(x)=a_0$ and $x\cap m\cap\mathcal Y=\emptyset$ for every
$m\in\mathcal M$: then $x$ is disjoint from every $m\in\mathcal M$ (as
$a(m)\ne a_0$), and $\mathcal M\cup\{x\}$ is an $\mathcal X$-saturating
matching.

\emph{Configurations.} Let $\mathcal M$ be good. A \emph{configuration} for
$\mathcal M$ is a sequence $(x_1,\mathcal S_1,\dots,x_k,\mathcal S_k)$ with
$k\ge0$, of edges $x_i$ of $\mathcal H$ and sets
$\mathcal S_i\subseteq\mathcal M$, with the following properties. Put
$\mathcal X^{(0)}:=\{a_0\}$, $Z^{(0)}:=\emptyset$ and, for $1\le i\le k$,
\[
\mathcal X^{(i)}:=\{a_0\}\cup\{a(m):m\in\mathcal S_1\cup\dots\cup\mathcal S_i\},
\qquad
Z^{(i)}:=\mathcal Y\cap\Bigl(x_1\cup\dots\cup x_i\cup
\bigcup_{m\in\mathcal S_1\cup\dots\cup\mathcal S_i}m\Bigr).
\]
The properties are, for every $1\le i\le k$: (C1) $a(x_i)\in\mathcal X^{(i-1)}$ and
$x_i\cap Z^{(i-1)}=\emptyset$; (C2)
$\mathcal S_i=\{m\in\mathcal M:m\cap x_i\cap\mathcal Y\ne\emptyset\}$ and
$\mathcal S_i\ne\emptyset$. The \emph{signature} of the configuration is the
sequence $(|\mathcal S_1|,\dots,|\mathcal S_k|)$ of positive integers.

The sets $\mathcal S_i$ are pairwise disjoint: if $j<i$ and
$m\in\mathcal S_j$, then $m\cap\mathcal Y\subseteq Z^{(i-1)}$, which $x_i$
avoids by (C1), so $m\notin\mathcal S_i$. As $m\mapsto a(m)$ is injective on
$\mathcal M$ and $a(m)\ne a_0$, we get
$|\mathcal X^{(k)}|-1=\sum_{i\le k}|\mathcal S_i|\le|\mathcal M|=|\mathcal X|-1$.
Each $m\in\mathcal S_i$ shares a vertex of $\mathcal Y$ with $x_i$, so
$|(m\cap\mathcal Y)\setminus x_i|\le q-1$. Hence, using $|\mathcal S_i|\ge1$,
\begin{align*}
|Z^{(k)}|&\le\sum_{i\le k}\Bigl(|x_i\cap\mathcal Y|
+\sum_{m\in\mathcal S_i}|(m\cap\mathcal Y)\setminus x_i|\Bigr)
\le\sum_{i\le k}\bigl(q+(q-1)|\mathcal S_i|\bigr)\\
&\le(2q-1)\sum_{i\le k}|\mathcal S_i|=(2q-1)\bigl(|\mathcal X^{(k)}|-1\bigr).
\end{align*}
So the hypothesis, with $\mathcal X':=\mathcal X^{(k)}$ and $Z:=Z^{(k)}$, gives
an edge $x$ with $a(x)\in\mathcal X^{(k)}$ and $x\cap Z^{(k)}=\emptyset$. $(*)$

\emph{Swap step.} Suppose that $\mathcal M$ is good, that
$(x_1,\mathcal S_1,\dots,x_j,\mathcal S_j)$ is a configuration for
$\mathcal M$, and that $y$ is an edge with $a(y)\in\mathcal X^{(j)}$,
$y\cap Z^{(j)}=\emptyset$ and $y\cap m\cap\mathcal Y=\emptyset$ for every
$m\in\mathcal M$. If $a(y)=a_0$, then $\mathcal M$ and $x:=y$ are as required,
and the swap step ends. Otherwise $a(y)=a(m)$ for exactly one $m\in\mathcal M$, and
$m\in\mathcal S_i$ for exactly one $i\le j$. Put
$\mathcal M':=(\mathcal M\setminus\{m\})\cup\{y\}$. It is good: $y$ shares no
vertex of $\mathcal Y$ with an edge of $\mathcal M$, the vertex $a(y)=a(m)$
lies in no edge of $\mathcal M\setminus\{m\}$, $a_0\notin y$, and $a(m)$ is
covered by $y$. For $i'\le j$ we have $x_{i'}\cap\mathcal Y\subseteq Z^{(j)}$,
so $y\cap x_{i'}\cap\mathcal Y=\emptyset$; and $m\in\mathcal S_{i'}$ only for
$i'=i$. Hence
$\{m'\in\mathcal M':m'\cap x_{i'}\cap\mathcal Y\ne\emptyset\}$ equals
$\mathcal S_{i'}$ for $i'<i$ and $\mathcal S_i\setminus\{m\}$ for $i'=i$, and
the sets $\mathcal X^{(i')}$ and $Z^{(i')}$ with $i'<i$ are the same for
$\mathcal M'$ as for $\mathcal M$. So, if $\mathcal S_i\ne\{m\}$, then
$(x_1,\mathcal S_1,\dots,x_{i-1},\mathcal S_{i-1},x_i,\mathcal S_i\setminus\{m\})$
is a configuration for $\mathcal M'$, and the swap step ends. If $\mathcal S_i=\{m\}$,
then $(x_1,\mathcal S_1,\dots,x_{i-1},\mathcal S_{i-1})$ is a configuration for
$\mathcal M'$, the edge $x_i$ satisfies $a(x_i)\in\mathcal X^{(i-1)}$ and
$x_i\cap Z^{(i-1)}=\emptyset$ by (C1), and $x_i$ shares no vertex of
$\mathcal Y$ with an edge of $\mathcal M'$; we repeat the swap step with
$\mathcal M'$, this configuration and $y:=x_i$. As $i-1<j$, and as
$\mathcal X^{(0)}=\{a_0\}$, the swap step ends after at most $j+1$ rounds.

\emph{One move.} Let $\mathcal M$ be good, with a configuration of signature
$(s_1,\dots,s_k)$, and let $x$ be as in $(*)$. If $x$ shares a vertex of
$\mathcal Y$ with some edge of $\mathcal M$, append $x$ and
$\mathcal S_{k+1}:=\{m\in\mathcal M:m\cap x\cap\mathcal Y\ne\emptyset\}$; by
$(*)$ this is a configuration for $\mathcal M$, of signature
$(s_1,\dots,s_k,s_{k+1})$. Otherwise apply the swap step with $j:=k$ and
$y:=x$. Either it finds $\mathcal M$ and $x$ as required, or it ends with a
good $\mathcal M'$ and a configuration for $\mathcal M'$ of signature
$(s_1,\dots,s_{i-1},s_i-1)$ for some $i\le k$ with $s_i\ge2$.

\emph{Termination.} For finite sequences $\sigma,\tau$ of positive integers
write $\sigma\prec\tau$ if the word $\sigma\infty$ precedes the word
$\tau\infty$ in the lexicographic order on words over
$\{1,2,\dots\}\cup\{\infty\}$. As $\infty$ occurs only as the last letter,
neither of two distinct such words is a prefix of the other, so $\prec$ is a
strict total order. A move that does not find $\mathcal M$ and $x$ as required
replaces the signature $\sigma$
by some $\sigma'\prec\sigma$: an extension $\sigma'$ of $\sigma$ has a finite
entry where $\sigma\infty$ has $\infty$, and $(s_1,\dots,s_{i-1},s_i-1)$ is
smaller at position $i$. Every signature has positive entries with sum at most
$|\mathcal X|-1$, so there are finitely many signatures, and a sequence that
is strictly decreasing for $\prec$ has no repetition; hence it is finite.
Start with a good $\mathcal M$ and the configuration with $k=0$, and make
moves as long as possible. After finitely many moves a swap step finds a good
$\mathcal M$ and an edge $x$ as required.
\end{proof}

\begin{citedresult}[{$T$-joins and Euler circuits}]\label{s1:citEuler}
\begin{enumerate}
\item[(a)] Let $H$ be a connected graph and $T\subseteq V(H)$ with $|T|$ even.
Every spanning tree $S$ of $H$ contains a $T$-join, i.e.\ a set $J\subseteq E(S)$
such that the vertices of odd $J$-degree are exactly the vertices of $T$.
\item[(b)] A connected graph (or multigraph) in which every vertex has even degree
has a closed trail using every edge exactly once.
\item[(c)] A graph in which every vertex has even degree is an edge-disjoint union
of cycles; a directed graph without loops in which every vertex has in-degree
equal to out-degree is an arc-disjoint union of directed cycles.
\end{enumerate}
The Lov\'asz path decomposition is used only through Cited
result~\ref{s1:citCor22}.
\deps{none (standard).}
\end{citedresult}
\begin{proof}[Derivation of (a) and (c)]
(a) Root $S$ at a vertex $\rho$, and put the edge from a vertex $v\ne\rho$ to its
parent into $J$ if and only if the subtree of $S$ rooted at $v$ contains an odd
number of vertices of $T$. For $v\ne\rho$ the $J$-degree of $v$ is congruent
modulo $2$ to the number of $T$-vertices in its subtree minus those in the
subtrees of its children, i.e.\ to $\ind{v\in T}$; for $\rho$ it is congruent to
$|T|-\ind{\rho\in T}\equiv\ind{\rho\in T}$, since $|T|$ is even. (b) is Euler's
theorem. (c) If some edge (arc) remains, start at one of its ends and walk along
unused edges (leaving along unused out-arcs); since every vertex has even degree
(balanced degrees), the walk can always continue until it revisits a vertex,
which closes a cycle (directed cycle); delete it and repeat.
\end{proof}

\subsection{Constants, their order, and the galactic conditions}\label{s1:ssecConstants}

\begin{definition}[the constants]\label{s1:defConstants}
\begin{enumerate}
\item[(i)] $\eps\defeq2^{-5}$, $\sigma\defeq100$, $\Cp\defeq103$ and
$\Aexp\defeq105$.
\item[(ii)] $\Nzero$ is an absolute constant, i.e.\ it depends only on $\sigma$,
$\Cp$ and $\Aexp$, with $\Nzero\ge2^{40}$ (redundant: size condition~(i),
$L\ge2^{10}$, already forces $\Nzero\ge2^{1024}$) and $\Nzero$ at least
each of the absolute size thresholds required in the proofs of
Lemma~\ref{s4:lemTPV}, Lemma~\ref{s4:lemPV} and Theorem~\ref{s4:thmVXp}.
Every requirement on $\Nzero$ in this item is an \emph{eventuality}: it holds
for every sufficiently large value of $\Nzero$. Indeed $\Nzero\ge2^{40}$ is a
lower bound, and every other requirement asks that an explicit inequality in
$N$, which holds for all sufficiently large $N$, hold for every $N\ge\Nzero$. So
such an $\Nzero$ exists (finitely many eventualities have a common threshold),
and only its existence is used; no numerical value of $\Nzero$ enters any proof.
(Each of these three statements carries the hypothesis that its vertex set has
at least $\Nzero$ vertices, and its proof records the size conditions it
needs. The lemmas of Section~\ref{s5:sec} carry the hypothesis $n\ge\Nzero$
through the setting paragraph of that section and use $\Nzero$ only to apply
these three statements. Lemma~\ref{s3:lemL15p} and Theorem~\ref{s3:thmT16s}
need no lower bound on the number of vertices.)
\item[(iii)] $\Dstar$ is a constant satisfying \ref{s1:condG1}--\ref{s1:condG4}
of Condition~\ref{s1:condGamma}. Each of these conditions is an eventuality in
$\Dstar$: once $\Nzero$ is fixed, every sufficiently large $\Dstar$ satisfies all
of them (Lemma~\ref{s7:lemGammaSat}). So such a constant exists; only its
existence is used, and no numerical value of $\Dstar$ enters any proof.
\item[(iv)] $\epsOne(\Dstar)$ and $\epsTwo(\Dstar)$ are the explicit functions
of $\Dstar$ defined in Proposition~\ref{s7:propCost} and
Lemma~\ref{s7:lemUHsplit} respectively, and $\thetaQ\defeq\epsTwo(\Dstar)$. Each is a finite sum of explicit terms
depending only on $\Dstar$, $\eps$, $\sigma$, $\Cp$ and $\Aexp$.
\item[(v)] $\Czero\defeq\Dstar/2+\cTot+\epsOne(\Dstar)$; the function
$\epsOne$ is counted exactly once. By \ref{s1:condG4},
$\Czero\le\Dstar/2+\cTotR$.
\item[(vi)] $\cEG\defeq\max\bigl\{\Czero/(1-2\thetaQ),\ \Nzero/2\bigr\}$.
\item[(vii)] (Generic notation, proofs only; permitted, but not used in the
present text.) Inside a proof, $\epsD$ may be
used as shorthand for an expression written out explicitly in the same proof
that depends only on $\Dstar$, $\eps$, $\sigma$, $\Cp$, $\Aexp$ and tends to $0$ as
$\Dstar\to\infty$. It never appears in the final bound of a statement; every
named small quantity ($\epsA$, $\epsU$, $\epsK=\epsChain$, $\epsCONC$,
$\varepsilon_M$, $\epsX$, $\epsOne$, $\epsTwo$) is given by an explicit formula
in its defining statement.
\end{enumerate}
\deps{none.}
{\small\noindent(The references of this definition to later statements are not
dependencies; see the paragraph on forward references in
Section~\ref{s1:ssecDAG}. For instance, $\epsOne$ and $\epsTwo$ are defined in
Proposition~\ref{s7:propCost} and Lemma~\ref{s7:lemUHsplit}.)\par}
\fixes{\RTb{J6}.}
\end{definition}

\begin{remark}[order of the constants]\label{s1:remOrder}
The constants are fixed in the following order:
\begin{enumerate}
\item $\sigma$, $\Cp$, $\Aexp$, $\eps$;
\item $\Nzero$;
\item $\Dstar$;
\item $\epsOne(\Dstar)$, $\epsTwo(\Dstar)=\thetaQ$ and $\Czero$;
\item $\cEG$.
\end{enumerate}
Nothing earlier in this list depends on anything later. In particular $\Czero$
and $\thetaQ$ depend on $\Dstar$ only (given the fixed $\sigma,\Cp,\Aexp,\eps$):
they do not depend on $G$, on $n$, on the run of the hierarchy, on the
designation, or on $\cEG$. The constant $\cEG$ enters only as the constant $c$
of the induction, Theorem~\ref{s7:thmHI}, which is applied last; this is why
that induction is not circular. The constants of steps~2 and~3 are only required to be sufficiently
large: every condition on $\Nzero$ (Definition~\ref{s1:defConstants}(ii)) and on
$\Dstar$ (Condition~\ref{s1:condGamma}) is an eventuality in that constant, given
the constants of the earlier steps. Table~\ref{s1:tabOrder}, after
Condition~\ref{s1:condGamma}, lists every such condition with the variable in
which it is eventual, the constants it may depend on, and the values at which it
is used.
\deps{\ref{s1:defConstants}.}
\end{remark}

\begin{condition}[the galactic conditions on $\Dstar$]\label{s1:condGamma}
The constant $\Dstar$ satisfies the following conditions. This is the single list
of galactic conditions of the manuscript; every use cites the specific item.
Every condition is an \emph{eventuality} in $\Dstar$: once $\Nzero$ is fixed, it
holds for all sufficiently large $\Dstar$ (Lemma~\ref{s7:lemGammaSat}). Only the
existence of a suitable $\Dstar$ is used: no proof evaluates any of the
inequalities below at a particular numerical value.
\begin{enumerate}
\item[\Gc{1}]\namedlabel{s1:condG1}{\Gc{1}} (\emph{Tower condition.})
$\Dstar>2$, so that $\log_2\log_2\Dstar$ is defined and positive; and for every real
$\mu\ge\log_2\log_2\Dstar$, with $\lambda\defeq2^\mu$, the following hold:
\begin{enumerate}
\item[(a)] $\mu\ge2^8$;
\item[(b)] $2^\mu\ge2^{14}\Aexp\mu^3$;
\item[(c)] $2\Aexp\log_2(\Aexp\mu)\le1.6\,\mu$;
\item[(d)] $2\log_2\mu+8\le\mu$ (this follows from (a) and~(b), and is listed
separately for direct citation);
\item[(e)]\namedlabel{s1:condGstar}{\Gstar} (the condition \Gstar)
$\lambda^{36}\ge2^{240}(\Aexp\mu)^{46\Aexp}$;
\item[(f)] for every row of the COL-JV table (Table~\ref{s3:tabCOLJV}), every
inequality in column~3 of that row holds at $\lambda$, where
$\bar M\defeq(\Aexp\mu)^{2\Aexp}$ and $\bar k\defeq192\lambda^3+\frac43\bar M^2$ as in
the caption of that table. (Row~13 is \Gstar, and the inequality of row~12 is
item~(c). This item points forward to explicit formulas, as \Gc{4} does for
$\epsOne$ and $\epsTwo$: column~3 consists of $18$ explicit inequalities, two in
row~1, five in row~8 and one in each other row, and they involve only $\mu$,
$\lambda=2^\mu$, $\bar M$ and $\bar k$.)
\end{enumerate}
Each item consists of explicit inequalities in the single real variable $\mu$,
and each of them holds for all sufficiently large $\mu$
(Lemma~\ref{s7:lemGammaSat}); \Gc{1} asks that
all of them hold on the whole ray $\mu\ge\log_2\log_2\Dstar$. So \Gc{1} is upward
closed in $\Dstar$. (The requirement $\Dstar>2$ only makes the ray well
defined: whenever $\log_2\log_2\Dstar$ is defined, (a) at
$\mu=\log_2\log_2\Dstar$ already gives $\Dstar\ge2^{2^{256}}$.)

Every round $r\le R$ of the hierarchy has average degree $d_r\ge\Dstar$
(Definition~\ref{s2:defHBtp}), so $\lambda_r\defeq\log_2d_r$ satisfies
$\log_2\lambda_r\ge\log_2\log_2\Dstar$, and (a)--(f) hold at
$\mu=\log_2\lambda_r$. They also hold at $\mu=\log_2\log_2d$ for every
$d\ge\Dstar$, and at $\mu=\log_2\Dstar$ (as $t\ge\log_2t$ for $t>0$). By (a) at
$\mu=\log_2\log_2\Dstar$, $\log_2\Dstar\ge2^{256}$. The consequences used later,
each proved where it is used, are: every requirement (column~2) and every
column-3 inequality of the COL-JV table (Lemma~\ref{s3:lemCOLJV}(ii)); the tower
facts of Proposition~\ref{s2:propDegRec}, Lemma~\ref{s2:lemLacunary} and
Lemma~\ref{s2:lemTower} (among them $\thGC_r>\tau_r$ and
$\lambda_r\ge2^{2\logs d_r+2}$); and Proposition~\ref{s2:propStructure} and
elementary estimates in Sections~\ref{s5:sec}--\ref{s7:sec}, each of which cites
the item it uses.

\item[\Gc{2}]\namedlabel{s1:condG2}{\Gc{2}}
\begin{enumerate}
\item[(a)] $\Dstar\ge2^{117}$; hence $\log_2(\tHB_ld_l)\ge117$ in every round
$l$, the fixed-point requirement of the Lemma-25 size cap
(Lemma~\ref{s2:lemCap});
\item[(b)] $P_{l-2}/2\ge M_l\log^4M_l$ for every $3\le l\le R$, including rounds
with $M_l=2^{40}$ (light parts two rounds back are admissible parents,
Proposition~\ref{s2:propParentless});
\item[(c)] every per-ancestor failure probability of Lemma~\ref{s3:lemCOL} and
of Lemma~\ref{s5:lemE1} is at most $|V(Y)|^{-2}$ (the failure-union lines of
Table~\ref{s3:tabCOLJV}).
\end{enumerate}
Each of (a)--(c) is implied by \ref{s1:condG1}: (a) is
immediate from \ref{s1:condG1}; (b) is proved in Lemma~\ref{s2:lemTower}(b); and
the bounds of (c) are proved in Lemma~\ref{s3:lemCOL} (with rows~6 and~7 of
Table~\ref{s3:tabCOLJV}) and in Lemma~\ref{s5:lemE1}(b),(c). None of these
proofs, and none of the results they use, uses (b) or (c). They are listed separately for traceability.

\item[\Gc{3}]\namedlabel{s1:condG3}{\Gc{3}} $(\log_2\Dstar)^{\Cp}\ge2\Nzero$.

Hence every pre-part has at least $P_R\ge(\log\Dstar)^{\Cp}\ge2\Nzero$ vertices,
every light part at least $P_R/2\ge\Nzero$, every ancestor, part and TPV,
P-vortex or VX$^+$ run meets the absolute size thresholds of its tools, and every
$o(1)$ term in a per-part success probability is at most $1/100$ (each proved
where used).

\item[\Gc{4}]\namedlabel{s1:condG4}{\Gc{4}} $\epsOne(\Dstar)\le6$ and
$\thetaQ=\epsTwo(\Dstar)\le1/4$, where $\epsOne$ and $\epsTwo$ are the explicit
expressions of Proposition~\ref{s7:propCost} and Lemma~\ref{s7:lemUHsplit}.

These expressions collect the following terms, with the coefficients stated;
here $\epsOne$ contains $3\,\epsX$ and $\epsTwo$ contains $12\,\epsX$, where
$\epsX=\epsU+10.3/\Dstar+2.1(6\log_2\Dstar+12)/\Dstar$
(Lemma~\ref{s7:lemEXprime}):
\begin{itemize}
\item the concentration term $1.5\,\epsCONC(\Dstar)$ of $\epsOne$ (it does
not occur in $\epsTwo$), where
\[
\epsCONC(\Dstar)=2^{\sigma+15}\frac{\logs\Dstar}{\log\Dstar}
+200\,\eps\,\frac{\logs\Dstar}{\Cp\log\log\Dstar}
+\frac{4(2\logs\Dstar+2)}{(\log\Dstar)^{1/2}}
\]
(Theorem~\ref{s6:thmCONCL}(iv); its first two terms bound the $\tau$- and
$d^*$-terms of all ancestors as in Theorem~\ref{s6:thmCONC}(iii), with a spare
factor $2$, and its last term bounds the guest-cap terms $\thGC_r(Y^0)$ of the
light parts);
\item $\epsA=31\eps/(\Cp\log\log\Dstar)$ (defined in
Lemma~\ref{s2:lemTower}(e)), with coefficient $169$ in $\epsOne$ and $12$ in
$\epsTwo$;
\item $\epsU$, the demotion and lost-parent expectations
(Lemma~\ref{s5:lemExpect}), with coefficient $3$ in $\epsOne$ and $12$ in
$\epsTwo$ (through $\epsX$);
\item $\epsK$, the U-chaining term (Lemma~\ref{s5:lemKRED}), with
coefficient $1$ in $\epsOne$ (it does not occur in $\epsTwo$);
\item the pool term
$\sum_l(6\log_2M_l+12)/M_l\le2.1(6\log_2\Dstar+12)/\Dstar$, whose bound
enters with coefficient $3$ in $\epsOne$ and $12$ in $\epsTwo$ (through
$\epsX$);
\item in $\epsTwo$ only, the quotient term $2F(\log_2\Dstar)$, with
coefficient $4$ (so $\epsTwo$ contains $8F(\log_2\Dstar)$), where
$F(x)=(3184+30400\log_2x)\,x^{-90.2}$ (Lemma~\ref{s7:lemUHsplit}(iv));
\item terms $O(1/\Dstar)$: $(252+9+3\cdot10.3)/\Dstar=291.9/\Dstar$ in
$\epsOne$ and $(12\cdot10.3+4\cdot60)/\Dstar=363.6/\Dstar$ in $\epsTwo$.
\end{itemize}
\end{enumerate}

\emph{Remarks.} (1) \emph{Eventualities.} Each of
\ref{s1:condG1}--\ref{s1:condG4} holds for all sufficiently large $\Dstar$, once
$\Nzero$ is fixed: \ref{s1:condG1} because each of its finitely many items holds
for all sufficiently large $\mu$ (and its range $\mu\ge\log_2\log_2\Dstar$
shrinks as $\Dstar$ grows); \ref{s1:condG2}(a) and \ref{s1:condG3} because they
are lower bounds on $\Dstar$, and \ref{s1:condG2}(b),(c) because they follow
from \ref{s1:condG1}; and \ref{s1:condG4} because
$\epsOne(\Dstar)\to0$ and $\epsTwo(\Dstar)\to0$ as $\Dstar\to\infty$. This is
proved in Lemma~\ref{s7:lemGammaSat}. Hence every sufficiently large $\Dstar$
satisfies all four conditions, and only the existence of such a $\Dstar$ is
used.

(2) \emph{Order of choice; no circularity.} Table~\ref{s1:tabOrder} lists the
constants in the order in which they are chosen (Remark~\ref{s1:remOrder}), the
conditions on them, the variable in which each condition is an eventuality, and
the constants it may depend on. A condition on the constant of a step depends
only on constants of earlier steps: the size conditions, \ref{s1:condG1},
\ref{s1:condG2}(a) and \ref{s1:condG4} depend on absolute constants only, and
\ref{s1:condG3} depends on $\Nzero$, which is chosen before $\Dstar$. Every
condition is used only at values of its variable that are at least the chosen
threshold: the size conditions at $N=|Z|$ for the vertex set $Z$ of a vortex
run, which is a pre-part or a light part and so has $N\ge P_R/2\ge\Nzero$ by
\ref{s1:condG3}; the items of \ref{s1:condG1} at values
$\mu\ge\log_2\log_2\Dstar$, as described there; and \ref{s1:condG2}(a),
\ref{s1:condG3} and \ref{s1:condG4} at $\Dstar$ itself. The quantities of a
run ($d_r$, $\lambda_r$, $|V(Y)|$, $L_Y$) are not chosen: they are universally
quantified, and they exceed the thresholds because $d_r\ge\Dstar$ in every round
$r\le R$ and by \ref{s1:condG3}. The hypotheses $n\ge\Nzero$ and
$d_1\ge\Dstar$ on the graph are not consequences of these conditions: a graph
with $n<\Nzero$ or with $d_1<\Dstar$ is handled directly in steps~(2) and~(3) of
the proof of Theorem~\ref{s7:thmHI}, because $\cEG\ge\Nzero/2$ and
$\cEG\ge\Czero\ge\Dstar/2$. Finally
$\epsOne$, $\epsTwo=\thetaQ$ and $\Czero$ are functions of $\Dstar$, and $\cEG$
is defined last; it is used only as the constant $c$ of
Theorem~\ref{s7:thmHI}, and no condition on $\Nzero$ or $\Dstar$ involves it.
So no choice is circular.

(3) Conditions of earlier formulations that are \emph{not used} on this route:
the condition written (G) in the standalone-junction analysis
($\lambda^{31}\ge2^{200}(\Aexp\log\lambda)^{46\Aexp}$), a lower bound on
$\log_2\lambda$ required by the colouring lemma for zone-set classes, and the
conditions (G$^+$) and (G$_{\mathrm{XL}}$) (lower bounds on
$\log_2\log_2\Dstar$ and on $\log_2\lambda$). They belong only to discarded
routes (Remark~\ref{s1:remNotUsed}) and are never assumed.
\deps{\ref{s1:defConstants}, \ref{s1:remOrder}.}
\end{condition}

\begin{table}[htbp]
\begin{center}
\small
\begin{tabular}{|p{0.75cm}|p{4.6cm}|p{2.3cm}|p{2.2cm}|p{3.4cm}|}
\hline
Step & Constant; conditions on it & Eventual in & Depends on & Used at\\\hline
1 & $\eps=2^{-5}$, $\sigma=100$, $\Cp=103$, $\Aexp=105$ & explicit & -- & everywhere\\\hline
2 & $\Nzero$: $\Nzero\ge2^{40}$ (redundant: size condition~(i), $L\ge2^{10}$, already forces $\Nzero\ge2^{1024}$), and the size conditions of Lemma~\ref{s4:lemTPV}, Lemma~\ref{s4:lemPV} and Theorem~\ref{s4:thmVXp} hold for every $N\ge\Nzero$ & $N$ (so every requirement is upward closed in $\Nzero$) & absolute constants only & $N=|Z|$ for the vertex set $Z$ of a vortex run; $N\ge\Nzero$ by \Gc{3}\\\hline
3 & $\Dstar$: \Gc{1}, i.e.\ $\Dstar>2$ and (a)--(f), including \Gstar{} and the column-3 inequalities of Table~\ref{s3:tabCOLJV} & $\mu$, on the ray $\mu\ge\log_2\log_2\Dstar$; hence $\Dstar$ & absolute constants only & $\mu=\log_2\lambda_r$ ($r\le R$), $\log_2\log_2d$ ($d\ge\Dstar$; e.g.\ $\log_2x$ for $x\ge\log_2\Dstar$), $\log_2\Dstar$\\\cline{2-5}
 & \Gc{2}(a): $\Dstar\ge2^{117}$ & $\Dstar$ & absolute constants only & $d_l\ge\Dstar$ (Lemma~\ref{s2:lemCap})\\\cline{2-5}
 & \Gc{2}(b),(c) & consequences of \Gc{1} & absolute constants only; not a condition & every round, every ancestor\\\cline{2-5}
 & \Gc{3}: $(\log_2\Dstar)^{\Cp}\ge2\Nzero$ & $\Dstar$ & $\Nzero$ (step~2) & sizes of parts, $P_R/2\ge\Nzero$\\\cline{2-5}
 & \Gc{4}: $\epsOne(\Dstar)\le6$ and $\epsTwo(\Dstar)\le1/4$ & $\Dstar$ & absolute constants only & $\Dstar$ (cost line, Section~\ref{s7:sec}; Definition~\ref{s1:defConstants}(v); Remark~\ref{s5:remConstants}(b))\\\hline
4 & $\epsOne(\Dstar)$, $\thetaQ=\epsTwo(\Dstar)$, $\Czero$ & defined & $\Dstar$ & Theorem~\ref{s7:thmHI}\\\hline
5 & $\cEG=\max\{\Czero/(1-2\thetaQ),\Nzero/2\}$ & defined & $\Czero$, $\thetaQ$, $\Nzero$ & Theorem~\ref{s7:thmHI}, as its constant $c$\\\hline
\end{tabular}
\end{center}
\caption{The order in which the constants are chosen, and the conditions on
them (Condition~\ref{s1:condGamma}, Remark~(2)). ``Eventual in $x$'' means:
holds for all sufficiently large $x$ once the constants of the earlier steps are
fixed. A condition on the constant of a step depends only on earlier steps. For
steps~2 and~3, the last column gives the values at which a condition is used;
each of them is at least the chosen threshold.}
\label{s1:tabOrder}
\end{table}

\begin{remark}[not used on this route]\label{s1:remNotUsed}
The following objects, statements and devices of earlier formulations of the
argument are \emph{not used} anywhere in this manuscript, and nothing here depends
on them:
\begin{itemize}
\item the zone sets ZS and the notion of rootedness built on them (every classed
port uses the expander of its own class, $H_{Y(u)}$), the routes XL, QI and
QI$^+$, the NZ functionals, the split and parity functionals SP$_\delta$ and
HC$_\delta$;
\item the star split of JS-LC and its concentration threshold $\tau'_l$ (the
cherry split is used for every giant pair, Remark~\ref{s6:remStar}); VX-parts
(here $\mathcal V=\emptyset$: every standalone pre-part is a MIX-part);
\item the Hall-type statements HALL and HALL$'$ (not to be confused with Hall's
theorem on systems of distinct representatives, a cited result of this
section that is used), and the ``corrected Theorem JV''
proposed by an earlier AI referee; these are superseded and never cited
(\RTa{M11}, \RTa{M12});
\item the collision lemma COLL and the proposition RAND-JV$^+$; the payments
(D5$'$) (concentrated pairs) and (D7$'$) (ultra excess); the rooted-first
deletion RULE; the functional $W$ and the coefficient $32/15$ (\RTa{M4});
\item the rooted-hub device RH, Theorem C and the calibration families.
\end{itemize}
Red team~2's items \RTb{I15} (the project's own ledger was stale) and
\RTb{I16} (an exponent in an off-route vortex variant) concern project
bookkeeping and an off-route statement; they do not apply to this manuscript and
are recorded here only for completeness.
\deps{none.}
\fixes{\RTa{M4}, \RTa{M11}, \RTa{M12}, \RTb{I15}, \RTb{I16}.}
\end{remark}

\subsection{The symbol table}\label{s1:ssecSymbols}

Every symbol that is used outside a single proof has exactly one meaning, given in
Table~\ref{s1:tabSymbols}. A symbol listed there is not reused locally. Symbols
local to one proof or one run are introduced there and have no meaning elsewhere:
the sans-serif symbols of a vortex run (Convention~\ref{s4:convVortex}); the
starred parameters of the proof of Theorem~\ref{s3:thmT16s}; the node notation of
Lemma~14$^\tau$ and Lemma SEP (Section~\ref{s2:sec}); the cluster, layer,
demand, padding and potential notation of Lemmas~\ref{s6:lemMED}
and~\ref{s6:lemHCCP}; and the letters $S_w$, $\mu$, $\mu^*$, $T$, $m_w$ of
Lemmas~\ref{s7:lemCC} and~\ref{s7:lemUltra}. The letter $\eta$ is deliberately
local as well and is not in the table: it denotes the failure terms
$\eta_{\mathrm{TPV}}$, $\eta_{\mathrm{PV}}$ and $\eta_{\mathrm{VX}}$ of the three
vortex statements of Section~\ref{s4:sec}; the bijections $\eta_h$ of
Definition~\ref{s7:defSchedule} and Construction~\ref{s7:consRound}; the
functions $\eta(x)$ of Lemma~\ref{s2:lemLacunary}(iv) and of the proof of
Lemma~\ref{s5:lemParent}(ii) (the latter also in the overview); and numbers
inside single proofs (for instance those of Lemmas~\ref{s1:lemBBD}
and~\ref{s6:lemJSLC}). The column ``macro'' gives the
\LaTeX{} macro of the preamble (``--'' if the symbol has none), and the column
``defined in'' points to the defining statement (forward pointers are allowed in
this table).

\begin{center}
\refstepcounter{table}\label{s1:tabSymbols}
\textbf{Table~\thetable.} The symbol table.\\[4pt]
{\small
\setlength{\tabcolsep}{4pt}
\begin{tabular}{p{3.7cm}p{3.0cm}p{5.6cm}p{2.5cm}}
\textbf{Symbol} & \textbf{Macro} & \textbf{Meaning} & \textbf{Defined in}\\
\hline
\multicolumn{4}{l}{\emph{General}}\\
$\f(F)$, $\f(G)$, $\f(n)$ & \texttt{\string\f} & least number of objects (cycles, single edges) partitioning $F$, $E(G)$; maximum over $n$-vertex $G$ & \ref{s1:defObject}\\
$n$ & -- & number of vertices of the input graph $G$ & \ref{s1:convGraphs}\\
$\log$, $\logs$ & \texttt{\string\log}, \texttt{\string\logs} & base $2$; iterated logarithm & \ref{s1:convGraphs}\\
$\Nbr_G(U)$, $\Ball{i}{G}{U,V}$ & \texttt{\string\Nbr}, \texttt{\string\Ball} & outside neighbourhood; ball through $V$ & \ref{s1:citNotation}\\
$\Prob$, $\Ex$, $\ind{\cdot}$, $\Bin$ & \texttt{\string\Prob}, \texttt{\string\Ex}, \texttt{\string\ind}, \texttt{\string\Bin} & probability, expectation, indicator, binomial & --\\
\hline
\multicolumn{4}{l}{\emph{Constants}}\\
$\eps=2^{-5}$; $\sigma=100$; $\Cp=103$; $\Aexp=105$ & \texttt{\string\eps}, --, \texttt{\string\Cp}, \texttt{\string\Aexp} & fixed parameters & \ref{s1:defConstants}\\
$\Nzero$, $\Dstar$, $\Czero$, $\cEG$ & \texttt{\string\Nzero}, \texttt{\string\Dstar}, \texttt{\string\Czero}, \texttt{\string\cEG} & size threshold; galactic degree threshold; main constants & \ref{s1:defConstants}\\
$\epsOne$, $\epsTwo=\thetaQ$ & \texttt{\string\epsOne}, \texttt{\string\epsTwo}, \texttt{\string\thetaQ} & cost and quotient-size functions of $\Dstar$ & \ref{s7:propCost}, \ref{s7:lemUHsplit}\\
$\epsA$, $\epsU$, $\epsK=\epsChain$, $\epsCONC$, $\varepsilon_M$, $\epsX$ & \texttt{\string\epsA}, \texttt{\string\epsU}, \texttt{\string\epsK}, \texttt{\string\epsChain}, \texttt{\string\epsCONC}, --, \texttt{\string\epsX} & named small functions of $\Dstar$ & \ref{s2:lemTower}, \ref{s5:lemExpect}, \ref{s5:lemParent}/\ref{s5:lemKRED}, \ref{s6:thmCONCL}, \ref{s6:thmMIXC}, \ref{s7:lemEXprime}\\
$\epsD$ & \texttt{\string\epsD} & generic, inside proofs only & \ref{s1:defConstants}(vii)\\
\Gc{1}--\Gc{4}, \Gstar & \texttt{\string\Gc}, \texttt{\string\Gstar} & galactic conditions & \ref{s1:condGamma}\\
$\cOV=\log_2(4/3)$ & \texttt{\string\cOV} & overlap constant & \ref{s2:lemOVgeneric}\\
$\psi(x)=(6\log_2x+12)/x$ & -- & tower sum weight & \ref{s2:lemTower}\\
\hline
\end{tabular}}
\end{center}

\begin{center}
{\small
\setlength{\tabcolsep}{4pt}
\begin{tabular}{p{3.7cm}p{3.0cm}p{5.6cm}p{2.5cm}}
\multicolumn{4}{l}{(Table~\ref{s1:tabSymbols}, continued.) \emph{Hierarchy}}\\
\hline
$\HBtp$ ($\HBs$, $\HBt$ only in comparisons) & \texttt{\string\HBtp} (\texttt{\string\HBs}, \texttt{\string\HBt}) & the hierarchy & \ref{s2:defHBtp}\\
$G_l$, $G'_l$, $d_l$, $\lambda_l=\log d_l$, $R$, $E_0$, $\Cyc_l$ & --, \texttt{\string\Cyc} & round-$l$ graph, the same after removing long cycles; average degree; its logarithm; last round; leftover; long cycles & \ref{s2:defHBtp}\\
$\tHB_l=\log^2d_l$; $M_l\in\mathbb N$ (a ceiling, (R2)); $\Lambda_l=\log M_l$; $s_l=\lceil\Lambda_l^\sigma\rceil$; $P_l=\lceil\lambda_l^{\Cp}\rceil$; $\tau_l=\lceil128s_l\log^2M_l\rceil$ & \texttt{\string\tHB} & round parameters & \ref{s2:defHBtp}\\
$\piece$ & \texttt{\string\piece} & piece of the $s=0$ run of Lemma~14 & \ref{s2:defHBtp}\\
$Z^0$, $X^0_Z$, $S_Z$, $\home(v)$, $D_l$ & \texttt{\string\home} & pre-part, its leaf expander, guests, home, duplicated vertices & \ref{s2:defHBtp}\\
light part $Z=Z^0\setminus S_Z$, $X_Z=X^0_Z-S_Z$; $\Std_l$; GC-part; $\thGC_l(Z^0)=\lceil|Z^0|\lambda_l^{-1/2}\rceil$ & \texttt{\string\Std}, \texttt{\string\thGC} & light/standalone split; guest cap & \ref{s2:defHBtp}\\
$E_l(Z)$ & -- & edges assigned at round $l$ to $Z$ & \ref{s2:defHBtp}\\
witness $(U,F)$, $N$, $F_0$, $\Hout$, $U'$, $N'$, $F_1$, $\Hin$, $N''$, $F''$, $G_1$, $G_2$ & \texttt{\string\Hout}, \texttt{\string\Hin} & $\tau$-rules (local to Section~\ref{s2:sec}) & \ref{s2:defWitness}, \ref{s2:defTauRules}\\
leaf $\Leaf$, $\Dup$, $\Dups_l$ & \texttt{\string\Leaf}, \texttt{\string\Dup}, \texttt{\string\Dups} & leaf of a recursion; vertices in at least two leaves (of one recursion; of round $l$'s two-level recursion) & \ref{s2:lemSEP}, \ref{s2:defHBtp}\\
ancestor $Y$; $V(Y)$; $H_Y$; $\eps_Y$; $s_Y$; $L_Y=\log|V(Y)|$ & -- & ancestor of round $r$ and its expander & \ref{s2:defAncestors}\\
$\anc_l(x)$; $A_Z$; $U_Z$; $F_Z$; $Q_Z$ & \texttt{\string\anc} & ancestors; hubs; ports; fresh ports; classed ports & \ref{s2:defAncestors}\\
$\mu_r(w)$, $\dup_r$, $\mult_r(w)$ & \texttt{\string\dup}, \texttt{\string\mult} & number of round-$r$ pre-parts containing $w$; duplication; number of round-$r$ ancestors containing $w$ & \ref{s2:defAncestors}\\
$j_0(x)$, $j_1(x)$ & -- & first round in which $x$ lies in a pre-part; in a light part & \ref{s2:defAncestors}\\
$\nu_l$ & -- & number of ancestors of rounds $\le l-2$ (at most $2.74n/P_{l-2}$) & \ref{s2:propOV} (bound: \ref{s2:lemTower}(b))\\
$(\alpha)$, $(\beta)$ & -- & ORIGIN edge types & \ref{s2:propOrigin}\\
\hline
\end{tabular}}
\end{center}

\begin{center}
{\small
\setlength{\tabcolsep}{4pt}
\begin{tabular}{p{3.7cm}p{3.0cm}p{5.6cm}p{2.5cm}}
\multicolumn{4}{l}{(Table~\ref{s1:tabSymbols}, continued.) \emph{Colouring, lending, zones, light parts}}\\
\hline
$\Own_Y$, $\Lend_Y$, $\klend(Y)$, $\kown$ & \texttt{\string\Own}, \texttt{\string\Lend}, \texttt{\string\klend}, \texttt{\string\kown} & two-stage colouring & \ref{s3:defCOL}\\
$\IU$, $\IJS$, $\IJV$ & \texttt{\string\IU}, \texttt{\string\IJS}, \texttt{\string\IJV} & lent index families & \ref{s3:defCOL}\\
$\LU_{Y,l,c,\uslot}$, $\LJS_{Y,l,j}$, $\LJV_{Y,l}$ & \texttt{\string\LU}, \texttt{\string\LJS}, \texttt{\string\LJV} & U-, JS- and JV-lent classes & \ref{s3:defCOL}\\
$\KJS_l=M_l^2$; $\tJS_l=2M_l+2$; $\rho_l=M_l^{-4}$; $\lab_{Y,l}(y)$; $T_j(Y,l)$ & \texttt{\string\KJS}, \texttt{\string\tJS}, \texttt{\string\lab} & JS layers, multiplicity, label probability, labels, junction sets & \ref{s3:defCOL}\\
$p_Y=1/(2\klend(Y))$ & -- & probability that an edge of $H_Y$ lies in $\LJV_{Y,l}$ & \ref{s3:defCOL}\\
$\Tslot_Y=\lceil L_Y^2\rceil$ & \texttt{\string\Tslot} & number of U-slots & \ref{s3:defCOL}\\
$t_Y=\lceil\lambda_r^{1.6}\rceil$ & -- & U-multiplicity parameter of a light ancestor $Y$ of round $r$ & \ref{s3:defCOL}\\
$\Aw(w)$ & \texttt{\string\Aw} & candidate sets of Lemma HB & \ref{s3:lemHB}\\
$\mathrm{choice}(v)$, $\zp_Y=L_Y^{-2}$, sublabel $(l,c,\uslot)$, $\Zone_{Y,l,c,\uslot}$, $\rho_Y$ & \texttt{\string\zp}, \texttt{\string\uslot}, \texttt{\string\Zone} & Theorem-U zones & \ref{s5:defZones}, \ref{s5:lemZones}\\
(E1), demoted, parent-bad, good parent, $\Bad$, $\Rt(Z)$, $\Pl(Z)$, $\parU(v)$, $\mathrm{dem}$, $\mathrm{lp}$ & \texttt{\string\Bad}, \texttt{\string\Rt}, \texttt{\string\Pl}, \texttt{\string\parU} & stage-2 statuses of light parts & \ref{s5:defStages}\\
$\extph$ & \texttt{\string\extph} & external phase map: phase of a vertex's round-$l$ zone, or $*$ & \ref{s4:lemPV}, \ref{s5:lemChild}\\
U-bundle $\Bdl_{l,c}$; U-quotient multigraph $\UQ_{l,c,i}$ & \texttt{\string\Bdl}, \texttt{\string\UQ} & Theorem-U chaining & \ref{s5:lemParent}\\
$\LentU(Y)$, $\Lentext(Y)$, $\Kstd$ & \texttt{\string\LentU}, \texttt{\string\Lentext}, \texttt{\string\Kstd} & used U-lent edges; externally used lent edges; standalone edges & \ref{s5:lemKRED}\\
\hline
\end{tabular}}
\end{center}

\begin{center}
{\small
\setlength{\tabcolsep}{4pt}
\begin{tabular}{p{3.7cm}p{3.0cm}p{5.6cm}p{2.5cm}}
\multicolumn{4}{l}{(Table~\ref{s1:tabSymbols}, continued.) \emph{Standalone parts and chaining}}\\
\hline
designation $\delta$, class $Y(u)$ & -- & class of a classed port & \ref{s6:defDesign}\\
$d_{Y,l}(h)$, $m_{Y,l}$, $d^*_{Y,l}$, $\alpha_{Y,l}$ & -- & class degree, its maximum, duplicated class ports, guest core degree & \ref{s6:defDesign}\\
$\Bead_{Y,l}$, $\gamma_l=\lfloor P_{l-2}/M_l\rfloor$, giant & \texttt{\string\Bead} & deterministic bead graph; giant threshold ($(Y,l)$ is giant iff a component has more than $2\gamma_l$ edges) & \ref{s6:defDesign}\\
$\cagg_{h,l}$, $c_x(Z)$, $\cpp(u)$ & \texttt{\string\cagg}, \texttt{\string\cpp} & class counts at hubs, fresh centres, ports & \ref{s6:defDesign}\\
cluster $\clu=(A_\clu,U_\clu,B_\clu)$, $b(\clu)$, $\exc(u)$, $\Phi(\clu)$, $\MED(<)$, $\pot$ & \texttt{\string\clu}, \texttt{\string\exc}, \texttt{\string\MED}, \texttt{\string\pot} & clusters & \ref{s6:defCluster}, \ref{s6:lemMED}\\
$\LayP_j$, $\pad(u)$, $\dem^{\pm}(u)$, $\Phi_j$, $\mathcal P_j$, $F_j$ & \texttt{\string\LayP}, \texttt{\string\pad}, \texttt{\string\dem} & HCC-P data & \ref{s6:lemHCCP}\\
lend-bad; $\Lost_Z$, $\Lost_l$; $\Ret_Z=A_Z\cup F_Z\cup\Lost_Z$; $\Qs_Z=Q_Z\setminus\Lost_Z$; $O_Z$; $\XU$ & \texttt{\string\Lost}, \texttt{\string\Ret}, \texttt{\string\Qs}, \texttt{\string\XU} & stage-2 lending data and the functional $X_U$ & \ref{s6:defLending}\\
$\sco_{Y,l}$, $g_{Y,l}$, $S_0(Y,l)$ & \texttt{\string\sco} & cherry split cost, number of cherry classes, non-giant system (local to the proof) & \ref{s6:lemJSLC} (proof)\\
$J_l=\Jlost_l+\Jhub_l+\Jfr_l+\Jpar_l$; \Jp{1}--\Jp{3} & \texttt{\string\Jlost}, \texttt{\string\Jhub}, \texttt{\string\Jfr}, \texttt{\string\Jpar}, \texttt{\string\Jp} & the J-set & \ref{s6:lemJplus}\\
J-consumer, $\Obj_l$, $\LentJV_l$, \JC{1}--\JC{3} & \texttt{\string\Obj}, \texttt{\string\LentJV}, \texttt{\string\JC} & consumer interface & \ref{s6:defJconsumer}\\
$\Lent(Y)$, $\LentU(Y)$, $\LentJS(Y)$, $\LentJV(Y)$, $\Erem(Y)$, $H_0(Y)$ & \texttt{\string\Lent}, \texttt{\string\LentU}, \texttt{\string\LentJS}, \texttt{\string\LentJV}, \texttt{\string\Erem} & used lent edges and remainders (defined recursively along the processing order) & \ref{s6:consOrder}, \ref{s6:lemLent}\\
$\EV(Z)$, $\EQ(Z)$, $B_Z=\EQ(Z)$ & \texttt{\string\EV}, \texttt{\string\EQ} & TPV output & \ref{s4:lemTPV}, \ref{s6:consOrder}\\
\hline
\end{tabular}}
\end{center}

\begin{center}
{\small
\setlength{\tabcolsep}{4pt}
\begin{tabular}{p{3.7cm}p{3.0cm}p{5.6cm}p{2.5cm}}
\multicolumn{4}{l}{(Table~\ref{s1:tabSymbols}, continued.) \emph{Junction-vertex quotient}}\\
\hline
$\plab(v)$, $q_l=M_l^{-2}$, $\pi_{l,r}=2^{-(l-1-r)}$, $\Pool_{l,r}$, $\Pool_l$, $r(w)$ & \texttt{\string\plab}, \texttt{\string\Pool} & round-split pool & \ref{s7:defPool}\\
$\Cand_l(u)$, $\Hcd_l=\lambda_{l-2}^{95}/(8M_l^2)$, JV-bad & \texttt{\string\Cand}, \texttt{\string\Hcd} & candidates & \ref{s7:defCand}\\
$\xi_l$, $\Past_l$ & \texttt{\string\Past} & round-$l$ randomness; the fixed past & \ref{s7:defSchedule}\\
items, live items, $\Live$, $\clive_h$, $\thult_l=\lfloor M_l\Hcd_l/7\rfloor$, ultra, $k_h$, $\KHUB_l=4M_l$ & \texttt{\string\Live}, \texttt{\string\clive}, \texttt{\string\thult}, \texttt{\string\KHUB} & round-step data & \ref{s7:consRound}\\
PAR object, cherry, unpaired fresh leg & -- & PAR data & \ref{s7:consRound}\\
$\Used(u)$, $\Used_\kappa(h)$ & \texttt{\string\Used} & junctions already used & \ref{s7:consRound}\\
$\BP_\kappa$, $\BH_\kappa$, $B^{(\iota)}$, $Q_l$, $\Dec_l$ & \texttt{\string\BP}, \texttt{\string\BH}, \texttt{\string\Dec} & layers, sub-layers, quotient, its decomposition & \ref{s7:consRound}\\
$\copies_l$, $\payrd_l$ & \texttt{\string\copies}, \texttt{\string\payrd} & round-$l$ quotient size and random payments & \ref{s7:consRound}\\
$m_\kappa(h,w)$, $m_\kappa(w,w')$ & -- & multiplicities in HUB and PAR layers & \ref{s7:lemMULT}, \ref{s7:lemCC}\\
$\tCC_l=\lceil2\log_2M_l\rceil$ & \texttt{\string\tCC} & parameter of Lemma CC & \ref{s7:lemCC}\\
$X_{V,l}$, $\XV$ & \texttt{\string\XV} & junction-copy weight & \ref{s7:lemVstar}\\
$\omega_l(v)$ & -- & pooled-vertex weight ($\cagg_{v,l}$ for $v\in D_l$, $M_l$ otherwise) & \ref{s7:lemPay}\\
$\Xpool$, $\Xp=\XU+\Xpool+\XV$ & \texttt{\string\Xpool}, \texttt{\string\Xp} & stage-1 functional & \ref{s7:defXprime}\\
$\dett_l$ & \texttt{\string\dett} & deterministic part of the quotient size & \ref{s7:lemUHsplit}\\
\hline
\end{tabular}}
\end{center}

\paragraph{Collisions resolved} (\RTa{M10}, \RTb{J7}). In earlier formulations
several symbols had more than one meaning; here each has exactly one. The symbol
$t_l$ had three meanings and is replaced by $\tHB_l$, $\tJS_l$ and $\tCC_l$; ``$K$''
is replaced by $\KJS_l$ and $\KHUB_l$; the JS giant threshold ``$\tau_l$'' is
replaced by $\gamma_l$, with $(Y,l)$ giant iff some component has more than
$2\gamma_l$ edges, and $\tau$ is reserved for $\tau_r$ (and for the generic
threshold of Lemma~14$^\tau$); ``$\theta$'' is replaced by $\thGC$, $\thult$ and
$\thetaQ$; the vertex count ``$a_r(w)$'' is renamed $\mult_r(w)$ (no per-vertex
bound on it is used); ``$\lambda$'' denotes only $\log_2d_l$ (the TPV level is
$\vlev$, the pool label is $\plab$, the overlap constant is $\cOV$, and the leaf
containing a vertex $u$ is $\Leaf_u$); the JV class ``$J_{Y,l}$'' is written
$\LJV_{Y,l}$ ($J_l$ is the J-set); the U-lent class ``$S_{Y,l,c,\sigma}$'' is
written $\LU_{Y,l,c,\uslot}$ ($\sigma$ is the robustness exponent); the chaining
``pools'' of Theorem~U are called U-bundles $\Bdl_{l,c}$ ($\Pool_l$ is the JV
pool); the JS-LC cost ``$\Gamma_{Y,l}$'' is written $\sco_{Y,l}$ (the
\Gc{k} are galactic conditions); the TPV retirement set ``$P_Z$'' is written
$\Ret_Z$ ($P_l$ is a size threshold); Theorem~U's ``$P(Z)$, $R(Z)$'' are written
$\Pl(Z)$, $\Rt(Z)$ and its parent ``$Y(v)$'' is written $\parU(v)$ ($Y(u)$ is a
class); the assigned edges ``$A(Y)$'' are written $E_r(Y)$, and the sets
``$A(w)$'' of Lemma HB are written $\Aw(w)$ ($A_Z$ are hubs and $\Aexp=105$);
``$H_l$'' is written $\Hcd_l$ ($H_Y$ and $H_0(Y)$ are graphs); ``$U(u)$,
$U_\kappa(h)$'' are written $\Used(u)$, $\Used_\kappa(h)$; the HCC-P demands
``$\delta^\pm$'' are written $\dem^\pm$ ($\delta$ is the designation); the padding
``$\pi$'' is written $\pad$ ($\pi_{l,r}$ is the pool split); the zone probability
``$\pi_Y$'' is written $\zp_Y$; the HCC-P layer ports ``$X_j$'' are written
$\LayP_j$; the Theorem-U quotient ``$\Gamma$'' is written $\UQ$; the functionals
``$X$, $Y_{\mathrm{pool}}$, $V^*$'' are written $\XU$, $\Xpool$, $\XV$; the fresh
centre ``$f$'' is written $x$; and the set of bad light parts is written $\Bad$.

\subsection{The fix ledger}\label{s1:ssecFixes}

Table~\ref{s1:tabFixes} lists every minor fix required by the two AI red teams
(\RTa{M1}--\RTa{M12} and \RTb{J1}--\RTb{J8}, \RTb{I1}--\RTb{I16}) and by the two
later AI audits of the routing engine inside JS-LC (written
\textsf{JSLC-A$k$} and \textsf{JSLC-B$k$} after the two audits), the repair of
the one major gap found by the first clean-room review (written
\textsf{CR1-PV}; see Remark~\ref{s1:remStatus}(ii)), and the statements where
each fix is applied. No fix required new mathematics; every
fix is applied in the statement or proof it concerns, and those statements carry
the corresponding ``Red-team fixes applied here'' line.

\begin{center}
\refstepcounter{table}\label{s1:tabFixes}
\textbf{Table~\thetable.} The fix ledger.\\[4pt]
{\small
\setlength{\tabcolsep}{4pt}
\begin{tabular}{p{2.1cm}p{7.6cm}p{5.6cm}}
\textbf{Fix} & \textbf{Content} & \textbf{Applied at}\\
\hline
\RTa{M1} & used JV junction edges enter $\Lent$, $H_0$, the external lent set $\Lentext$ and the coverage statement & \ref{s6:consOrder}, \ref{s6:lemLent}, \ref{s6:thmMIXC}(b), \ref{s5:lemKRED} ($\Lentext$), \ref{s7:lemLift}\\
\RTa{M2} & Lemma CC uses $\tCC_l\defeq\lceil2\log_2M_l\rceil$, not $\tJS_l$ & \ref{s1:tabSymbols}, \ref{s7:lemCC}, \ref{s7:lemVstar}\\
\RTa{M3} & the UH$^*$ display is split into a stage-1 weight plus a conditional part & \ref{s7:lemVstar}, \ref{s7:lemUHsplit}\\
\RTa{M4} & JV$^{+*}$ written in full; (D7$'$), the RULE, $W$, HALL$'$ and $32/15$ deleted; (D5$'$) dropped; $X(u)=H_{Y(u)}$; one randomness schedule & \ref{s7:defSchedule}, \ref{s7:consRound}, \ref{s7:defCand}, \ref{s1:remNotUsed}\\
\RTa{M5} & Markov factors $12\Ex\Xp+4\sum_l\dett_l$; the bound (K6) (where $2n/M_l^2$ suffices, red team~2) & \ref{s7:lemUHsplit}(iii), \ref{s7:lemEXprime}, \ref{s6:lemLost}\\
\RTa{M6} & base case $d_1<\Dstar$ by the trivial bound $\f\le|E|$ & \ref{s7:thmHI}, step (3)\\
\RTa{M7} & the MIX-C assembly (previously checked once) written in full & \ref{s6:consOrder}, \ref{s6:thmMIXC}\\
\RTa{M8} & exponent $2^{l-r-2}$; independence from simplicity of $G$, not from \Jp{1}; ``no collision payments'' & \ref{s7:lemCand}, \ref{s7:lemUltra}, \ref{s7:lemPay}, \ref{s7:thmJVps}\\
\RTa{M9} & complete list of interface facts of the J-set & \ref{s6:lemJplus} (i)--(v)\\
\RTa{M10} & symbol collisions resolved & \ref{s1:tabSymbols}, conventions, \ref{s2:defHBtp}, \ref{s5:defZones}, \ref{s5:lemParent}, \ref{s6:defDesign}\\
\RTa{M11}, \RTa{M12} & superseded HALL statements; never used & \ref{s1:remNotUsed}\\
\hline
\RTb{J1} & JS-LC uses the cherry split only & \ref{s6:lemJSLC}, \ref{s6:remStar}\\
\RTb{J2} & $\Lent(Z)=\emptyset$ for lend-bad standalone $Z$ & \ref{s6:lemLent}(ii), \ref{s7:lemLift}\\
\RTb{J3} & the numerical claim $K^5\Prob\le0.040$ is false and is not used; only the analytic bound $K^{-4}$ is used & \ref{s7:lemWellDef}(iii)\\
\RTb{J4} & a looped cherry costs two edges & \ref{s7:consRound}(e3), \ref{s7:lemPay}\\
\RTb{J5} & wording: payments bounded by stage-1 weights; $\Past_l$ is the fixed past; Markov alone gives probability $\ge1/2$ & \ref{s7:defSchedule}, \ref{s7:consRound}(a),(f), \ref{s7:lemOneOutcome}\\
\RTb{J6} & $\Czero$ counts $\epsOne$ once & \ref{s1:defConstants}, \ref{s7:propCost}\\
\RTb{J7} & corrected candidate-count display; pool name collision (Theorem~U's pools renamed U-bundles $\Bdl_{l,c}$) & \ref{s7:lemCand}, \ref{s1:tabSymbols}\\
\RTb{J8} & group size $\lceil\Hcd_l/8\rceil$ & \ref{s7:consRound}(c), \ref{s7:lemWellDef}(iv)\\
\textsf{RT2} (status) & stale status line of the JV$^{+*}$ write-up & \ref{s1:remStatus}\\
\hline
\end{tabular}}
\end{center}

\begin{center}
{\small
\setlength{\tabcolsep}{4pt}
\begin{tabular}{p{2.1cm}p{7.6cm}p{5.6cm}}
\multicolumn{3}{l}{(Table~\ref{s1:tabFixes}, continued.)}\\
\hline
\RTb{I1} & the same as \RTa{M1}, in the input write-ups & \ref{s6:consOrder}, \ref{s6:lemLent}, \ref{s6:thmMIXC}, \ref{s5:lemKRED}\\
\RTb{I2} & the JS-LC hypothesis refers to the classes of the current round $l$ & \ref{s6:lemJSLC}\\
\RTb{I3} & invalid proof of the cuckoo bound replaced by an analytic proof & \ref{s7:lemWellDef}(iii)\\
\RTb{I4} & conditioning in Lemma CC; $\tCC_l$; the round-by-round order written out & \ref{s7:lemCC}, \ref{s7:lemVstar}, \ref{s6:consOrder}, \ref{s7:lemOneOutcome}\\
\RTb{I5} & the false claim about $(\alpha)$-edges and the undefined $(\gamma)$ removed & \ref{s2:propOrigin}\\
\RTb{I6} & $R-r\le2\logs d_r+2$ (not $+1$) & \ref{s2:lemTower}(a), \ref{s6:thmCONC}\\
\RTb{I7} & the split depends only on $(U,N)$ & \ref{s2:defTauRules}\\
\RTb{I8} & the $\HBtp$ constants are used throughout & \ref{s2:propOV}, \ref{s2:remTauConstants}, all of Section~\ref{s2:sec}; \ref{s5:lemExpect}, \ref{s5:lemParent}\\
\RTb{I9} & parentless means: no parent that is neither demoted nor parent-bad & \ref{s2:propParentless}(iii), \ref{s5:defStages}, \ref{s5:lemExpect}, \ref{s5:lemKRED}\\
\RTb{I10} & the JV class family has $2\logs d_r+1$ members (not ``$+16$'') & \ref{s3:lemCOLJV}\\
\RTb{I11} & the failure union needs \ref{s1:condG1}, not merely $\Cp>34$ & \ref{s3:lemCOLJV}, \ref{s3:lemCOL}\\
\RTb{I12} & constant and final-stage regime of Theorem~16$^*$ explicit; $\theta+1$ in Proposition~18 & \ref{s3:thmT16s}, \ref{s3:lemP18s}, \ref{s1:citProp18}\\
\RTb{I13} & TPV multiplicity bound $2+b$ (not $2+2b$), with $b$ the load bound of Lemma HB in the run, and the length-$\ge2$ argument; VX$^+$ uses $N+o(N)$ single edges & \ref{s4:lemTPV}, \ref{s4:thmVXp}\\
\RTb{I14} & Definition~7 of \cite{BM} is a multiset statement & \ref{s1:citDef7}, \ref{s3:remMultiset}, \ref{s5:lemParent}\\
\RTb{I15}, \RTb{I16} & project bookkeeping; off-route exponent: not applicable & \ref{s1:remNotUsed}\\
\hline
\textsf{JSLC-A2} & keep $k_i\defeq\min(\KJS_l,\max(1,\lfloor c_i/2\rfloor))$ with its three cases & \ref{s6:lemJSLC}\\
\textsf{JSLC-A3}, \textsf{JSLC-B2} & state the joint-routing facts: a port lies in at most $M_l-1$ systems; the multiplicity is counted over the union multiset of junction-$j$ pairs of all systems of $(Y,l)$; HCC-P needs vertex-disjointness only within a system; $g_{Y,l}\le m_{Y,l}/2+2M_l$ & \ref{s6:lemHCCglob}, \ref{s6:lemJSLC}\\
\textsf{JSLC-B3} & wording: each bead has exactly one end that is a port of its class & \ref{s6:lemJSLC}\\
\textsf{JSLC-A4} & no scratch computation is cited as evidence; toy checks are mentioned only as non-proof checks & whole manuscript; the final subsection ``\ref{s7:ssecNotEstablished}'' of Section~\ref{s7:sec}\\
\textsf{JSLC-A1}, \textsf{JSLC-B1} & confirm \RTb{I11} & \ref{s3:lemCOLJV}\\
\hline
\end{tabular}}
\end{center}

\begin{center}
{\small
\setlength{\tabcolsep}{4pt}
\begin{tabular}{p{2.1cm}p{7.6cm}p{5.6cm}}
\multicolumn{3}{l}{(Table~\ref{s1:tabFixes}, continued.)}\\
\hline
\textsf{CR1-PV} & clean-room review~1, major gap PV-ARC-END: the per-vertex arc-end bound of Lemma~PV(c) is $\deg_{H_0}(x)\le|Z|-1$ (the earlier bound $8\vJ+8+\vb\le2^9L^8$ missed arcs left by stripping in the finish); each arc end lies in at most one transition, also after visit capping, so a vertex lies in at most $M_l-1$ transition pairs; the U-lent classes use the multiplicity parameter $t_Y\defeq\lceil\lambda_r^{1.6}\rceil\ge M_l$ in place of $2^9L_Y^8$; COL-JV rows~4 and~7 re-derived ($\lambda^{42.1}\ge2^{193}\cdot12^5$ with row~9; $\lambda^{64.4}\ge2^{127}\cdot12^4$); no new mathematics; re-reviewed (clean-room review~2, verdict (A)) & \ref{s4:lemPV}(c), \ref{s5:lemChild}(c), \ref{s5:defStages}, \ref{s5:lemParent} (F-d), Steps~5--6, \ref{s3:defCOL}, \ref{s3:lemCOL}(c), \ref{s3:lemCOLJV} rows~4, 7\\
\hline
\end{tabular}}
\end{center}

\subsection{The dependency table}\label{s1:ssecDAG}

Table~\ref{s1:tabDAG} has one row for every numbered statement of
Sections~\ref{s1:sec}--\ref{s7:sec}, in the order in which the statements appear
(sectioning units such as the overview and the final subsection
``\ref{s7:ssecNotEstablished}'' of Section~\ref{s7:sec} are not statements and
have no row). The columns are: the label; a short name
and the number of the statement; its section; the statements it depends on
(this column copies the ``Dependencies'' line printed at the end of the statement); its
AI referee status; and its provenance.

\emph{Referee status.} \rc{2} means that the statement and its proof were checked
by two independent AI referees, or by one AI referee together with an independent
AI re-derivation; ``+RT'' means that in addition an AI red team re-derived it;
``+RA'' means that the two AI audits of the routing engine inside JS-LC
re-derived it in the form used for cherry systems; \rc{0} marks a new formulation
made for this manuscript (a repackaging of refereed material) that has had no
referee in this form; ``elem.'' means elementary and re-derived at least three
times; ``pub.'' means a published result and ``standard'' a standard textbook
result; ``--'' marks conventions and remarks that assert nothing beyond the
statements they cite; ``changed by \textsf{CR1-PV}: \rc{1}'' marks the part of a
statement changed by the repair of the major gap found by clean-room review~1
(Table~\ref{s1:tabFixes}), which has had one AI re-review (clean-room review~2)
since the change. A status ending in ``one clean-room AI review (G0,
2026-09-26)'' (for instance ``new in v6 (R1); one clean-room AI review (G0,
2026-09-26)'') marks a statement, or the named part of it, that was
written or changed for the present version~6 (v6) of this text, which
simplifies the text for a formal verification; the item in parentheses names
the change (item R$k$ is documented in the project file
\texttt{proofs/\allowbreak{}manuscript/\allowbreak{}v6work/\allowbreak{}R$k$.md};
``integration'' names a correction made when the items R1--R7 were combined
into this text, documented in the project file
\texttt{proofs/\allowbreak{}manuscript/\allowbreak{}v6work/\allowbreak{}INTEGRATION\_v6.md}).
Such text postdates the six review rounds mentioned below. Before G0, a
change of items R1--R7 had been checked by the AI reviews of that change (their
final repair rounds excepted), and an ``integration'' correction had had no
review. G0 (2026-09-26) is one clean-room AI review of the version-6 changes:
eight AI referee lenses that read only the manuscript and the cited papers, and
an AI judge (Remark~\ref{s1:remStatus}(iii)). It covered every row so marked
and returned verdict~(A), no fatal or major error; it is an AI review, not a
human one. A status ending in ``no referee yet'' marks text added after G0: the proof
of the stated form of Cited result~\ref{s1:citHaxell}, and the changes of
version~6.1 described next. The other changes made after G0 are the G0 fixes,
which are unmarked and have not been re-reviewed. A status ending in ``changed in v6.1 (P2
triage): no referee yet'' marks a statement, or the named part of it, changed
in version~6.1 after the triage of the formalization in Lean
(Remark~\ref{s1:remStatus}(iii); project file
\texttt{proofs/\allowbreak{}manuscript/\allowbreak{}v61/\allowbreak{}PATCHES.md});
these changes have had no review.
The rest of such an entry gives the
status of the statement before the change, which remains valid for its
unchanged parts.
\emph{Every referee is an AI agent; no statement has had a human referee. The
manuscript as a whole has had six AI review rounds, clean-room reviews~1--5
and the cross-model round (Remark~\ref{s1:remStatus}(ii)), which this column does not record
separately.}

\emph{Revision R6 of version~6} (the galactic conditions as eventualities). In
the status column, ``R6: one clean-room AI review (G0, 2026-09-26)'' marks a
statement that R6 added
(Lemma~\ref{s3:lemCOLJVev}) or whose statement or proof R6 rewrote
(Condition~\ref{s1:condGamma}, Proposition~\ref{s2:propDegRec},
Lemmas~\ref{s2:lemLacunary}, \ref{s2:lemTower}, \ref{s3:lemCOLJV} and
\ref{s7:lemGammaSat}); the status before the marker refers to version~5. R6 also
changed single proof steps or explanatory text of the following statements,
without changing what they state. Their rows carry no marker (except the row
of Lemma~\ref{s3:lemCOL}, which since the fixes of G0 is marked for step~(b) of
its proof, changed by R6 and by the integration); their status
refers to the text of version~5, and the changed steps have had one clean-room
AI review (G0) since the change:
Definition~\ref{s1:defConstants}(ii),(iii) and Remark~\ref{s1:remOrder}
(existence of the constants; order of choice);
Proposition~\ref{s2:propStructure}; Lemma~\ref{s3:lemCOL}; the setting
paragraph of Section~\ref{s5:sec} (the proof of \eqref{s5:eqZp}, used by
Definition~\ref{s5:defZones} and Lemma~\ref{s5:lemZones}); Lemmas~\ref{s5:lemE1},
\ref{s5:lemExpect}, \ref{s5:lemParent} and~\ref{s5:lemDemoted};
Remark~\ref{s5:remConstants}(b); the facts on $\logs$ before
Theorem~\ref{s6:thmCONC} (\eqref{s6:eqTowerHalf} and \eqref{s6:eqTowerEnd}),
used by Theorems~\ref{s6:thmCONC} and~\ref{s6:thmCONCL};
Lemma~\ref{s7:lemUHsplit}(iv); and Corollary~\ref{s7:thmMainProof}. Each of
these changes removes a use of the former numerical threshold (citing an item of
\Gc{1} where a bound is still needed; Lemmas~\ref{s5:lemExpect}
and~\ref{s7:lemUHsplit} now also declare \ref{s1:condG1} among their
dependencies), adds explanatory text on the existence of the constants
(Definition~\ref{s1:defConstants}, Remark~\ref{s1:remOrder}), deletes an unused
aside (Remark~\ref{s5:remConstants}(b)), or (in
Corollary~\ref{s7:thmMainProof}) takes the constant $\Dstar$ of
Theorem~\ref{s1:thmMain} instead of choosing one.

\emph{Provenance} names the project working file in which the statement and its
refereed proof were first written. It records history only: no proof in this
manuscript is delegated to any file, and every argument is written out in the
statement's own section.

\emph{What an entry means.} The entries of the row of a statement $A$ are
its logical inputs: the statements that the statement or the proof of $A$ uses
through a reference, by name (for instance ``by Hall's theorem'' or ``by
Markov's inequality''), or through a named item such as \Gc{3}, (K1), \Jp{2},
COL(g) or (R5). A fact that is proved in the running text between two
statements (for instance \eqref{s5:eqLY}, \eqref{s5:eqZp},
\eqref{s6:eqTowerHalf}, the observations of Section~\ref{s4:ssecConv}, or the
bounds on the numbers of PAR objects and hub items in the opening paragraph of
the subsection on quotient size and payments in Section~\ref{s7:sec}) has no
row; a statement that
uses such a fact lists the statements from which the fact is proved.
Convention~\ref{s1:convGraphs} (graphs are finite and simple, unless they are
called multigraphs) holds throughout; it is listed only in the rows of the
statements that cite it, directly or through an observation of
Section~\ref{s4:ssecConv}. Definition~\ref{s1:defConstants} fixes the constants
used throughout. It is listed in the row of every later statement that cites it
or invokes a requirement that it places on a constant (for instance ``by the
choice of $\Nzero$''); every other statement that uses one of these constants
reaches it through the entries of its row. The row of a remark lists the earlier
statements that the remark is about. A few rows also list the result of
\cite{BM} whose proof the statement modifies and writes out in full.
Theorem~\ref{s3:thmT16s} does not list \cite[Theorem~16]{BM}: its proof follows
that of \cite[Theorem~16]{BM}, but it is written out in full in
Section~\ref{s3:sec} and uses neither that theorem nor a step of its proof. A few
rows list a statement that is used only where the statement of the row is
applied, or only through another entry of the same row: for instance \Gc{3} in
the rows of Lemma~\ref{s4:lemTPV}, Lemma~\ref{s4:lemPV} and
Theorem~\ref{s4:thmVXp}, and \Gc{2} in the row of Lemma~\ref{s2:lemTower}. These
entries are harmless over-declarations. The following are not entries: pointers
that only say where a result is proved, used or modified, or that it is not
used (except in the row of Theorem~\ref{s1:thmMain}, whose one entry is its
proof, Corollary~\ref{s7:thmMainProof}; see below); comparisons with \cite{BM}; and the mere use of notation introduced
elsewhere. The items whose name is marked ``not used'' (or ``not used
directly'') are entries of no row.

\emph{Forward references.} Apart from pointers (see above) and the reference of
Theorem~\ref{s1:thmMain} to its proof, Corollary~\ref{s7:thmMainProof}, four
kinds of forward reference occur in the text. They are not entries, and none of
them is circular.
(1) Constants. Theorem~\ref{s1:thmMain}, Definition~\ref{s1:defConstants},
Remark~\ref{s1:remOrder} and Condition~\ref{s1:condGamma} use explicit
functions of $\Dstar$ ($\epsOne$, $\epsTwo$ and the terms they collect),
explicit size thresholds and, in \Gc{1}(f), the explicit inequalities of
column~3 of Table~\ref{s3:tabCOLJV}; these are written out in later
statements, proofs and tables.
These are explicit formulas; no lemma of the manuscript is needed to state them.
What is used about them is that finitely many explicit inequalities hold
eventually: in $N$ for $\Nzero$ (Definition~\ref{s1:defConstants}(ii) and the
size conditions of Section~\ref{s4:sec}), and in $\Dstar$ for $\Dstar$
(Lemma~\ref{s7:lemGammaSat}). The existence of $\Nzero$ and of such a $\Dstar$
is used only in Corollary~\ref{s7:thmMainProof}; every other statement uses only
the properties that Definition~\ref{s1:defConstants} requires of the fixed
constants. Remark~\ref{s1:remOrder} also says where
constants enter later (for instance, that $\cEG$ enters only through the
induction of Theorem~\ref{s7:thmHI}); this is a description, verified in those
statements.
(2) Announced consequences. Condition~\ref{s1:condGamma} records consequences
of its items that are proved later, where they are used. For instance,
\Gc{2}(b) is proved in Lemma~\ref{s2:lemTower}(b), and the bounds of \Gc{2}(c)
in Lemma~\ref{s3:lemCOL} (with rows~6 and~7 of Table~\ref{s3:tabCOLJV}) and in
Lemma~\ref{s5:lemE1}(b),(c). None of these lemmas uses \Gc{2}(b) or (c), and
\Gc{2}(a) enters them only through Lemma~\ref{s2:lemCap}; the only statement
that uses \Gc{2}(b),(c) as such, Lemma~\ref{s7:lemGammaSat}, lists all three
lemmas. Definition~\ref{s7:defCand} notes that every JV-good port $u$ has
$|\Cand_l(u)|\ge\Hcd_l$. This is Lemma~\ref{s7:lemCand}(iv), which is proved
from the definition of JV-good and from Lemma~\ref{s7:lemCand}(ii), not from the
note. Definition~\ref{s3:defCOL} notes, in its paragraph ``Consumers'', that
for $r\ge R-1$ the set $\Lent(Y)$ is empty, so that $\Erem(Y)=E_r(Y)$ or
$H_0(Y)=E_r(Y)$ (Construction~\ref{s6:consOrder}); and
Definition~\ref{s5:defStages} notes that for $r\ge R-1$
Lemma~\ref{s5:lemE1}(c) reduces to~(b). None of these conditions and
definitions uses these consequences.
(3) Constructions. That Construction~\ref{s6:consOrder} applies every device
within its hypotheses is shown in Lemma~\ref{s6:lemLent} and
Theorem~\ref{s6:thmMIXC}(a); that Construction~\ref{s7:consRound} is well defined
is Lemma~\ref{s7:lemWellDef}(i), (ii), (iv) and~(v). These results are proved
about the constructions. Besides earlier statements, each part uses only the
definitions of the steps it concerns and of the steps before them (for
Construction~\ref{s6:consOrder}, the steps at the rounds already processed), and
not the well-definedness of the step that it proves.
(4) Definition~\ref{s7:defSchedule} describes the choice of the stage-1 outcome
through the functional $\Xp$ of Definition~\ref{s7:defXprime}, and the choice of
each $\xi_l$ through step~(f) of Construction~\ref{s7:consRound}. These choices
are made in Lemma~\ref{s7:lemOneOutcome}, whose row lists
Construction~\ref{s7:consRound} and Lemma~\ref{s7:lemEXprime} (the latter rests
on Definition~\ref{s7:defXprime}). Moreover, $\Xp$ is a function of the run,
$\delta$ and stage~1 only (Definition~\ref{s7:defXprime}), so the selection~(1e)
of Definition~\ref{s7:defSchedule} depends neither on the later stages of the
schedule nor on the round step.

\begin{center}
\refstepcounter{table}\label{s1:tabDAG}
\textbf{Table~\thetable.} Dependency table of all numbered statements.\\[4pt]
{\footnotesize
\setlength{\tabcolsep}{3pt}
\begin{tabular}{p{3.0cm}p{2.4cm}p{0.5cm}p{3.8cm}p{2.6cm}p{2.7cm}}
\textbf{Label} & \textbf{Name} & \textbf{S.} & \textbf{Depends on} & \textbf{AI referee status} & \textbf{Provenance}\\
\hline
\texttt{s1:thmMain} & Main Theorem (candidate) (\ref{s1:thmMain}) & 1 & \ref{s7:thmMainProof} & composition \rc{2} (JV$^{+*}$ referees); this text \rc{0} & \prov{jv\_\allowbreak{}plus\_\allowbreak{}star.md \S0.1, P14}\\
\texttt{s1:remStatus} & status of this manuscript (\ref{s1:remStatus}) & 1 & -- & --; item (iii) (paragraph on version~6.1) added in v6.1 (P2 triage): no referee yet & \prov{redteam1.md; redteam2.md; jslc\_\allowbreak{}routing\_\allowbreak{}audits.md; cleanroom1.md}\\
\texttt{s1:convGraphs} & graph conventions (\ref{s1:convGraphs}) & 1 & -- & -- & \prov{jv\_\allowbreak{}plus\_\allowbreak{}star.md \S1}\\
\texttt{s1:defObject} & objects; $\f$ (\ref{s1:defObject}) & 1 & -- & -- & \prov{jv\_\allowbreak{}plus\_\allowbreak{}star.md \S1}\\
\texttt{s1:factAdd} & properties of $\f$ (\ref{s1:factAdd}) & 1 & \ref{s1:defObject} & elem. & \prov{jv\_\allowbreak{}plus\_\allowbreak{}star.md \S1}\\
\texttt{s1:factEG0} & long-cycle bound (\ref{s1:factEG0}) & 1 & \ref{s1:convGraphs}, \ref{s1:defObject} & new in v6 (R1); one clean-room AI review (G0, 2026-09-26) & \prov{v6work/\allowbreak{}R1.md}\\
\texttt{s1:citNotation} & B--M notation (\ref{s1:citNotation}) & 1 & -- & pub. & \prov{papers/\allowbreak{}2211.07689.txt l.123--145, 379--385}\\
\texttt{s1:citChernoff} & B--M Thm 4 (Chernoff) (\ref{s1:citChernoff}) & 1 & -- & pub. & \prov{papers/\allowbreak{}2211.07689.txt l.318--322}\\
\texttt{s1:citDef7} & B--M Def 7 (\ref{s1:citDef7}) & 1 & -- & pub. & \prov{papers/\allowbreak{}2211.07689.txt l.341--355}\\
\texttt{s1:citProp8} & B--M Prop 8 (\ref{s1:citProp8}) & 1 & -- & pub. & \prov{papers/\allowbreak{}2211.07689.txt l.386--392}\\
\texttt{s1:citLem9} & B--M Lemma 9 (\ref{s1:citLem9}) & 1 & -- & pub. & \prov{papers/\allowbreak{}2211.07689.txt l.449--488}\\
\texttt{s1:citDef11} & B--M Def 11 (\ref{s1:citDef11}) & 1 & -- & pub.; remark: $\varepsilon>0$ added in v6 (integration): one clean-room AI review (G0, 2026-09-26) & \prov{papers/\allowbreak{}2211.07689.txt l.579--594}\\
\texttt{s1:citProp12} & B--M Prop 12 (\ref{s1:citProp12}) & 1 & -- & pub. & \prov{papers/\allowbreak{}2211.07689.txt l.620--642}\\
\texttt{s1:citProp13} & B--M Prop 13 (\ref{s1:citProp13}) & 1 & -- & pub. & \prov{papers/\allowbreak{}2211.07689.txt l.643--713}\\
\hline
\end{tabular}}
\end{center}

\begin{center}
{\footnotesize
\setlength{\tabcolsep}{3pt}
\begin{tabular}{p{3.0cm}p{2.4cm}p{0.5cm}p{3.8cm}p{2.6cm}p{2.7cm}}
\multicolumn{6}{l}{(Table~\ref{s1:tabDAG}, continued.)}\\
\textbf{Label} & \textbf{Name} & \textbf{S.} & \textbf{Depends on} & \textbf{AI referee status} & \textbf{Provenance}\\
\hline
\texttt{s1:citLem14} & B--M Lemma 14 and its proof (\ref{s1:citLem14}) & 1 & -- & pub.; justification of (9) for small $n$ in the proof sketch changed in v6.1 (P2 triage): no referee yet & \prov{papers/\allowbreak{}2211.07689.txt l.714--821}\\
\texttt{s1:citLem15} & B--M Lemma 15 (not used) (\ref{s1:citLem15}) & 1 & -- & pub. & \prov{papers/\allowbreak{}2211.07689.txt l.822--836}\\
\texttt{s1:citThm16} & B--M Thm 16 (not used; see Thm 16$^*$) (\ref{s1:citThm16}) & 1 & -- & pub. & \prov{papers/\allowbreak{}2211.07689.txt l.944--951, 1327--1360}\\
\hline
\end{tabular}}
\end{center}

\begin{center}
{\footnotesize
\setlength{\tabcolsep}{3pt}
\begin{tabular}{p{3.0cm}p{2.4cm}p{0.5cm}p{3.8cm}p{2.6cm}p{2.7cm}}
\multicolumn{6}{l}{(Table~\ref{s1:tabDAG}, continued.)}\\
\textbf{Label} & \textbf{Name} & \textbf{S.} & \textbf{Depends on} & \textbf{AI referee status} & \textbf{Provenance}\\
\hline
\texttt{s1:citLem17} & B--M Lemma 17 (\ref{s1:citLem17}) & 1 & -- & pub. & \prov{papers/\allowbreak{}2211.07689.txt l.1046--1233}\\
\texttt{s1:citProp18} & B--M Prop 18 (\ref{s1:citProp18}) & 1 & -- & pub. & \prov{papers/\allowbreak{}2211.07689.txt l.1234--1261}\\
\texttt{s1:citLem19} & B--M Lemma 19 (\ref{s1:citLem19}) & 1 & -- & pub. & \prov{papers/\allowbreak{}2211.07689.txt l.1262--1326}\\
\texttt{s1:citThm21} & B--M Thm 21 (not used directly) (\ref{s1:citThm21}) & 1 & -- & pub. & \prov{papers/\allowbreak{}2211.07689.txt l.1419--1422}\\
\texttt{s1:citCor22} & B--M Cor 22 (\ref{s1:citCor22}) & 1 & -- & pub. & \prov{papers/\allowbreak{}2211.07689.txt l.1423--1442}\\
\texttt{s1:citLem25} & B--M Lemma 25 (\ref{s1:citLem25}) & 1 & -- & pub.; hypothesis $m\ge2$ added in v6 (integration): one clean-room AI review (G0, 2026-09-26) & \prov{papers/\allowbreak{}2211.07689.txt l.1502--1558}\\
\texttt{s1:remBMused} & where B--M is used (\ref{s1:remBMused}) & 1 & \ref{s1:citNotation}, \ref{s1:citChernoff}, \ref{s1:citDef7}, \ref{s1:citProp8}, \ref{s1:citLem9}, \ref{s1:citDef11}, \ref{s1:citProp12}, \ref{s1:citProp13}, \ref{s1:citLem14}, \ref{s1:citLem17}, \ref{s1:citProp18}, \ref{s1:citLem19}, \ref{s1:citCor22}, \ref{s1:citLem25}, \ref{s1:factEG0} & --; table and last sentence rewritten in v6 (R1): one clean-room AI review (G0, 2026-09-26) & \prov{papers/\allowbreak{}2211.07689.txt}\\
\hline
\end{tabular}}
\end{center}

\begin{center}
{\footnotesize
\setlength{\tabcolsep}{3pt}
\begin{tabular}{p{3.0cm}p{2.4cm}p{0.5cm}p{3.8cm}p{2.6cm}p{2.7cm}}
\multicolumn{6}{l}{(Table~\ref{s1:tabDAG}, continued.)}\\
\textbf{Label} & \textbf{Name} & \textbf{S.} & \textbf{Depends on} & \textbf{AI referee status} & \textbf{Provenance}\\
\hline
\texttt{s1:citChernoffGen} & Chernoff (Poisson-binomial) (\ref{s1:citChernoffGen}) & 1 & -- & standard & \prov{standard; lemma\_\allowbreak{}uh.md P5}\\
\texttt{s1:citMarkov} & Markov (\ref{s1:citMarkov}) & 1 & -- & standard & \prov{standard; jv\_\allowbreak{}plus\_\allowbreak{}star.md \S5.4}\\
\texttt{s1:lemBBD} & Lemma BBD (Bernstein, bounded differences) (\ref{s1:lemBBD}) & 1 & -- & new in v6 (R2), a standard result; one clean-room AI review (G0, 2026-09-26) & \prov{standard; v6work/\allowbreak{}R2.md}\\
\texttt{s1:citHall} & Hall's theorem (\ref{s1:citHall}) & 1 & -- & new in v6 (R4), a standard result; one clean-room AI review (G0, 2026-09-26) & \prov{standard; v6work/\allowbreak{}R4.md}\\
\texttt{s1:citHaxell} & Haxell 1995 (\ref{s1:citHaxell}) & 1 & -- & pub.\ (not yet checked against \cite{Hax95}; agrees with Haxell's theorem as quoted in the literature, with $r=q+1$); derivation new in v6 (R3): one clean-room AI review (G0, 2026-09-26); proof of the stated form added after G0: no referee yet & \prov{published; v6work/\allowbreak{}R3.md}\\
\texttt{s1:citEuler} & $T$-joins, Euler (\ref{s1:citEuler}) & 1 & -- & standard & \prov{standard; js\_\allowbreak{}sync.md P3}\\
\texttt{s1:defConstants} & the constants (\ref{s1:defConstants}) & 1 & -- & \rc{0} (collected from refereed sources); (ii) changed and former item (vii) removed in v6 (R1): one clean-room AI review (G0, 2026-09-26) & \prov{jv\_\allowbreak{}plus\_\allowbreak{}star.md \S1}\\
\texttt{s1:remOrder} & order of constants (\ref{s1:remOrder}) & 1 & \ref{s1:defConstants} & \rc{0} (as above) & \prov{jv\_\allowbreak{}plus\_\allowbreak{}star.md \S1}\\
\texttt{s1:condGamma} & galactic conditions (\ref{s1:condGamma}) & 1 & \ref{s1:defConstants}, \ref{s1:remOrder} & \rc{0}; R6: one clean-room AI review (G0, 2026-09-26); where \ref{s1:condG2}(b),(c) are proved corrected in v6 (R7): one clean-room AI review (G0, 2026-09-26) & \prov{jv\_\allowbreak{}plus\_\allowbreak{}star.md \S1}\\
\texttt{s1:remNotUsed} & not used on this route (\ref{s1:remNotUsed}) & 1 & -- & -- & \prov{jv\_\allowbreak{}plus\_\allowbreak{}star.md \S0.2, \S3}\\
\hline
\end{tabular}}
\end{center}

\begin{center}
{\footnotesize
\setlength{\tabcolsep}{3pt}
\begin{tabular}{p{3.0cm}p{2.4cm}p{0.5cm}p{3.8cm}p{2.6cm}p{2.7cm}}
\multicolumn{6}{l}{(Table~\ref{s1:tabDAG}, continued.)}\\
\textbf{Label} & \textbf{Name} & \textbf{S.} & \textbf{Depends on} & \textbf{AI referee status} & \textbf{Provenance}\\
\hline
\texttt{s2:defWitness} & witness (\ref{s2:defWitness}) & 2 & \ref{s1:citDef11}, \ref{s1:defConstants} & \rc{2}+RT & \prov{conc.md P1}\\
\texttt{s2:defTauRules} & $\tau$-rules and split (\ref{s2:defTauRules}) & 2 & \ref{s2:defWitness} & \rc{2}+RT & \prov{conc.md P1; mix\_\allowbreak{}c.md \S1}\\
\texttt{s2:lemSEP} & Lemma SEP (\ref{s2:lemSEP}) & 2 & \ref{s2:defTauRules} & \rc{2}+RT & \prov{conc.md P2}\\
\texttt{s2:lemThinCut} & thin cut (\ref{s2:lemThinCut}) & 2 & \ref{s2:lemSEP}, \ref{s2:defTauRules} & \rc{2}+RT & \prov{conc.md P2--P3; redteam2.md}\\
\texttt{s2:lemOVgeneric} & Lemma OV (generic overlap) (\ref{s2:lemOVgeneric}) & 2 & \ref{s2:lemSEP}, \ref{s2:defWitness}, \ref{s1:citLem14} & \rc{2}+RT & \prov{hierarchy.md (3); conc.md P1}\\
\texttt{s2:lem14tau} & Lemma 14$^\tau$ (\ref{s2:lem14tau}) & 2 & \ref{s1:citLem14}, \ref{s1:citDef11}, \ref{s2:defTauRules}, \ref{s2:lemThinCut}, \ref{s2:lemOVgeneric}, \ref{s2:lemSEP}, \ref{s2:defWitness} & \rc{2}+RT (two deep audits, with simulations) & \prov{conc.md P1; second\_\allowbreak{}referee\_\allowbreak{}conc.md}\\
\texttt{s2:defHBtp} & the hierarchy $\HBtp$ (\ref{s2:defHBtp}) & 2 & \ref{s1:citLem14}, \ref{s2:lem14tau}, \ref{s1:defConstants}, \ref{s2:lemOVgeneric}, \ref{s2:lemSEP} & \rc{2}+RT; (R2): $M_l$ made an integer (ceiling) changed in v6.1 (P2 triage): no referee yet & \prov{mix\_\allowbreak{}c.md \S1; theorem\_\allowbreak{}u.md step 1}\\
\hline
\end{tabular}}
\end{center}

\begin{center}
{\footnotesize
\setlength{\tabcolsep}{3pt}
\begin{tabular}{p{3.0cm}p{2.4cm}p{0.5cm}p{3.8cm}p{2.6cm}p{2.7cm}}
\multicolumn{6}{l}{(Table~\ref{s1:tabDAG}, continued.)}\\
\textbf{Label} & \textbf{Name} & \textbf{S.} & \textbf{Depends on} & \textbf{AI referee status} & \textbf{Provenance}\\
\hline
\texttt{s2:defAncestors} & ancestors, typing (\ref{s2:defAncestors}) & 2 & \ref{s2:defHBtp} & \rc{2} & \prov{jv\_\allowbreak{}plus\_\allowbreak{}star.md \S4; uni\_\allowbreak{}mix.md \S0}\\
\texttt{s2:lemCap} & Lemma-25 size cap (\ref{s2:lemCap}) & 2 & \ref{s1:citLem25}, \ref{s1:condG2}, \ref{s2:defHBtp}, \ref{s2:lem14tau} & \rc{2}+RT; proof of (ii): ``$M_l\ge\max(\dots)$'' changed in v6.1 (P2 triage): no referee yet & \prov{u\_\allowbreak{}inputs\_\allowbreak{}check.md (1)}\\
\texttt{s2:lemHS} & Lemma HS (\ref{s2:lemHS}) & 2 & \ref{s1:citDef11} & \rc{2}+RT & \prov{u\_\allowbreak{}inputs\_\allowbreak{}check.md (3)}\\
\texttt{s2:propStructure} & structure of a run (\ref{s2:propStructure}) & 2 & \ref{s2:defHBtp}, \ref{s2:defAncestors}, \ref{s2:lem14tau}, \ref{s2:lemCap}, \ref{s2:lemHS}, \ref{s1:citDef11}, \ref{s1:citLem14}, \ref{s2:lemSEP}, \ref{s1:condG1} & \rc{2}+RT & \prov{mix\_\allowbreak{}c.md P1; jv\_\allowbreak{}plus\_\allowbreak{}star.md \S4}\\
\texttt{s2:lemEL} & edge laminarity EL (\ref{s2:lemEL}) & 2 & \ref{s2:defHBtp}, \ref{s2:defAncestors} & \rc{2}+RT & \prov{mix\_\allowbreak{}c.md P1}\\
\texttt{s2:propOV} & overlap constants (\ref{s2:propOV}) & 2 & \ref{s2:lemOVgeneric}, \ref{s2:lem14tau}, \ref{s2:propStructure}, \ref{s2:defAncestors}, \ref{s2:defHBtp}, \ref{s2:lemSEP} & \rc{2}+RT & \prov{mix\_\allowbreak{}c.md \S1; lemma\_\allowbreak{}uh.md P0}\\
\texttt{s2:propDegRec} & degree recursion (\ref{s2:propDegRec}) & 2 & \ref{s2:lem14tau}, \ref{s2:lemCap}, \ref{s2:propOV}, \ref{s2:defHBtp}, \ref{s2:lemSEP}, \ref{s1:condG1}, \ref{s2:lemOVgeneric} & \rc{2}+RT; R6: one clean-room AI review (G0, 2026-09-26); proof of $M_l\le d_l^2$ for the integral $M_l$ changed in v6.1 (P2 triage): no referee yet & \prov{u\_\allowbreak{}inputs\_\allowbreak{}check.md (2); hierarchy.md (1)}\\
\hline
\end{tabular}}
\end{center}

\begin{center}
{\footnotesize
\setlength{\tabcolsep}{3pt}
\begin{tabular}{p{3.0cm}p{2.4cm}p{0.5cm}p{3.8cm}p{2.6cm}p{2.7cm}}
\multicolumn{6}{l}{(Table~\ref{s1:tabDAG}, continued.)}\\
\textbf{Label} & \textbf{Name} & \textbf{S.} & \textbf{Depends on} & \textbf{AI referee status} & \textbf{Provenance}\\
\hline
\texttt{s2:propExists} & existence, termination (\ref{s2:propExists}) & 2 & \ref{s2:lem14tau}, \ref{s2:propDegRec}, \ref{s2:defHBtp}, \ref{s2:lemCap}, \ref{s2:lemOVgeneric}, \ref{s2:defWitness}, \ref{s2:propOV} & \rc{2}+RT; standing assumption on $\Dstar$ stated at the start of Section~\ref{s2:sec}, changed in v6.1 (P2 triage): no referee yet & \prov{conc.md P1; mix\_\allowbreak{}c.md P1}\\
\texttt{s2:lemLacunary} & tower-lacunary sums (\ref{s2:lemLacunary}) & 2 & \ref{s2:propDegRec}, \ref{s1:condG1} & \rc{2}; R6: one clean-room AI review (G0, 2026-09-26) & \prov{conc.md P4; mix\_\allowbreak{}c.md P5}\\
\texttt{s2:lemTower} & tower facts (\ref{s2:lemTower}) & 2 & \ref{s2:propDegRec}, \ref{s2:lemLacunary}, \ref{s2:defHBtp}, \ref{s1:condG1}, \ref{s1:condG2}, \ref{s2:propStructure}, \ref{s2:lemCap}, \ref{s2:propOV}, \ref{s2:defAncestors} & \rc{2}+RT; R6: one clean-room AI review (G0, 2026-09-26) & \prov{conc.md P4; jv\_\allowbreak{}plus\_\allowbreak{}star.md \S1, \S4}\\
\texttt{s2:lemGC} & rule GC (\ref{s2:lemGC}) & 2 & \ref{s2:defHBtp}, \ref{s2:lemTower}, \ref{s2:lemEL}, \ref{s2:defAncestors} & \rc{2}+RT & \prov{mix\_\allowbreak{}c.md \S1, P1}\\
\texttt{s2:propOrigin} & ORIGIN$^\tau$ (\ref{s2:propOrigin}) & 2 & \ref{s2:lemThinCut}, \ref{s2:lemSEP}, \ref{s2:lemCap}, \ref{s2:propOV}, \ref{s2:lemEL}, \ref{s2:defHBtp}, \ref{s2:lemGC}, \ref{s2:lemOVgeneric}, \ref{s2:propStructure} & \rc{2}+RT & \prov{conc.md P3}\\
\texttt{s2:propParentless} & fresh, parentless mass (\ref{s2:propParentless}) & 2 & \ref{s2:propStructure}, \ref{s2:lemTower}, \ref{s2:lemCap}, \ref{s2:defAncestors}, \ref{s2:defHBtp} & \rc{2}+RT & \prov{theorem\_\allowbreak{}u.md step 4; jv\_\allowbreak{}plus\_\allowbreak{}star.md \S4}\\
\texttt{s2:remTauConstants} & $\tau{+}$ constants (\ref{s2:remTauConstants}) & 2 & \ref{s2:propOV}, \ref{s2:lem14tau}, \ref{s2:lemTower}, \ref{s2:defHBtp}, \ref{s2:lemOVgeneric} & -- & \prov{mix\_\allowbreak{}c.md \S1}\\
\hline
\end{tabular}}
\end{center}

\begin{center}
{\footnotesize
\setlength{\tabcolsep}{3pt}
\begin{tabular}{p{3.0cm}p{2.4cm}p{0.5cm}p{3.8cm}p{2.6cm}p{2.7cm}}
\multicolumn{6}{l}{(Table~\ref{s1:tabDAG}, continued.)}\\
\textbf{Label} & \textbf{Name} & \textbf{S.} & \textbf{Depends on} & \textbf{AI referee status} & \textbf{Provenance}\\
\hline
\texttt{s3:lemMonotone} & monotonicity (\ref{s3:lemMonotone}) & 3 & \ref{s1:citDef7}, \ref{s1:citDef11} & \rc{2}+RT & \prov{t16star.md; uni\_\allowbreak{}mix\_\allowbreak{}section.tex}\\
\texttt{s3:lemL15p} & Lemma 15$^+$ (\ref{s3:lemL15p}) & 3 & \ref{s1:citDef11}, \ref{s1:citChernoff} & \rc{2}+RT & \prov{t16star.md A}\\
\texttt{s3:remMultiset} & multiset remark (\ref{s3:remMultiset}) & 3 & \ref{s1:citDef7}, \ref{s1:citLem9} & -- & \prov{fresh\_\allowbreak{}eyes\_\allowbreak{}r14.md; redteam2.md}\\
\texttt{s3:propP13s} & Prop 13$^*$ (\ref{s3:propP13s}) & 3 & \ref{s1:citProp12}, \ref{s1:citProp13}, \ref{s1:citDef11} & \rc{2}+RT & \prov{t16star.md B}\\
\texttt{s3:lemL17s} & Lemma 17$^*$ (\ref{s3:lemL17s}) & 3 & \ref{s3:propP13s}, \ref{s1:citLem17}, \ref{s1:lemBBD}, \ref{s1:citChernoff}, \ref{s1:citChernoffGen} & \rc{2}+RT; case~(b) re-proved via Lemma BBD in v6 (R2): one clean-room AI review (G0, 2026-09-26) & \prov{t16star.md C; v6work/\allowbreak{}R2.md}\\
\texttt{s3:lemP18s} & Prop 18$^*$, Lemma 19$^*$ (\ref{s3:lemP18s}) & 3 & \ref{s1:citProp18}, \ref{s1:citLem19}, \ref{s3:lemL17s}, \ref{s1:citDef11} & \rc{2}+RT & \prov{t16star.md D}\\
\hline
\end{tabular}}
\end{center}

\begin{center}
{\footnotesize
\setlength{\tabcolsep}{3pt}
\begin{tabular}{p{3.0cm}p{2.4cm}p{0.5cm}p{3.8cm}p{2.6cm}p{2.7cm}}
\multicolumn{6}{l}{(Table~\ref{s1:tabDAG}, continued.)}\\
\textbf{Label} & \textbf{Name} & \textbf{S.} & \textbf{Depends on} & \textbf{AI referee status} & \textbf{Provenance}\\
\hline
\texttt{s3:lemL9rho} & Lemma 9$_\rho$ (\ref{s3:lemL9rho}) & 3 & \ref{s1:citLem9}, \ref{s1:citProp8}, \ref{s1:citHaxell}, \ref{s1:citDef7}, \ref{s3:remMultiset} & \rc{2}+RT for the statement; claim and final step re-proved via Haxell's theorem in v6 (R3): one clean-room AI review (G0, 2026-09-26) & \prov{t16star.md E; v6work/\allowbreak{}R3.md}\\
\texttt{s3:thmT16s} & Theorem 16$^*$ (\ref{s3:thmT16s}) & 3 & \ref{s3:lemL15p}, \ref{s3:propP13s}, \ref{s3:lemL17s}, \ref{s3:lemP18s}, \ref{s3:lemL9rho}, \ref{s3:lemMonotone}, \ref{s1:citDef11}, \ref{s1:citChernoff} & \rc{2}+RT; item (c) (``apart from Step~0'') and the paragraph on (c) in the proof changed in v6.1 (P2 triage): no referee yet & \prov{t16star.md F}\\
\texttt{s3:lemHB} & Lemma HB (\ref{s3:lemHB}) & 3 & \ref{s1:citHall}, \ref{s1:citDef11} & \rc{2}+RT for the statement; re-proved via Hall's theorem in v6 (R4): one clean-room AI review (G0, 2026-09-26) & \prov{vortex.md (Lemma HB); v6work/\allowbreak{}R4.md}\\
\texttt{s3:defCOL} & lending data COL-JV (\ref{s3:defCOL}) & 3 & \ref{s2:defAncestors}, \ref{s2:propStructure}, \ref{s2:lemTower} & \rc{2}; $t_Y$ changed by \textsf{CR1-PV}: \rc{1} & \prov{mix\_\allowbreak{}c.md \S3; jv\_\allowbreak{}plus.md \S3}\\
\texttt{s3:lemCOLJV} & COL-JV table (\ref{s3:lemCOLJV}) & 3 & \ref{s1:condG1}, \ref{s1:condGstar}, \ref{s2:lemTower}, \ref{s3:defCOL}, \ref{s2:defAncestors}, \ref{s2:propStructure}, \ref{s3:lemL15p}, \ref{s3:thmT16s} & \rc{2} (table recomputed three times independently); rows~4, 7 changed by \textsf{CR1-PV}: \rc{1}; R6: one clean-room AI review (G0, 2026-09-26); (ii) restated at every $\mu\ge\log\log\Dstar$ in v6 (integration): one clean-room AI review (G0, 2026-09-26); proof of (B6) for the integral $M_l$ changed in v6.1 (P2 triage): no referee yet & \prov{jv\_\allowbreak{}plus.md \S4, P2}\\
\hline
\end{tabular}}
\end{center}

\begin{center}
{\footnotesize
\setlength{\tabcolsep}{3pt}
\begin{tabular}{p{3.0cm}p{2.4cm}p{0.5cm}p{3.8cm}p{2.6cm}p{2.7cm}}
\multicolumn{6}{l}{(Table~\ref{s1:tabDAG}, continued.)}\\
\textbf{Label} & \textbf{Name} & \textbf{S.} & \textbf{Depends on} & \textbf{AI referee status} & \textbf{Provenance}\\
\hline
\texttt{s3:lemCOLJVev} & eventual form of COL-JV (\ref{s3:lemCOLJVev}) & 3 & -- & new; R6: one clean-room AI review (G0, 2026-09-26) & \prov{v6 revision R6}\\
\texttt{s3:lemCOL} & Lemma COL (\ref{s3:lemCOL}) & 3 & \ref{s3:defCOL}, \ref{s3:lemL15p}, \ref{s3:thmT16s}, \ref{s3:lemMonotone}, \ref{s3:lemCOLJV}, \ref{s2:lemTower}, \ref{s2:propStructure}, \ref{s2:defAncestors}, \ref{s1:condG1} & \rc{2}; (c) changed by \textsf{CR1-PV}: \rc{1}; step (b) ($\rho_lN\ge L^2$) changed in v6 (integration, R6): one clean-room AI review (G0, 2026-09-26) & \prov{mix\_\allowbreak{}c.md \S3, P2}\\
\hline
\end{tabular}}
\end{center}

\begin{center}
{\footnotesize
\setlength{\tabcolsep}{3pt}
\begin{tabular}{p{3.0cm}p{2.4cm}p{0.5cm}p{3.8cm}p{2.6cm}p{2.7cm}}
\multicolumn{6}{l}{(Table~\ref{s1:tabDAG}, continued.)}\\
\textbf{Label} & \textbf{Name} & \textbf{S.} & \textbf{Depends on} & \textbf{AI referee status} & \textbf{Provenance}\\
\hline
\texttt{s4:convVortex} & vortex notation (\ref{s4:convVortex}) & 4 & -- & --; level range and proof-local symbols added in v6 (R5): one clean-room AI review (G0, 2026-09-26) & \prov{vortex.md; uni\_\allowbreak{}mix\_\allowbreak{}section.tex; v6work/\allowbreak{}R5.md}\\
\texttt{s4:lemTPV} & Lemma TPV (\ref{s4:lemTPV}) & 4 & \ref{s3:lemL15p}, \ref{s3:thmT16s}, \ref{s3:lemHB}, \ref{s3:lemMonotone}, \ref{s1:citCor22}, \ref{s1:factEG0}, \ref{s1:citChernoff}, \ref{s1:citMarkov}, \ref{s1:condG3}, \ref{s4:convVortex}, \ref{s1:convGraphs}, \ref{s1:defConstants} & \rc{2}+RT; per-step bound \eqref{s4:eqTPVstep}, finish and final count changed in v6 (R1): one clean-room AI review (G0, 2026-09-26); level law truncated, (G2) re-proved via the new observation (MC), and (B) re-proved, in v6 (R5): one clean-room AI review (G0, 2026-09-26) & \prov{uni\_\allowbreak{}mix.md \S1, P1; v6work/\allowbreak{}R5.md}\\
\hline
\end{tabular}}
\end{center}

\begin{center}
{\footnotesize
\setlength{\tabcolsep}{3pt}
\begin{tabular}{p{3.0cm}p{2.4cm}p{0.5cm}p{3.8cm}p{2.6cm}p{2.7cm}}
\multicolumn{6}{l}{(Table~\ref{s1:tabDAG}, continued.)}\\
\textbf{Label} & \textbf{Name} & \textbf{S.} & \textbf{Depends on} & \textbf{AI referee status} & \textbf{Provenance}\\
\hline
\texttt{s4:lemPV} & Lemma PV (P-vortex) (\ref{s4:lemPV}) & 4 & \ref{s3:thmT16s}, \ref{s3:lemHB}, \ref{s3:lemMonotone}, \ref{s1:citCor22}, \ref{s1:factEG0}, \ref{s1:citChernoff}, \ref{s1:citMarkov}, \ref{s1:condG3}, \ref{s4:convVortex}, \ref{s4:lemTPV}, \ref{s1:convGraphs}, \ref{s1:defConstants} & \rc{2}+RT for the constant 369; the $\extph$-formulation \rc{0}; (c) changed by \textsf{CR1-PV}: \rc{1}; per-step bound \eqref{s4:eqPVstep}, finish and count (a) changed in v6 (R1): one clean-room AI review (G0, 2026-09-26); level law truncated, (G2) re-proved via the new observation (MC), and (B) re-proved, in v6 (R5): one clean-room AI review (G0, 2026-09-26) & \prov{theorem\_\allowbreak{}p.md \S1; theorem\_\allowbreak{}u.md step 5; v6work/\allowbreak{}R5.md}\\
\texttt{s4:thmVXp} & Theorem VX$^+$ (\ref{s4:thmVXp}) & 4 & \ref{s3:lemL15p}, \ref{s3:thmT16s}, \ref{s3:lemHB}, \ref{s3:lemMonotone}, \ref{s1:citCor22}, \ref{s1:factEG0}, \ref{s1:citChernoff}, \ref{s1:condG3}, \ref{s4:convVortex}, \ref{s1:convGraphs}, \ref{s1:defConstants} & \rc{2}+RT; displayed bound, finish and count changed in v6 (R1): one clean-room AI review (G0, 2026-09-26); level law truncated, (G2) re-proved via the new observation (MC), and (B) re-proved, in v6 (R5): one clean-room AI review (G0, 2026-09-26) & \prov{vortex.md; v6work/\allowbreak{}R5.md}\\
\texttt{s4:remTPVVX} & TPV versus VX$^+$ (not used) (\ref{s4:remTPVVX}) & 4 & \ref{s4:lemTPV}, \ref{s4:thmVXp} & -- & \prov{uni\_\allowbreak{}mix\_\allowbreak{}section.tex}\\
\hline
\end{tabular}}
\end{center}

\begin{center}
{\footnotesize
\setlength{\tabcolsep}{3pt}
\begin{tabular}{p{3.0cm}p{2.4cm}p{0.5cm}p{3.8cm}p{2.6cm}p{2.7cm}}
\multicolumn{6}{l}{(Table~\ref{s1:tabDAG}, continued.)}\\
\textbf{Label} & \textbf{Name} & \textbf{S.} & \textbf{Depends on} & \textbf{AI referee status} & \textbf{Provenance}\\
\hline
\texttt{s5:defZones} & zones (\ref{s5:defZones}) & 5 & \ref{s2:defAncestors}, \ref{s3:defCOL}, \ref{s2:propStructure}, \ref{s2:lemTower}, \ref{s1:condG1} & \rc{2} & \prov{theorem\_\allowbreak{}u.md step 4}\\
\texttt{s5:lemZones} & zone lemma (\ref{s5:lemZones}) & 5 & \ref{s5:defZones}, \ref{s2:propStructure}, \ref{s2:lemTower}, \ref{s1:condG1}, \ref{s3:defCOL} & \rc{2}+RT & \prov{theorem\_\allowbreak{}u.md step 4}\\
\texttt{s5:defStages} & stage-2 statuses (\ref{s5:defStages}) & 5 & \ref{s5:defZones}, \ref{s3:defCOL}, \ref{s3:lemHB}, \ref{s2:propParentless}, \ref{s2:defAncestors}, \ref{s2:propStructure}, \ref{s2:lemTower}, \ref{s1:citDef7} & \rc{2}; parent-bad changed by \textsf{CR1-PV}: \rc{1} & \prov{theorem\_\allowbreak{}u.md steps 4--5}\\
\texttt{s5:lemE1} & bad probabilities (\ref{s5:lemE1}) & 5 & \ref{s3:lemCOL}, \ref{s3:lemHB}, \ref{s5:lemZones}, \ref{s5:defStages}, \ref{s1:citChernoffGen}, \ref{s3:defCOL}, \ref{s3:lemL15p}, \ref{s3:lemCOLJV}, \ref{s2:propStructure}, \ref{s2:lemTower}, \ref{s1:condG1} & \rc{2}+RT & \prov{theorem\_\allowbreak{}u.md step 5; mix\_\allowbreak{}c.md P2}\\
\texttt{s5:lemExpect} & expected bad mass (\ref{s5:lemExpect}) & 5 & \ref{s5:lemE1}, \ref{s5:defStages}, \ref{s2:propOV}, \ref{s2:lemTower}, \ref{s2:lemLacunary}, \ref{s1:condG1}, \ref{s2:propStructure} & \rc{2}+RT & \prov{theorem\_\allowbreak{}u.md step 7; one\_\allowbreak{}shot.md (1)}\\
\texttt{s5:lemChild} & child side (\ref{s5:lemChild}) & 5 & \ref{s4:lemPV}, \ref{s5:defStages}, \ref{s5:lemZones}, \ref{s3:lemCOL}, \ref{s2:propStructure}, \ref{s3:defCOL}, \ref{s1:condG3} & \rc{2}+RT; (c) changed by \textsf{CR1-PV}: \rc{1} & \prov{theorem\_\allowbreak{}u.md step 5; theorem\_\allowbreak{}p.md \S1}\\
\texttt{s5:lemParent} & parent side (\ref{s5:lemParent}) & 5 & \ref{s5:lemChild}, \ref{s5:defStages}, \ref{s5:lemZones}, \ref{s3:defCOL}, \ref{s3:lemCOLJV}, \ref{s3:lemMonotone}, \ref{s1:citEuler}, \ref{s2:propOV}, \ref{s2:lemTower}, \ref{s2:propStructure}, \ref{s1:condG1}, \ref{s1:citDef7} & \rc{2}+RT; (F-d), Steps~5--6 changed by \textsf{CR1-PV}: \rc{1} & \prov{theorem\_\allowbreak{}u.md step 6; theorem\_\allowbreak{}p.md \S2}\\
\texttt{s5:lemDemoted} & demoted parts (\ref{s5:lemDemoted}) & 5 & \ref{s4:thmVXp}, \ref{s2:propStructure}, \ref{s3:lemCOLJV}, \ref{s2:lemTower}, \ref{s1:condG1}, \ref{s1:condG3}, \ref{s5:defStages}, \ref{s5:lemParent} & \rc{2}+RT & \prov{theorem\_\allowbreak{}u.md step 3; vortex.md}\\
\texttt{s5:lemKRED} & Lemma K-RED (\ref{s5:lemKRED}) & 5 & \ref{s5:lemChild}, \ref{s5:lemParent}, \ref{s5:lemDemoted}, \ref{s2:propStructure}, \ref{s2:propParentless}, \ref{s2:lemEL}, \ref{s3:defCOL}, \ref{s5:defStages} & \rc{2}+RT; for-all form in $\Lentext$ \rc{2} (JV$^{+*}$ referees) & \prov{one\_\allowbreak{}shot.md (1); uni\_\allowbreak{}mix.md P4}\\
\texttt{s5:remConstants} & 369 and 745 (\ref{s5:remConstants}) & 5 & \ref{s4:lemPV}, \ref{s5:lemKRED}, \ref{s2:propStructure}, \ref{s2:propParentless}, \ref{s2:propOV}, \ref{s2:lemTower}, \ref{s4:thmVXp}, \ref{s1:condG4}, \ref{s1:factEG0} & --; (a) rewritten in v6 (R1): one clean-room AI review (G0, 2026-09-26) & \prov{theorem\_\allowbreak{}u.md (referee 2)}\\
\hline
\end{tabular}}
\end{center}

\begin{center}
{\footnotesize
\setlength{\tabcolsep}{3pt}
\begin{tabular}{p{3.0cm}p{2.4cm}p{0.5cm}p{3.8cm}p{2.6cm}p{2.7cm}}
\multicolumn{6}{l}{(Table~\ref{s1:tabDAG}, continued.)}\\
\textbf{Label} & \textbf{Name} & \textbf{S.} & \textbf{Depends on} & \textbf{AI referee status} & \textbf{Provenance}\\
\hline
\texttt{s6:lemGATE} & Lemma GATE (\ref{s6:lemGATE}) & 6 & \ref{s1:defObject} & \rc{2}; +RA & \prov{hcc.md S0.1; hpc.md}\\
\texttt{s6:defCluster} & clusters (\ref{s6:defCluster}) & 6 & -- & \rc{2} & \prov{hcc.md S1}\\
\texttt{s6:lemMED} & Lemma MED (\ref{s6:lemMED}) & 6 & \ref{s6:defCluster} & \rc{2}; +RA & \prov{hcc.md S2}\\
\texttt{s6:lemEQLPT} & Lemma EQ-LPT (\ref{s6:lemEQLPT}) & 6 & -- & \rc{2}; +RA & \prov{hcc.md S6.1}\\
\texttt{s6:lemPAR} & Lemma PAR (\ref{s6:lemPAR}) & 6 & \ref{s1:citEuler} & \rc{2}; +RA & \prov{js\_\allowbreak{}sync.md P3}\\
\texttt{s6:lemHCCP} & Lemma HCC-P (\ref{s6:lemHCCP}) & 6 & \ref{s6:lemGATE}, \ref{s6:lemMED}, \ref{s6:defCluster} & \rc{2}; +RA (flagged in \ref{s1:remStatus}) & \prov{js\_\allowbreak{}sync.md S1, P1}\\
\texttt{s6:lemHCCglob} & union of systems (\ref{s6:lemHCCglob}) & 6 & \ref{s1:factAdd}, \ref{s6:lemHCCP} & \rc{2}; +RA (joint-routing form); citation in the proof changed in v6.1 (P2 triage): no referee yet & \prov{hcc.md S5; jslc\_\allowbreak{}routing\_\allowbreak{}audits.md}\\
\texttt{s6:defDesign} & designation, classes (\ref{s6:defDesign}) & 6 & \ref{s2:defAncestors}, \ref{s2:defHBtp} & \rc{2} & \prov{mix\_\allowbreak{}c.md \S0}\\
\texttt{s6:thmCONC} & Theorem CONC (\ref{s6:thmCONC}) & 6 & \ref{s2:propOrigin}, \ref{s2:lemTower}, \ref{s2:propOV}, \ref{s2:lemLacunary}, \ref{s6:defDesign}, \ref{s2:defHBtp}, \ref{s1:condG1} & \rc{2}+RT & \prov{conc.md P4}\\
\texttt{s6:thmCONCL} & Theorem CONC-L (\ref{s6:thmCONCL}) & 6 & \ref{s2:propOrigin}, \ref{s2:lemSEP}, \ref{s2:lemThinCut}, \ref{s2:lemGC}, \ref{s2:lemEL}, \ref{s2:lemTower}, \ref{s2:lemLacunary}, \ref{s6:thmCONC}, \ref{s2:propOV}, \ref{s2:defHBtp}, \ref{s6:defDesign}, \ref{s1:condG1} & \rc{2}+RT (four AI passes, two with simulators) & \prov{mix\_\allowbreak{}c.md \S6, P5}\\
\hline
\end{tabular}}
\end{center}

\begin{center}
{\footnotesize
\setlength{\tabcolsep}{3pt}
\begin{tabular}{p{3.0cm}p{2.4cm}p{0.5cm}p{3.8cm}p{2.6cm}p{2.7cm}}
\multicolumn{6}{l}{(Table~\ref{s1:tabDAG}, continued.)}\\
\textbf{Label} & \textbf{Name} & \textbf{S.} & \textbf{Depends on} & \textbf{AI referee status} & \textbf{Provenance}\\
\hline
\texttt{s6:defLending} & lending statuses (\ref{s6:defLending}) & 6 & \ref{s3:defCOL}, \ref{s3:lemCOL}, \ref{s5:defStages}, \ref{s6:defDesign}, \ref{s5:defZones} & \rc{2} & \prov{mix\_\allowbreak{}c.md \S5; uni\_\allowbreak{}mix\_\allowbreak{}section.tex}\\
\texttt{s6:lemLost} & the sets Lost (\ref{s6:lemLost}) & 6 & \ref{s6:defLending}, \ref{s3:lemCOL}, \ref{s5:lemE1}, \ref{s5:lemExpect}, \ref{s2:propStructure}, \ref{s2:lemTower}, \ref{s3:defCOL} & \rc{2}+RT & \prov{mix\_\allowbreak{}c.md P4; jv\_\allowbreak{}plus\_\allowbreak{}star.md \S4}\\
\texttt{s6:lemJSLC} & Lemma JS-LC (\ref{s6:lemJSLC}) & 6 & \ref{s6:lemPAR}, \ref{s6:lemMED}, \ref{s6:lemEQLPT}, \ref{s6:lemHCCP}, \ref{s6:lemHCCglob}, \ref{s6:defDesign}, \ref{s6:defLending}, \ref{s3:lemCOL}, \ref{s3:lemMonotone}, \ref{s2:lemEL}, \ref{s2:lemCap}, \ref{s2:propOV}, \ref{s2:lemTower}, \ref{s2:defHBtp}, \ref{s2:propStructure}, \ref{s3:defCOL}, \ref{s1:citDef7} & \rc{2}+RT (hypotheses, multiplicities); +RA; cycle bound stated as $126n/M_l+1.5\sum m_{Y,l}$, $\sco_{Y,l}$ and $g_{Y,l}$ moved into the proof, changed in v6.1 (P2 triage): no referee yet & \prov{mix\_\allowbreak{}c.md \S4, P3; jslc\_\allowbreak{}routing\_\allowbreak{}audits.md}\\
\texttt{s6:remStar} & star split (not used) (\ref{s6:remStar}) & 6 & \ref{s6:lemJSLC} & -- & \prov{mix\_\allowbreak{}c.md P3}\\
\texttt{s6:lemJplus} & Lemma J$^+$ (\ref{s6:lemJplus}) & 6 & \ref{s6:lemJSLC}, \ref{s6:lemPAR}, \ref{s6:defDesign}, \ref{s2:propStructure}, \ref{s2:lemEL}, \ref{s2:lemCap}, \ref{s3:defCOL}, \ref{s2:defHBtp}, \ref{s6:lemHCCP}, \ref{s6:defLending} & \rc{2}+RT & \prov{qi\_\allowbreak{}plus.md \S3, P2; jv\_\allowbreak{}plus\_\allowbreak{}star.md P1}\\
\texttt{s6:defJconsumer} & J-consumer (\ref{s6:defJconsumer}) & 6 & \ref{s3:defCOL}, \ref{s6:lemJplus} & \rc{0} (interface made explicit here) & \prov{jv\_\allowbreak{}plus\_\allowbreak{}star.md \S5.6}\\
\texttt{s6:consOrder} & processing order (\ref{s6:consOrder}) & 6 & \ref{s4:lemTPV}, \ref{s5:lemChild}, \ref{s5:lemParent}, \ref{s5:lemDemoted}, \ref{s6:lemJSLC}, \ref{s6:defJconsumer}, \ref{s6:defLending}, \ref{s5:defStages}, \ref{s4:thmVXp}, \ref{s3:defCOL} & \rc{2} (JV$^{+*}$ referees, assembly audit) & \prov{jv\_\allowbreak{}plus\_\allowbreak{}star.md \S5.6, \S5.8}\\
\texttt{s6:lemLent} & Lent sets (\ref{s6:lemLent}) & 6 & \ref{s6:consOrder}, \ref{s6:defLending}, \ref{s6:defJconsumer}, \ref{s6:lemJSLC}, \ref{s5:lemParent}, \ref{s5:lemKRED}, \ref{s3:defCOL}, \ref{s6:lemJplus}, \ref{s5:defStages}, \ref{s4:lemTPV}, \ref{s5:lemChild}, \ref{s5:lemDemoted}, \ref{s4:thmVXp} & \rc{2} (JV$^{+*}$ referees) & \prov{jv\_\allowbreak{}plus\_\allowbreak{}star.md \S5.8}\\
\texttt{s6:thmMIXC} & Theorem MIX-C (\ref{s6:thmMIXC}) & 6 & \ref{s6:consOrder}, \ref{s4:lemTPV}, \ref{s5:lemKRED}, \ref{s5:lemDemoted}, \ref{s6:lemJSLC}, \ref{s6:lemJplus}, \ref{s6:defJconsumer}, \ref{s6:lemLent}, \ref{s6:thmCONCL}, \ref{s2:propStructure}, \ref{s2:propParentless}, \ref{s2:propOV}, \ref{s2:lemTower}, \ref{s3:defCOL}, \ref{s3:lemCOL}, \ref{s4:lemPV}, \ref{s5:lemChild}, \ref{s4:thmVXp}, \ref{s2:defHBtp}, \ref{s1:condG3}, \ref{s6:defDesign}, \ref{s6:defLending} & \rc{2} (assembly audit, JV$^{+*}$ referees); this formulation \rc{0}; the $\mathcal O_{\mathrm S}$ line of the proof of (c) changed in v6.1 (P2 triage): no referee yet & \prov{mix\_\allowbreak{}c.md \S5, P4; jv\_\allowbreak{}plus\_\allowbreak{}star.md P11--P12}\\
\hline
\end{tabular}}
\end{center}

\begin{center}
{\footnotesize
\setlength{\tabcolsep}{3pt}
\begin{tabular}{p{3.0cm}p{2.4cm}p{0.5cm}p{3.8cm}p{2.6cm}p{2.7cm}}
\multicolumn{6}{l}{(Table~\ref{s1:tabDAG}, continued.)}\\
\textbf{Label} & \textbf{Name} & \textbf{S.} & \textbf{Depends on} & \textbf{AI referee status} & \textbf{Provenance}\\
\hline
\texttt{s7:defPool} & pool labels (\ref{s7:defPool}) & 7 & \ref{s2:lemTower}, \ref{s3:defCOL} & \rc{2} & \prov{lemma\_\allowbreak{}uh.md \S0; jv\_\allowbreak{}plus\_\allowbreak{}star.md \S5.2}\\
\texttt{s7:defCand} & candidates (\ref{s7:defCand}) & 7 & \ref{s7:defPool}, \ref{s3:defCOL}, \ref{s6:defDesign} & \rc{2} & \prov{jv.md \S3; jv\_\allowbreak{}plus\_\allowbreak{}star.md \S5.2}\\
\texttt{s7:lemCand} & candidate counts (\ref{s7:lemCand}) & 7 & \ref{s7:defCand}, \ref{s7:defPool}, \ref{s3:defCOL}, \ref{s3:lemCOLJV}, \ref{s2:propStructure}, \ref{s2:lemTower}, \ref{s1:citChernoffGen} & \rc{2}+RT & \prov{lemma\_\allowbreak{}uh.md P1; jv\_\allowbreak{}plus\_\allowbreak{}star.md P2}\\
\texttt{s7:defSchedule} & randomness schedule (\ref{s7:defSchedule}) & 7 & \ref{s3:defCOL}, \ref{s5:defZones}, \ref{s7:defPool}, \ref{s6:defLending}, \ref{s5:defStages}, \ref{s7:defCand}, \ref{s4:lemTPV}, \ref{s4:lemPV}, \ref{s4:thmVXp}, \ref{s5:lemChild}, \ref{s5:lemDemoted} & \rc{2} & \prov{jv\_\allowbreak{}plus\_\allowbreak{}star.md \S5.0}\\
\texttt{s7:consRound} & JV$^{+*}$ round step (\ref{s7:consRound}) & 7 & \ref{s6:lemJplus}, \ref{s6:defJconsumer}, \ref{s7:defCand}, \ref{s7:defPool}, \ref{s7:defSchedule}, \ref{s1:citMarkov} & \rc{2} (JV$^{+*}$ referees); citation in step~(f) added in v6 (R7): one clean-room AI review (G0, 2026-09-26); paragraph \emph{Fixed rules} and the claims of steps (c), (e2) changed in v6.1 (P2 triage): no referee yet & \prov{jv\_\allowbreak{}plus\_\allowbreak{}star.md \S5.7; jv.md \S3}\\
\texttt{s7:lemWellDef} & round step well defined (\ref{s7:lemWellDef}) & 7 & \ref{s7:consRound}, \ref{s6:lemJplus}, \ref{s7:lemCand}, \ref{s3:lemCOLJV}, \ref{s1:citHall}, \ref{s2:defHBtp} & \rc{2}+RT; citation in the proof of (iii) added in v6 (R7): one clean-room AI review (G0, 2026-09-26); proof of (iii): grouping reads $\Past_l$ only, changed in v6.1 (P2 triage): no referee yet & \prov{jv\_\allowbreak{}plus\_\allowbreak{}star.md P3}\\
\texttt{s7:lemMULT} & Lemma MULT (\ref{s7:lemMULT}) & 7 & \ref{s7:consRound}, \ref{s6:lemJplus}, \ref{s7:defCand}, \ref{s7:defPool}, \ref{s2:propOV} & \rc{2}+RT & \prov{lemma\_\allowbreak{}uh.md P2; jv\_\allowbreak{}plus\_\allowbreak{}star.md P4}\\
\texttt{s7:lemSimple} & $Q_l$ is simple (\ref{s7:lemSimple}) & 7 & \ref{s7:consRound} & \rc{2}+RT & \prov{jv\_\allowbreak{}plus\_\allowbreak{}star.md P5}\\
\hline
\end{tabular}}
\end{center}

\begin{center}
{\footnotesize
\setlength{\tabcolsep}{3pt}
\begin{tabular}{p{3.0cm}p{2.4cm}p{0.5cm}p{3.8cm}p{2.6cm}p{2.7cm}}
\multicolumn{6}{l}{(Table~\ref{s1:tabDAG}, continued.)}\\
\textbf{Label} & \textbf{Name} & \textbf{S.} & \textbf{Depends on} & \textbf{AI referee status} & \textbf{Provenance}\\
\hline
\texttt{s7:lemLift} & Lemma JV-L (lift) (\ref{s7:lemLift}) & 7 & \ref{s7:consRound}, \ref{s7:lemSimple}, \ref{s7:lemWellDef}, \ref{s6:lemJplus}, \ref{s6:defJconsumer}, \ref{s3:defCOL}, \ref{s2:propStructure}, \ref{s6:defLending} & \rc{2}+RT & \prov{jv.md P6; lemma\_\allowbreak{}uh.md P3; jv\_\allowbreak{}plus\_\allowbreak{}star.md P6}\\
\texttt{s7:lemCC} & Lemma CC (\ref{s7:lemCC}) & 7 & \ref{s7:consRound}, \ref{s7:lemWellDef}, \ref{s1:citChernoffGen}, \ref{s6:lemJplus} & \rc{2}+RT; proof of (i): citation of the fixed-rule restrictions changed in v6.1 (P2 triage): no referee yet & \prov{jv.md P5; jv\_\allowbreak{}plus\_\allowbreak{}star.md P7}\\
\texttt{s7:lemUltra} & ultra-hub copies (\ref{s7:lemUltra}) & 7 & \ref{s7:consRound}, \ref{s7:lemWellDef}, \ref{s6:lemJplus}, \ref{s1:citChernoffGen}, \ref{s7:lemMULT}, \ref{s7:lemCand} & \rc{2}+RT & \prov{lemma\_\allowbreak{}uh.md P5; jv\_\allowbreak{}plus\_\allowbreak{}star.md P7}\\
\texttt{s7:lemVstar} & the weight $\XV$ (\ref{s7:lemVstar}) & 7 & \ref{s7:defPool}, \ref{s2:propOV}, \ref{s2:lemTower} & \rc{2}+RT & \prov{jv\_\allowbreak{}plus\_\allowbreak{}star.md \S5.4, P7}\\
\texttt{s7:lemPay} & payments (\ref{s7:lemPay}) & 7 & \ref{s7:consRound}, \ref{s7:lemWellDef}, \ref{s6:lemJplus}, \ref{s2:propParentless}, \ref{s2:lemTower} & \rc{2}+RT & \prov{jv\_\allowbreak{}plus\_\allowbreak{}star.md P8}\\
\texttt{s7:defXprime} & the functional $\Xp$ (\ref{s7:defXprime}) & 7 & \ref{s6:defLending}, \ref{s7:lemVstar}, \ref{s7:lemPay} & \rc{2} & \prov{jv\_\allowbreak{}plus\_\allowbreak{}star.md \S5.4}\\
\texttt{s7:lemEXprime} & expectation of $\Xp$ (\ref{s7:lemEXprime}) & 7 & \ref{s7:defXprime}, \ref{s6:lemLost}, \ref{s7:lemCand}, \ref{s7:lemVstar}, \ref{s2:propOV}, \ref{s2:propStructure}, \ref{s2:lemTower}, \ref{s5:lemExpect}, \ref{s2:lemCap} & \rc{2}+RT & \prov{jv\_\allowbreak{}plus\_\allowbreak{}star.md P9}\\
\texttt{s7:lemUHsplit} & Lemma UH$^*$-split (\ref{s7:lemUHsplit}) & 7 & \ref{s7:lemCC}, \ref{s7:lemMULT}, \ref{s7:lemUltra}, \ref{s7:lemVstar}, \ref{s7:lemEXprime}, \ref{s7:defSchedule}, \ref{s2:propOV}, \ref{s2:lemTower}, \ref{s2:lemLacunary}, \ref{s1:condG4}, \ref{s1:condG1}, \ref{s7:consRound}, \ref{s7:lemWellDef}, \ref{s2:defHBtp}, \ref{s6:lemJplus} & \rc{2}+RT & \prov{jv\_\allowbreak{}plus\_\allowbreak{}star.md P7}\\
\hline
\end{tabular}}
\end{center}

\begin{center}
{\footnotesize
\setlength{\tabcolsep}{3pt}
\begin{tabular}{p{3.0cm}p{2.4cm}p{0.5cm}p{3.8cm}p{2.6cm}p{2.7cm}}
\multicolumn{6}{l}{(Table~\ref{s1:tabDAG}, continued.)}\\
\textbf{Label} & \textbf{Name} & \textbf{S.} & \textbf{Depends on} & \textbf{AI referee status} & \textbf{Provenance}\\
\hline
\texttt{s7:lemOneOutcome} & one joint outcome (\ref{s7:lemOneOutcome}) & 7 & \ref{s7:defSchedule}, \ref{s7:consRound}, \ref{s6:consOrder}, \ref{s4:lemTPV}, \ref{s4:lemPV}, \ref{s4:thmVXp}, \ref{s5:lemChild}, \ref{s5:lemDemoted}, \ref{s7:lemEXprime}, \ref{s7:lemPay}, \ref{s7:lemUHsplit}, \ref{s1:citMarkov}, \ref{s1:condG3} & \rc{2}+RT & \prov{jv\_\allowbreak{}plus\_\allowbreak{}star.md P10}\\
\texttt{s7:propCost} & cost (\ref{s7:propCost}) & 7 & \ref{s6:thmMIXC}, \ref{s6:thmCONCL}, \ref{s7:lemLift}, \ref{s7:lemPay}, \ref{s7:lemEXprime}, \ref{s7:lemOneOutcome}, \ref{s5:lemKRED}, \ref{s2:propParentless}, \ref{s2:propOV}, \ref{s7:consRound}, \ref{s2:lemTower}, \ref{s1:condG4}, \ref{s6:defLending}, \ref{s7:defXprime} & \rc{2}+RT & \prov{jv\_\allowbreak{}plus\_\allowbreak{}star.md P12}\\
\texttt{s7:thmJVps} & Theorem JV$^{+*}$ (\ref{s7:thmJVps}) & 7 & \ref{s7:propCost}, \ref{s7:lemUHsplit}, \ref{s7:lemSimple}, \ref{s7:lemLift}, \ref{s7:lemOneOutcome}, \ref{s6:thmMIXC}, \ref{s2:propExists}, \ref{s7:lemPay}, \ref{s7:lemUltra}, \ref{s1:condGamma} & \rc{2} (written JV$^{+*}$); this text \rc{0} & \prov{jv\_\allowbreak{}plus\_\allowbreak{}star.md \S6, P13}\\
\texttt{s7:thmHI} & Theorem HI$''$ (\ref{s7:thmHI}) & 7 & \ref{s1:factAdd}, \ref{s1:defObject}, \ref{s1:defConstants} & elem. & \prov{NOTES/critic\_\allowbreak{}r18.md \S2.10; jv\_\allowbreak{}plus\_\allowbreak{}star.md P14}\\
\texttt{s7:lemGammaSat} & satisfiability of the conditions (\ref{s7:lemGammaSat}) & 7 & \ref{s1:condGamma}, \ref{s7:propCost}, \ref{s7:lemUHsplit}, \ref{s5:lemExpect}, \ref{s5:lemKRED}, \ref{s2:lemTower}, \ref{s7:lemEXprime}, \ref{s6:thmCONCL}, \ref{s1:defConstants}, \ref{s3:lemCOL}, \ref{s5:lemE1}, \ref{s3:lemCOLJVev} & \rc{2}; R6: one clean-room AI review (G0, 2026-09-26) & \prov{jv\_\allowbreak{}plus\_\allowbreak{}star.md \S1}\\
\texttt{s7:thmMainProof} & proof of the Main Theorem (\ref{s7:thmMainProof}) & 7 & \ref{s7:thmHI}, \ref{s7:thmJVps}, \ref{s2:propExists}, \ref{s7:lemGammaSat}, \ref{s6:defDesign}, \ref{s1:condGamma}, \ref{s1:defConstants} & composition \rc{2}; this text \rc{0} & \prov{jv\_\allowbreak{}plus\_\allowbreak{}star.md P14}\\
\texttt{s7:remNonCirc} & non-circularity (\ref{s7:remNonCirc}) & 7 & \ref{s7:thmHI}, \ref{s7:thmJVps}, \ref{s7:propCost}, \ref{s7:lemUHsplit}, \ref{s7:lemCC}, \ref{s7:lemUltra}, \ref{s7:lemPay}, \ref{s7:lemSimple}, \ref{s7:consRound}, \ref{s4:lemTPV}, \ref{s4:lemPV}, \ref{s4:thmVXp}, \ref{s1:condG4}, \ref{s1:factEG0} & --; (4) rewritten in v6 (R1): one clean-room AI review (G0, 2026-09-26) & \prov{jv\_\allowbreak{}plus\_\allowbreak{}star.md P14}\\
\hline
\end{tabular}}
\end{center}

Every edge of this table points from a statement to one appearing earlier in the
manuscript; the only exception is the Main Theorem, stated in
Section~\ref{s1:sec} and proved in Section~\ref{s7:sec}.
